%#!latex main.tex
\chapter{$p$-adic orbifolds and monodromy.}
\addtocontents{toc}{\protect\par\vskip2mm\hskip20mm(Group theory)\par\vskip5mm}
\minitoc
\newpage

\section{\'Etale coverings and fundamental groups in $p$-adic geometry.}
\label{sec::1}

\markboth{\thechapter. $p$-ADIC ORBIFOLDS AND MONODROMY.}%
{\thesection. \'ETALE COVERINGS AND FUNDAMENTAL GROUPS.}

\begin{abst}
 In this section, inspired by De Jong's work \cite{efgonas}, we discuss various
 kinds of finite or infinite \'etale coverings in non-archimedean
 analytic geometry, and the associated fundamental groups.
 The temperate fundamental group of a $p$-adic manifold encapsulates
 information on the bad reduction properties of all its finite etale
 coverings.
\end{abst}
   

\subsection{More on the topological fundamental group.}\label{sub::1.1} 

\begin{para}\label{para::1.1.1}
 We shall work throughout in the context of Berkovich analytic geometry,
 over a fixed complete ultrametric field $(K,|\;|)$.

 What has been said about Berkovich spaces in I.\ref{sec-1} remains true
 in this broader context (\ie without any further restriction on
 $K$). 

 More precisely, we shall work with {\it paracompact strictly
 $K$-analytic spaces} in the sense of \cite{ecfnas}. We recall that these
 form a category equivalent to the category of quasi-separated rigid
 spaces over $K$ having an admissible covering of finite type.

 The main
 advantage of these spaces over rigid spaces lies in their topology: this
 is a topology in the usual sense (not a Grothendieck topology);
 moreover, each point has a fundamental system of open neighborhoods
 which are locally compact, countable at infinity and arcwise connected.

 We shall deal mostly with what we have called, for convenience,
 (analytic) {\it K-manifolds}, \ie 
 paracompact strictly $K$-analytic spaces which satisfy the following
 condition (I.\ref{para:p-adic-manifold}):
 
\begin{quote}
 {\it any $s\in S$ has a neighborhood $U(s)$ which is isomorphic 
 to an affinoid subdomain of some smooth space (in the sense of \cite{ecfnas}
 \ie a space which admits locally an \'etale morphism to the affine space 
  $\A^{\dim S}$).}
\end{quote}

\begin{rems}\label{rems:1.1.1.5}
% \medskip\noindent {\it Remarks}
 a) Any affinoid domain of a $K$-manifold is
 rigid-analytically smooth, \ie satisfies the jacobian criterion (\cf
 \cite[2.2]{ecfnas}; note also that ``smooth'' in the sense of Berkovich
 amounts to ``rigid-analytically smooth and having no boundary''); a
 fortiori, an affinoid $K$-manifold is rigid-analytically
 smooth.  Conversely, a rigid-analytically smooth (strictly) $K$-affinoid
 space $S$ is a $K$-manifold, \cf \cite[rem. 9.7]{spasalc}.

 b) This notion is slightly stronger than the notion of spaces
 locally embeddable in a smooth space considered in \cite{spasalc}. The
 main difference is that we impose that $S$ is a {\it good space} in the
 sense of \cite{ecfnas}, \ie that any point has an affinoid
 neighborhood, which is not automatic for non-classical points (this
 will be useful in \secref{sec::3}).

 c) It follows immediately from remark a) that {\it a
 $K$-manifold is nothing but a good and rigid-analytically smooth
 (paracompact strictly $K$-analytic) space.}

 d) Let $\phi:S'\to S$ be an \'etale morphism of paracompact
 strictly $K$-analytic spaces.  If $S$ is a $K$-manifold, so is $S'$.

 e) For an algebraic $K$-variety $X$, one has $X$ smooth
 $\iff$ $X^{\mathrm{an}}$ smooth $\iff$  $X^{\mathrm{an}}$ is a
 $K$-manifold.  A morphism $f$ is \'etale $\iff$ $f^{an}$ is \'etale.
\end{rems}

 \medskip 
 It follows from \cite[cor. 9.5]{spasalc} that $K$-manifolds are locally
 arcwise-connected and {\it locally simply-connected},
 hence subject to the familiar theory of topological coverings (the
 difficulty is to show that non-classical points admit contractible 
 neighborhoods).   
  In particular, any pointed $K$-manifold $(S,s)$ admits a {\it universal
 (pointed) topological covering} $(\tilde S,\tilde s)$, such that
% \begin{equation*}
 $S \cong \tilde S/\pi_1^{\mathrm{top}}(S,s)$; 
% \end{equation*}
 $\tilde S$ is a (pointed) $K$-manifold, and  $\pi_1^{\mathrm{top}}(S,s)$
 is a discrete group of automorphisms of $\tilde S$.

 Let us also recall that in dimension one, $\pi_1^{\mathrm{top}}(S,s)$
 is a free group isomorphic to the fundamental group of the dual graph
 of the semistable reduction of $S$ (\cf \cite[5.3]{efgonas}). In higher
 dimension, there is no such restriction on the discrete groups
 $\pi_1^{\mathrm{top}}(S,s)$: actually, any differentiable manifold is
 homotopy equivalent to some $K$-manifold \cite[6.1.8]{staagonf}.
\end{para}

\begin{para}\label{para::1.1.2}
 Berkovich's analysis of the homotopy of analytic spaces of higher
 dimension makes use of De Jong's alteration theorem. This allows him to
 reduce the problem to the case of the generic fiber of so-called
 polystable fibrations of formal schemes over $K^0$ (the ring of
 integers in $K$). He associates to such fibrations a simplicial set
 which encodes the combinatorics of incidence of the intersection strata
 in the special fiber, and shows that the geometric realization of this
 simplicial set is homotopy equivalent to the generic fiber (viewed as
 an analytic $K$-manifold) \cite[5.4]{spasalc}.

 In particular, in the case of (the generic fiber of) an algebraic
 polystable fibration, the 
 simplicial set is finite, hence the topological fundamental group is
 finitely generated. Using De Jong's theorem and the following result,
 one can deduce that 
 {\it the topological fundamental group of an algebraic $K$-manifold is
 finitely generated}.
\end{para}

\begin{pro}\label{pro::1.1.3}
 Let $\ovl{S}$ be a $K$-manifold, and let $Z$ be a Zariski-closed,
 nowhere dense, reduced analytic subset. Then any topological covering
 of $S:=\ovl{S}\setminus Z$ extends uniquely to a topological covering of
 $\ovl{S}$. Equivalently (since $S$ and $\ovl{S}$ are locally
 simply-connected), the natural homomorphism $\pi_1^{\mathrm{top}}(S,s)\to
 \pi_1^{\mathrm{top}}(\ovl{S},s)$ is an isomorphism.
\end{pro}
\begin{proof}
 Extensions of topological coverings from $S$ to $\ovl{S}$ are unique if
 they exist, since $\ovl{S}$ is normal (note that $S$ is connected, 
 like $\ovl{S}$, by normality of $\ovl{S}$ \cite[3.3.16]{staagonf}).
 For the existence, we argue by induction on the dimension of $Z$
 (starting from the trivial case where $Z$ is empty).  This reduces the
 question to extending topological coverings from $S$ to $\ovl{S}\setminus
 Z^{\mathrm{sing}}$ (where $Z^{\mathrm{sing}}$ denotes the non-smooth
 locus of $Z$, which 
 is closed and of dimension less than the dimension of $Z$, \cf 
 \cite[2.2]{ecfnas}); equivalently, this reduces the question to the case
 where $Z$ is (rigid-analytically) smooth.

 By unicity, the question of existence of extensions is local, \ie is a
 question at the neighborhood in $S$ of each point $z$ of $Z$; we may
 assume that $\ovl{S}$ and $Z$ are (strictly) affinoid,
 rigid-analytically smooth and equidimensional ($Z$ of codimension $r>0$
 in $\ovl{S}$). The statement is clear if $z$ has a neighborhood $U(z)$
 which is a finite union of affinoid domains $U_i$'s such that
 $U_i\setminus (U_i\cap Z)$ is simply-connected. This is indeed the
 case, by Kiehl's theorem on the existence of tubular neighborhoods
 \cite{defeaidnf}: translated into Berkovich geometry, it asserts that (for
 $S$ affinoid and $Z$ a Zariski-closed subset, both rigid-analytically
 smooth and equidimensional) $Z$ has a neighborhood which is a finite
 union of affinoid domains $U_i$'s isomorphic to $(U_i \cap Z)\times
 \matheur{D}^r$, where $\matheur{D}^r$ is an $r$-dimensional polydisk,
 and $ (U_i\cap Z)$ corresponds to $(U_i \cap Z)\times \{0\}.$ Since
 $U_i\cap Z$ is a rigid-analytically smooth affinoid domain, it is
 locally simply connected according to Berkovich; on the other hand, in the
 non-archimedean situation, $\matheur{D}^r\setminus \{0\}$ is
 simply-connected. Whence the result.
\end{proof}


\subsection{\'Etale versus topological coverings (again).}\label{sub::1.2}  

\begin{para}\label{para::1.2.1}
 Let $S$ be a connected $K$-manifold. Let $S' \xrightarrow{f} S$ be a
 finite \'etale morphism ($S'$ is then a $K$-manifold).  Its degree $n$ may be
 defined as the rank of the locally free $\mathcal{O}_S$-module
 $f_\ast\mathcal{O}_{S'}$. Assume that $K$ is algebraically closed. Then
 one has the formula 

 \begin{equation*}
  n=\sum_{y\in f^{-1}(x)}[\mathscr{H}(y):\mathscr{H}(x)]
 \end{equation*}
 where $\mathscr{H}(x)$ stands for the residue field of
 the point $x$ (if $S=\matheur{M}(\mathcal{A})$ is affinoid and
 $x$ corresponds to the bounded multiplicative seminorm $|\;|_x$, this is
 just $\mathscr{H}(x)= 
 \widehat{Q(\mathcal{A}/\Ker|\;|_x)}$).

 Then $f$ is a {\it topological covering if and only if all fibers have
 the same cardinality}, which 
 amounts to the algebraic notion of being ``completely decomposed'': $
 \mathscr{H}(y)=\mathscr{H}(f(y))$ for any $y\in S'$, 
 \cf \cite[3.2.7]{staagonf}, \cite[6.3.1]{ecfnas}.

 This criterion shows at once that the Kummer \'etale covering
 $z\mapsto z^n$ (with $\mathrm{char}(K) \nmid n$) of the 
 punctured disk is {\it not} a topological covering. 
\end{para}

\begin{para}\label{para::1.2.2}   A {\it geometric point} $\ovl{s}$ of $S$ is a
 point with value in some complete algebraically closed extension 
 $(\Omega,|\;|)$ of
 $(K,|\;|)$,  \ie a morphism ${\ovl{s}} : \matheur{M}(\Omega) \to S$ in
 the category of analytic spaces over 
 $K$ \cite[p.48]{staagonf}.  For any finite \'etale covering there are
 exactly $n$ liftings ${\ovl{s}'}:\matheur{M}(\Omega)\to S'$ of $\ovl{s}$,
 \ie $n$ geometric points of $S'$ above $S$. We usually denote by $s$ the
 point of $S$ which is the image of $\ovl{s}$.  
\end{para}

\begin{dfn} [Berkovich, De Jong] \label{def::1.2.3}
 A morphism $f:S'\to S$ is a {\it covering} (\resp {\it \'etale 
 covering}, \resp {\it topological covering}) if $S$ is covered by open
 subsets $U$ such that $\coprod V_j = f^{-1}U \to U$ and the restriction 
 of $f$ to every $V_j$ is finite (\resp \'etale finite, \resp an
 isomorphism).
\end{dfn}

\noindent There is an obvious notion of morphism
 between coverings of $S$.  The category of \'etale coverings, \resp
 finite \'etale coverings, \resp topological coverings, will be denoted
 by $\Cov^{\mathrm{et}}_S$, \resp $\Cov^{\mathrm{alg}}_S$, \resp
 $\Cov^{\mathrm{top}}_S$.
 
 \medskip There are inclusions (fully faithful embeddings) 
 \begin{equation*}
  \Cov^{\mathrm{alg}}_S \hookrightarrow  \Cov^{\mathrm{et}}_S,
   \quad \Cov^{\mathrm{top}}_S \hookrightarrow  \Cov^{\mathrm{et}}_S.
 \end{equation*}

\begin{rems}
 (i) If $S'\to S$ is an \'etale covering, then $S'$ is a $K$-manifold.

 (ii) It is worth noticing that the notions of \'etale and of 
 topological coverings are local on the base $S$, but not the
 notion of algebraic coverings (except if $S$ is compact).

 (iii) Even if $S$ is covered by finitely many $U$'s as in the
 definition, it does not follow that $S'$ is a disjoint union of
 finite algebraic coverings of $S$. For instance, a Tate elliptic
 curve (with universal covering $\tilde S\cong \G_m$) has a
 covering by two open subsets $U_1,U_2$ such that $\tilde S_{|
 U_i}$ is a disjoint union of copies of $S$.

  (iv) It is an immediate but important fact that these
 categories are stable under taking fiber product (over $S$). 
 It is also stable under taking finite disjoint unions.

  (v) One should pay attention to the fact that unlike
 $\Cov^{\mathrm{top}}_S$, $\Cov^{\mathrm{et}}_S$ is not stable 
 under arbitrary disjoint unions, nor under composition (in the sense
 that the composition of an \'etale covering $S'/S$ and an \'etale
 covering $S''/S'$ may not be an \'etale covering
 $S''/S$). However, we have:
\end{rems}

\begin{lem}\label{lem::1.2.4}
 Any morphism which is the composition of a covering (\resp \'etale
 covering) followed or preceded by a finite (\resp finite \'etale)
 morphism is a covering (\resp \'etale covering).
\end{lem}
\begin{proof}
 Let us consider the case of a covering $g:S''\to S'$
 followed by a finite covering $f:S'\to S$ (the reverse situation is
 immediate from the definition). Let $s$ be a point of $S$. Then for any
 $s'$ in the finite set $f^{-1}(s)$, there exists an open neighborhood
 $U_{s'}\subset S'$ of $s'$ such that $g^{-1}(U_{s'})$ is a disjoint
 union of spaces $U_{s',i}$, each being finite over $U_{s'}$ via $g$. We
 may assume that the $U_{s'}$ are pairwise disjoint. Since $f$ is finite,
 it is in particular a closed continuous map, hence 
 $S'$ admits a basis of open neighborhoods of $f^{-1}(\bar{s})$ of
 the form $f^{-1}(V)$. 
% $f^{-1}(s)$ admits a basis of open neighborhoods of the form $f^{-1}(V)$.
 In particular,
 there is an open neighborhood $V_s$ of $s$ such that $\displaystyle
 \coprod_{s'\in f^{-1}(s)} U_{s'}$ contains $f^{-1}(V_s)$. We may
 replace each $U_{s'}$ by its intersection with $f^{-1}(V_s)$ (which is a
 union of connected components of $f^{-1}(V_s)$). It is then clear that
 $g\circ f$ induces a finite morphism from $U_{s',i}$ to $V_s$. Hence
 $g\circ f$ is a covering.  The \'etale variant is similar.
\end{proof}

\begin{exs}\label{exs::1.2.5}
  (i) The prototype of an \'etale covering is given by the logarithm map
	\begin{equation*}
	 \log:\; \matheur{D}_{\C_p}(1,1^-)\longrightarrow
	 (\A^1_{\C_p})^{\mathrm{an}}.
	\end{equation*}
	For any $m\in \N$, the inverse image of the open disk
	\begin{equation*}
	 \matheur{D}_m=\matheur{D}\left(0,
	 \bigl(|p|^{-m+\frac{1}{p-1}}\bigr)^-\right)
	 \subset (\A^1_{\C_p})^{\mathrm{an}}
	\end{equation*}
	is a disjoint union
	\begin{equation*}
	 \coprod_{\zeta \in \mu_{p^\infty}}\;\Bigl\{q\in
	  \matheur{D}(1,1^-) \Bigm| \zeta q^{p^m}\in 
	  \matheur{D}\left(1, \bigl(|p|^{\frac{1}{p-1}}\bigr)^-\right)\Bigr\}
	\end{equation*}
	and each of these components is isomorphic to $\matheur{D}_m$ if $m=0$,
	and finite \'etale onto $\matheur{D}_m$ in general (we are using
	here the 
	factorization $q\mapsto \zeta q^{p^m}=t\mapsto \log t \mapsto
	p^{-m}\log t =\log q$ of the logarithm).

  (ii) This example may be interpreted as a Dwork period
	mapping (II.\ref{para:6.3.1}). In fact, all examples of period
	mappings considered in II.\ref{sub:6.3} are \'etale coverings
	of the period domain 
	$\mathcal{D}$. For instance, the Gross-Hopkins period mapping  
	(II.\ref{para:6.3.2}) is an \'etale covering of $
	(\P^1_{\C_p})^{\mathrm{an}}$, \cf \cite[7.2]{efgonas}. 

	\noindent Any period mapping is \'etale, but it may happen in
	some cases that it is surjective without being an \'etale
	covering, \cf \cite[5.41]{psfpg} (this somewhat mysterious
	phenomenon is due to the existence of bijective
	non-quasi-compact \'etale morphisms).

  (iii) Any (possibly infinite) topological covering of a
	finite \'etale covering is an \'etale covering after \ref{lem::1.2.4}.
\end{exs}

The following two lemmas are due to V. Berkovich and J. De Jong (and not
stated in the most general form).
 
\begin{lem}\label{lem::1.2.6}
 \begin{enumerate}
  \item An \'etale presheaf $\mathcal{F}$ on $S$ which is representable
	by an \'etale covering $S'$ is a sheaf.  For any subdomain $U
	\subset S$, $\mathcal{F}_{|U}$ is representable by $S'\times_S
	U$.
  \item Let $\mathcal{F}$ be an \'etale sheaf. Let $U_i$ an
	open covering of $S$. If for every $i$, $\mathcal{F}_{|U_i}$ is
	representable by an \'etale covering, so is $\mathcal{F}$.
 \end{enumerate}
\end{lem}
 
 For (i), \cf \cite[4.1.3, 4.1.4, 4.1.5]{ecfnas}.  Here we take advantage
 of working with {\it strict} analytic spaces.  Gluing --- using 
 \cite[1.3.3]{ecfnas} --- gives (ii) (\cf also \cite[2.3]{efgonas}).

\begin{lem}\label{lem::1.2.7} Let $S'\to S$ be an \'etale covering, and
 let $R\subset S'\times_S S'$ be a union of connected components which
 is an equivalence relation on $S'$ over $S$. Then $S'/R$ (viewed as an
 \'etale sheaf on $S$) is representable by an \'etale covering $S''\to S$. 
\end{lem}


\cf \cite[2.4]{efgonas}: by the previous lemma, the question is local on
 $S$; the proof is by reduction to the case of a finite \'etale covering
 of affinoid spaces. 

\medskip In the situation of \ref{lem::1.2.7}, we say that the $S''\to
 S$ is a {\it quotient \'etale covering of} $S'\to S$.


\subsection{Existence of \'etale paths.}\label{sub::1.3}  
 
\begin{para}\label{para::1.3.1} Let us fix a geometric point $\ovl{s}$ of
 $S$. We consider the fiber functor 
 \begin{equation*}
  F^{\mathrm{et}}_{S,\ovl{s}}=F^{\mathrm{et}}_{\ovl{s}} :
   \Cov^{\mathrm{et}}_S \longrightarrow \Sets,\;\; S' \mapsto
   \{\text{geometric points $\ovl{s}'$ of $S'$ above $\ovl{s}$}\}
 \end{equation*}
and its restrictions

\begin{equation*}
 F^{\mathrm{alg}}_{S,\ovl{s}}=F^{\mathrm{alg}}_{\ovl{s}}:
 \Cov^{\mathrm{alg}}_S\longrightarrow \Sets,
\end{equation*}
(considered in \cite{sga1}) and
\begin{equation*}
 F^{\mathrm{top}}_{S,\ovl{s}}=F^{\mathrm{top}}_{\ovl{s}}:
  \Cov^{\mathrm{top}}_S\longrightarrow \Sets.
\end{equation*} 

 An {\it \'etale path} (\resp {\it algebraic path}) from
 $\ovl{s}$ to another geometric point $\ovl{t}$ of $S$ is an isomorphism
 of functors $F^{\mathrm{et}}_{\ovl{s}} \xrightarrow{\sim}
 F^{\mathrm{et}}_{\ovl{t}}$ (\resp $F^{\mathrm{alg}}_{\ovl{s}}
 \xrightarrow{\sim} F^{\mathrm{alg}}_{\ovl{t}}$). 

 The set of \'etale paths is topologized by taking as fundamental open
 neighborhoods of an \'etale path $\alpha$ the set 
 $\Stab_{S',\ovl{s}'}\circ \alpha$, where $\Stab_{S',\ovl{s}'}$ runs among
 the stabilizers in $\Aut(F_{\bar{s}}^{\mathrm{et}})$ of arbitrary geometric
 points $\ovl{s}'$ above 
 $\ovl{s}$ in arbitrary \'etale coverings $S'/S$.

 If we do the same for algebraic paths, we get the profinite topology
 considered in \cite{sga1}. 

 On the other hand, an isomorphism of functors
 $F^{\mathrm{top}}_{\ovl{s}} \xrightarrow{\sim}
 F^{\mathrm{top}}_{\ovl{t}}$ amounts to a path up to homotopy in the
 usual sense between $s$ and $t$ in the arcwise connected space $S$
 (this is compatible with composition of paths, which is juxtaposition
 in the reverse order according to our convention); they form a
 discrete set --- in fact, a principal homogeneous space under 
 $\pi^{\mathrm{top}}(S,s)$ acting on the left. 

 \medskip The inclusions $ \Cov^{\mathrm{alg}}_S \hookrightarrow
 \Cov^{\mathrm{et}}_S ,\;\Cov^{\mathrm{top}}_S \hookrightarrow
 \Cov^{\mathrm{et}}_S$ induce continuous maps from the space of \'etale
 paths between $\ovl{s}$ and $\ovl{t}$ to the space of algebraic paths
 between $\ovl{s}$ and $\ovl{t}$ (\resp to the discrete set of paths up
 to homotopy between $s$ and $t$).

 \medskip\noindent If $f: S\to X$ is any morphism of $K$-manifolds, and
 if $\alpha$ is an \'etale path from $\ovl{s}$ to $\ovl{t}$, there is an 
 obvious notion of push-forward $f_\ast\alpha$, which is an \'etale path
 from $f\circ \ovl{s}$ to $f\circ \ovl{t}$ ; it is compatible with
 composition of \'etale paths (as composition of isomorphisms of functors).
\end{para}

\begin{para}\label{para::1.3.2} It follows from \cite[exp. V]{sga1} that
 algebraic paths exist. Using the nice topology of Berkovich spaces,
 A.J. De Jong was able to prove \cite[2.9]{efgonas}: 
\end{para}

\begin{klem}\label{lem::1.3.3}  Etale paths exist. More precisely,
 any Berkovich path $\gamma$ (up to homotopy) between $s$ and 
 $t$ lifts to an \'etale path $\ovl{\gamma}$ between $\ovl{s}$ and $\ovl{t}$.
\end{klem}

\begin{proof}
 Up to introducing intermediate geometric points, we may assume
 that $S$ is affinoid (hence compact). For any {\it finite} open cover
 $\mathcal{U}=(U_i)$ of $S$, let us denote by
 $\Cov^{\mathrm{et}}_{S,\mathcal{U}}$ the full subcategory of
 $\Cov^{\mathrm{et}}_{S}$ of \'etale coverings $f:S'\to S$ such that
 $f^{-1}(U_i)= \coprod V_{ij}$ and $V_{ij}$ is \'etale finite over $U_i$
 via $f$.  We thus have
 \begin{equation*}
 \Cov^{\mathrm{et}}_{S}=\underset{\mathcal{U}}{\liminj}\;
  \Cov^{\mathrm{et}}_{S,\mathcal{U}}, \text{ and }
  \Isom(F^{\mathrm{et}}_{\ovl{s}}, F^{\mathrm{et}}_{\ovl{t}})
  =\underset{\mathcal{U}}{\limproj}\;
  \Isom(F^{\mathrm{et}}_{\ovl{s},\mathcal{U}},
  F^{\mathrm{et}}_{\ovl{t},\mathcal{U}})
 \end{equation*}

 \medskip\noindent where $F^{\mathrm{et}}_{\ovl{s},\mathcal{U}},
 F^{\mathrm{et}}_{\ovl{t},\mathcal{U}}$ denote the restrictions of
 $F^{\mathrm{et}}_{\ovl{s} }, F^{\mathrm{et}}_{\ovl{t} }$ to
 $\Cov^{\mathrm{et}}_{S,\mathcal{U}}$.

 We have to show that $\Isom(F^{\mathrm{et}}_{\ovl{s}},
 F^{\mathrm{et}}_{\ovl{t}}) \neq \emptyset$. The idea is to construct 
 non-empty compact subsets $K_\mathcal{U}$ of
 $\Isom(F^{\mathrm{et}}_{\ovl{s},\mathcal{U}},
 F^{\mathrm{et}}_{\ovl{t},\mathcal{U}})$ (topologized as above) such that
 if $\mathcal{U}'$ refines $\mathcal{U}$, $K_{\mathcal{U}'}$ maps to
 $K_\mathcal{U} $ under the natural map
 $\Isom(F^{\mathrm{et}}_{\ovl{s},\mathcal{U}'},
 F^{\mathrm{et}}_{\ovl{t},\mathcal{U}'})\to
 \Isom(F^{\mathrm{et}}_{\ovl{s},\mathcal{U}},
 F^{\mathrm{et}}_{\ovl{t},\mathcal{U}})$; it will follow that
 $\Isom(F^{\mathrm{et}}_{\ovl{s}}, F^{\mathrm{et}}_{\ovl{t}}) \supset
 \underset{\mathcal{U}}{\limproj} K_\mathcal{U}\neq \emptyset$.

 In fact, it suffices to do so for a cofinal system of $\mathcal{U}$'s,
 which are chosen as follows. Since $S$ is {\it arcwise connected}, there
 is a continuous embedding $I=[0,1]\hookrightarrow S$ with $0\mapsto
 s,\;1\mapsto t$.
 One considers finite open coverings $\mathcal{U} =(U_1,\ldots, U_n)$
 such that
 \begin{equation*}
 U_1\cap I = [0,r_1[,\;\ldots, U_i\cap I=]r_i,t_i[,\;\ldots, U_m\cap
  I=]r_m,1],\;U_{m+1}\cap 
  I=\emptyset,\ldots
 \end{equation*}
 \noindent for some points $r_i<t_{i-1}<r_{i+1}<t_i$ in $I \;(i=2,\ldots,
 m-1)$ and some $ m\leq n$.

 Given $\mathcal{U}$, one chooses $s_i$ in $U_i\cap U_{i+1}\cap I$ (which
 amounts to a real number between $r_{i+1}$ and $t_i$), and a geometric
 point $\ovl{s}_i$ above $s_i\; (i=1,\ldots, m-1)$; one completes this
 collection of geometric points by setting $\ovl{s}_0=\ovl{s},\;
 \ovl{s}_m=\ovl{t}$.  Let us look at the diagram

 \begin{equation*}
 \Cov^{\mathrm{et}}_{S,\mathcal{U}}
 \xrightarrow{\sigma_i}
 \Ind \Cov^{\mathrm{alg}}_{\mathcal{U}_{i+1}}
 \begin{array}{c}
  \xrightarrow{F^{\mathrm{alg}}_{\ovl{s}_i, U_{i+1}}}\\
  \xrightarrow{F^{\mathrm{alg}}_{\ovl{s}_{i+1}, U_{i+1}}}\\
 \end{array}
 \Sets.
 \end{equation*}               
 One has $F^{\mathrm{alg}}_{\ovl{s}_j, U_{i+1}}\circ \sigma_i =
 F^{\mathrm{et}}_{\ovl{s}_i, \mathcal{U} },\;\;j=i,i+1$, whence
 continuous maps 
 \begin{equation*}
 \Isom(F^{\mathrm{alg}}_{\ovl{s}_i,
 U_{i+1}},F^{\mathrm{alg}}_{\ovl{s}_{i+1}, U_{i+1}})\to
 \Isom(F^{\mathrm{et}}_{\ovl{s}_i, \mathcal{U}
 },F^{\mathrm{et}}_{\ovl{s}_{i+1}, \mathcal{U}}).
 \end{equation*} 
 Since algebraic paths exist, $\Isom(F^{\mathrm{alg}}_{\ovl{s}_i,
 U_{i+1}},F^{\mathrm{alg}}_{\ovl{s}_{i+1}, U_{i+1}})$ is a non-empty
 compact set for any $i=0,\ldots, m$.  On composing such isomorphisms for
 $i=0,\ldots, m$, one gets a continuous map
 \begin{equation*}
 \sigma_{\mathcal{U},(\ovl{s}_i)}: \prod_{i=0}^m
  \Isom(F^{\mathrm{alg}}_{\ovl{s}_{i},
  U_{i+1}},F^{\mathrm{alg}}_{\ovl{s}_{i+1}, U_{i+1}})\to
  \Isom(F^{\mathrm{et}}_{\ovl{s},\mathcal{U}},
  F^{\mathrm{et}}_{\ovl{t},\mathcal{U}})
 \end{equation*} 
 whose image is a non-empty compact set $K_\mathcal{U}$.

 In fact, while the map $\sigma_{\mathcal{U},(\ovl{s}_i)}$ depends on the
 choice of the geometric points $\ovl{s}_i$, its image $K_\mathcal{U}$
 does not. Indeed, for any other choice $\ovl{s}'_i$, there is an
 algebraic path
 \begin{equation*}
  \gamma_i:\;F^{\mathrm{alg}}_{\ovl{s}_i, U_{i+1}\cap U_i} \xrightarrow{\sim}
   F^{\mathrm{alg}}_{\ovl{s}'_i, U_{i+1}\cap
   U_i},\;\;\gamma_0=\id,\;\gamma_m=\id.
 \end{equation*} 
 We denote by the same letter the induced algebraic
 paths on $U_i$ and on $U_{i+1}$. Then 
 $\sigma_{\mathcal{U},(\ovl{s}_i)}(\alpha_0,\ldots,\alpha_m) 
 =
 \sigma_{\mathcal{U},(\ovl{s}'_i)}(\gamma_1\alpha_0\gamma_0^{-1},\ldots,
 \gamma_{m}\alpha_m\gamma_{m-1}^{-1})$, whence the independence.

 From this observation, it is straightforward to conclude that if
 $\mathcal{U}'$ refines $\mathcal{U}$, $K_{\mathcal{U}'}$ maps to 
 $K_\mathcal{U} $ under the natural map
 $\Isom(F^{\mathrm{et}}_{\ovl{s},\mathcal{U}'},
 F^{\mathrm{et}}_{\ovl{t},\mathcal{U}'})\to
 \Isom(F^{\mathrm{et}}_{\ovl{s},\mathcal{U}},
 F^{\mathrm{et}}_{\ovl{t},\mathcal{U}})$. This establishes the first
 assertion.

 For the second one (the surjectivity of
 $\Isom(F_{\ovl{s}}^{\mathrm{et}}), F_{\ovl{t}}^{\mathrm{et}}) \to
 \Isom (F_{\ovl{s}}^{\mathrm{top}}, F_{\ovl{t}}^{\mathrm{top}})$), 
 one may argue as follows: let $\ovl{\tilde s}$ be a
 lifting of $\ovl{s}$ to the universal topological covering $\pi:\tilde
 S \to S$. Then $\gamma.\bar{\tilde s}$ is a lifting of $\ovl{t}$
 to the universal topological covering.  Let $\alpha$ be an \'etale path
 between $\bar{\tilde s}$ and $\gamma.\bar{\tilde s}$. Then
 $\ovl{\gamma}:=\pi_\ast(\alpha)$ is an \'etale path between $\ovl{s}$
 and $\ovl{t}$ which maps to the ordinary path $\gamma$ up to homotopy,
 viewed as an isomorphism $F_{\ovl{s}}^{\mathrm{top}}\to
 F_{\ovl{t}}^{\mathrm{top}}$.
\end{proof}


\subsection{The formalism of fundamental groups.}\label{sub::1.4}  
 
\begin{para}\label{para::1.4.1} Let $\Cov^\bullet_S$ be a full
 subcategory of $\Cov^{\mathrm{et}}_S$ which is {\it stable under taking
 connected components, fiber products} (over $S$) {\it and quotients}
 (in the sense of \ref{lem::1.2.7}.). Examples:  $\Cov^{\mathrm{et}}_S$,
 $\Cov^{\mathrm{top}}_S$, $\Cov^{\mathrm{alg}}_S$.

 We denote by $F^\bullet_{S,\ovl{s}}=F^\bullet_{\ovl{s}}$ the
 restriction of $F^{\mathrm{et}}_{\ovl{s}}$ to $\Cov^\bullet_S$, and set 

 \begin{equation*}
  \pi_1^{\bullet}(S,\ovl{s}) = \Aut F^\bullet_{\ovl{s}}
 \end{equation*}
 topologized by considering as fundamental
 open neighborhoods of $1$ the stabilizers $\Stab^\bullet_{S',\ovl{s}'}$
 in $\pi_1^{\bullet}(S,\ovl{s})$ of arbitrary geometric points $\ovl{s}'$
 above $\ovl{s}$ in arbitrary \'etale coverings $S'\to S$ in
 $\Cov^\bullet_S\;$ (here we see $\pi_1^{\bullet}(S,\ovl{s})$ as acting
 on the left on   $F^\bullet_{\ovl{s}}(S')$).

 That the topology is indeed compatible with the group law is due to
 the fact that the system of subgroups $\Stab^\bullet_{S',\ovl{s}'}$ is
 stable under intersection (since $\Cov^\bullet_S$ is stable under
 taking fiber products) and under conjugation (which amounts to a change
 of $\ovl{s}'$).

 \medskip In the case $\Cov^\bullet_S = \Cov^{\mathrm{et}}_S,\;
 \pi_1^{\mathrm{et}}(S,\ovl{s})$ is the \'etale fundamental group
 introduced and studied by De Jong \cite{efgonas}.

 In the case $\Cov^\bullet_S = \Cov^{\mathrm{top}}_S,\;
 \pi_1^{\mathrm{top}}(S,\ovl{s})=\pi_1^{\mathrm{top}}(S, s)$ is (an
 incarnation of) the topological fundamental group.

 In the case $\Cov^\bullet_S = \Cov^{\mathrm{alg}}_S,\;
 \pi_1^{\mathrm{alg}}(S,\ovl{s})$ is a profinite group, which will be
 called the {\it algebraic fundamental group}. The theory is embodied in
 \cite[V]{sga1}.
 If $S$ is the analytification of an algebraic smooth
 $K$-variety $S^{\mathrm{alg}}$ and if $\mathrm{char}(K)=0$, then this
 group coincides with Grothendieck's algebraic fundamental 
 group. This follows from the {\it Gabber-L\"utkebohmert theorem},
 according to which the GAGA functor is an equivalence between the
 categories of finite \'etale coverings of $S^{\mathrm{alg}}$ and $S$
 respectively \cite{repfapf}.
\end{para}

\begin{lem}\label{lem::1.4.2} The natural continuous map 
 \begin{equation*}
  \phi: \pi_1^\bullet(S,\ovl{s})\to
  \underset{\Stab^\bullet_{S',\ovl{s}'}}{\limproj} 
  \pi_1^\bullet(S,\ovl{s})/\Stab^\bullet_{S',\ovl{s}'}
 \end{equation*}
 is a homeomorphism. In particular,
 $\pi_1^\bullet(S,\ovl{s})$ is a separated pro-discrete (hence totally
 discontinuous) topological space.
\end{lem}

\begin{proof}
 $\phi$ sends any neighborhood of $\id\in
 \pi_1^\bullet(S,\ovl{s})$ to a neighborhood of $\phi(\id)\in
 im(\phi)$. Since $\phi$ is equivariant with respect to the continuous
 action of $\pi_1^\bullet(S,\ovl{s})$ by left translation, it is open
 onto its image. Hence it suffices to show that it is bijective. Let
 $\underline{\gamma}:=(\gamma_{S',\ovl{s}'})$ be a compatible family of
 elements of $\pi_1^\bullet(S,\ovl{s})/\Stab^\bullet_{S',\ovl{s}'}$, \ie
 represent an element of $\limproj\;
 \pi_1^\bullet(S,\ovl{s})/\Stab^\bullet_{S',\ovl{s}'}$.

 \noindent For any $S'$ in $\Cov^{\bullet}_S$ and any $\ovl{s}'\in
 F^\bullet_{\ovl{s}}(S')$, $\gamma_{S',\ovl{s}'}.\ovl{s}'$ is another
 well-defined element of $F^\bullet(S')$, which we denote by
 $\gamma.\ovl{s}'$. From the compatibility of the
 $\gamma_{S',\ovl{s}'}$'s, it follows that $\gamma$ defines an
 endomorphism of the functor $F^\bullet_{\ovl{s}}$.

 On the other hand, one has a homeomorphism 
 \begin{equation*}
 \inv:\; \limproj\; \pi_1^\bullet(S,\ovl{s}) /\Stab^\bullet_{S',\ovl{s}'}
  \to \limproj\;
  \Stab^\bullet_{S',\ovl{s}'}\backslash\pi_1^\bullet(S,\ovl{s})
 \end{equation*}

 \noindent induced by $g\mapsto g^{-1}$. If we copy the construction
 $\underline{\gamma}\mapsto \gamma$ starting from
 $\inv(\underline{\gamma})$ and using right actions instead of left
 actions, the resulting endomorphism of $F^\bullet_{\ovl{s}}$ is nothing
 but the inverse of $\gamma$. Hence $\gamma$ is an element of $
 \pi_1^\bullet(S,\ovl{s})$, and it is clear that
 $\underline{\gamma}\mapsto \gamma$ is inverse to the natural map
 $\displaystyle\pi_1^\bullet(S,\ovl{s})\to \limproj\; 
 \pi_1^\bullet(S,\ovl{s})/\Stab^\bullet_{S',\ovl{s}'}$.
\end{proof}

\begin{para}\label{para::1.4.3} It is clear from the definitions that 
 $F^\bullet_{\ovl{s}}$ may be enriched to a functor
 \begin{equation*}
  \Cov^\bullet_S \longrightarrow
   \text{$\pi_1^\bullet(S,\ovl{s})$-$\Sets$},\;\;S'\mapsto
   F^\bullet_{\ovl{s}}(S')
 \end{equation*}
 where $\pi_1^\bullet(S,\ovl{s})$-$\Sets$ is the category of discrete
 sets endowed with a continuous 
 left action of $\pi_1^\bullet(S,\ovl{s})$. 
\end{para}

\begin{pro}\label{pro::1.4.4} Up to isomorphism (unique up to composition
 by inner automorphisms), $\pi_1^\bullet(S,\ovl{s})$ does not depend on
 $\ovl{s}$.
\end{pro}
\begin{proof}
 The space
 $\Isom(F^\bullet_{\ovl{s}},F^\bullet_{\ovl{t}})$, non-empty due to
 \ref{lem::1.3.3}, is formally a
 $\pi_1^{\mathrm{et}}(S,\ovl{t})-\pi_1^{\mathrm{et}}(S,\ovl{s})$-bitorsor,
 and the result follows.
\end{proof}

\begin{thm}\label{thm::1.4.5} The enriched functor $F^\bullet_{\ovl{s}}$
\begin{equation*}
 \Cov^\bullet_S \longrightarrow \text{$\pi_1^\bullet(S,\ovl{s})$-$\Sets$}
\end{equation*}
is fully faithful, and extends to an equivalence of categories  

\begin{equation*}
\{\text{disjoint unions of objects of }\Cov^\bullet_S\} \longrightarrow
\text{$\pi_1^\bullet(S,\ovl{s})$-$\Sets$}.
\end{equation*}
Connected coverings correspond to $\pi_1^\bullet(S,\ovl{s})$-orbits.
\end{thm}

\begin{proof}
 (\cf \cite[2.10]{efgonas} in the \'etale case). If
 $f:S'\to S$ is an \'etale covering in $\Cov^\bullet_S$ and
 $\ovl{s}',\;\ovl{s}''$ are elements of $F_{\ovl{s}}(S')$ in the same
 connected component of $S'$, proposition \ref{lem::1.3.3} ensures the existence of
 an \'etale path $\alpha$ from $\ovl{s}'$ to $\ovl{s}''$. Then
 $f_\ast(\alpha)\in \pi_1^{\mathrm{et}}(S,\ovl{s})$, and
 $f_\ast(\alpha).\ovl{s}'=\ovl{s}''$; of course, we may replace here
 $f_\ast(\alpha)$ by its image in $\pi_1^\bullet(S,\ovl{s})$. This shows
 that $\pi_1^\bullet(S,\ovl{s})$-orbits in $F^\bullet_{\ovl{s}}(S')$
 correspond bijectively to connected components of $S'$. In particular

 \begin{equation*}
 F^\bullet_{\ovl{s}}(S') = 
  \pi_1^\bullet(S,\ovl{s})/{\Stab^\bullet_{S',\ovl{s}'}}
 \end{equation*}
 if $S'$ is connected.

 This implies that the enriched functor $F^\bullet_{\ovl{s}}$ is fully
 faithful, because if $S'$ is connected, a morphism between \'etale
 coverings $S'\to S''$ of $S$ may be identified via its graph with a
 connected component $T$ of the \'etale covering $S'\times_S S''$ (which
 belongs to $\Cov^\bullet_S$ by assumption) such that
 $F^\bullet_{\ovl{s}}(T)\xrightarrow{\sim}
 F^\bullet_{\ovl{s}}(S')$ (this can be checked locally on $S$, and this
 reduces to a well-known fact about finite algebraic coverings).

 It remains to show that any $\pi_1^\bullet(S,\ovl{s})$-orbit comes from
 a (connected) \'etale covering. By definition, such an orbit may be
 written $ \pi_1^\bullet(S,\ovl{s})/H$, where $H$ is a subgroup
 containing the stabilizer ${\Stab^\bullet_{S',\ovl{s}'}}$ of some
 geometric point $\ovl{s}'$ above $\ovl{s}$ in some connected \'etale
 covering $S'$ of $S$. As we have seen, the connected components of
 $S'\times_S S'$ correspond bijectively to
 $\pi_1^\bullet(S,\ovl{s})$-orbits in $F^\bullet_{\ovl{s}}(S'\times_S
 S')= F^\bullet_{\ovl{s}}(S')^2.$ The union $R\subset S'\times_S S'$ of
 connected components which corresponds to the union of points
 $(\ovl{s}',\ovl{s}'.h),\;h\in H$, is an equivalence relation on $S'$
 over $S$. By lemma \ref{lem::1.2.7}, the quotient $S''=S'/R$ is an
 \'etale covering 
 of $S$ which belongs to $\Cov^\bullet_S$ by assumption. If $\ovl{s}''$
 denotes the image of $\ovl{s}'$ in $S''$, it is clear that
 $H=\Stab^\bullet_{S'',\ovl{s}''}$.
\end{proof}

\begin{para}\label{para::1.4.6} A connected \'etale covering $S'\to S$ is
 {\it Galois} if $S$ is the quotient $S'/(\Aut_S S')$; $\Aut_S S'$ is
 called the {\it Galois group} of the covering. By \cite[4.1.9]{ecfnas},
 an \'etale Galois covering of $S$ with group $G$ is the same as a
 connected principal homogeneous space under the constant 
 $S$-group $G_S$\footnote{{\it loc.~cit.}  4.1.4, which is used in the proof,
 is available since we are dealing with strictly analytic spaces, \cf
 {\it loc.~cit.} 4.1.5.}.

\noindent Through the dictionary
 of \ref{thm::1.4.5}, a Galois \'etale covering in $\Cov^{\bullet}_S$ 
 corresponds to a surjective continuous 
 homomorphism 

 \begin{equation*}
  \pi_1^\bullet(S,\ovl{s})\to G,
 \end{equation*}
 with $G$ discrete ($\cong \Aut_S S'$). In other words, a (pointed) covering
 $(S',\ovl{s}')$ in $\;\Cov^{\bullet}_S\;$ is Galois if and only if 
 ${\Stab^\bullet_{S',\ovl{s}'}}\subset
 \pi_1^\bullet(S,\ovl{s})$ is normal. 

 Example: the logarithm (\ref{exs::1.2.5}. (i)) is a Galois \'etale covering of
 $(\A^1_{\C_p})^{\mathrm{an}}$ with group $\mu_{p^\infty}\cong \Q_p/\Z_p$. 

The following is also clear:


\end{para}

\begin{cor}\label{cor::1.4.7} $\pi_1^\bullet(S,\ovl{s})$ is a
 pro-discrete group if and only if in $\Cov^\bullet_S$ any 
 connected covering is dominated by a Galois covering.
\end{cor}

\par\noindent {\it Remark}. This property holds in
$\Cov^{\mathrm{top}}_S$ and $\Cov^{\mathrm{alg}}_S$, but not in
$\Cov^{\mathrm{et}}_S$ in general. Let for instance $S=\P^1_{\C_p}$, and
let us consider the Galois covering given by the 
Gross-Hopkins period mapping (II.\ref{para:6.3.3}). According to 
\cite[7.4]{efgonas} 
(using the viewpoint sketched in II.\ref{sub:6.4}), the corresponding
discrete quotient of $\pi_1^{\mathrm{et}}(\P^1,\ovl{s})$ is the composition 

\begin{equation*}
 \pi_1^{\mathrm{et}}(\P^1,\ovl{s})\to SL_2(\Q_p) \to SL_2(\Q_p)/SL_2(\Z_p)
\end{equation*}
where the first map is a continuous surjective homomorphism,
and the quotient topology on $SL_2(\Q_p)$ is non-discrete ({\it
loc.~cit.} p. 116). Since the biggest normal subgroup of $SL_2(\Q_p)$
contained in $SL_2(\Z_p)$ is $\{\pm \id\}$, we see that its preimage in
$\pi_1^{\mathrm{et}}(\P^1,\ovl{s})$ cannot be open. Thus
$\pi_1^{\mathrm{et}}(\P^1,\ovl{s})$ is {\it not a pro-discrete group}.

\begin{cor}\label{cor::1.4.8} Let $\Cov^{\bullet\bullet}_S$ be a full
 subcategory of $\Cov^{\bullet}_S$ stable under taking connected
 components, fiber products (over $S$) and quotients. Then the
 corresponding continuous homomorphism
 $\pi_1^\bullet(S,\ovl{s})\to\pi_1^{\bullet\bullet}(S,\ovl{s})$ has
 dense image (it is neither strict\footnote{recall that a continuous
 homomorphism between topological groups is strict if the induced
 continuous bijective homomorphism between the coimage and the image is
 an isomorphism, i.e. bicontinuous; the coimage is the quotient of the
 source by the kernel, with the quotient topology.} nor surjective in
 general).  \hfill\break In particular, if
 $\;\Cov^{\mathrm{alg}}_S\subset \Cov^{\bullet}_S,\;$ then
 $\;\pi_1^{\mathrm{alg}}(S,\ovl{s})$ is the profinite
 completion\footnote{the profinite completion of a topological group is
 the inverse limit of its quotients by {\it open} normal subgroups of
 finite index.} of $\;\pi_1^\bullet(S,\ovl{s})$. 
\end{cor}

\begin{proof}
 By lemma \ref{lem::1.4.2}, it suffices to prove that for any $S'$ in
 $\Cov^{\bullet\bullet}_S$ and any $\ovl{s}'\in
 F^\bullet_{\ovl{s}}(S')=F^{\bullet\bullet}_{\ovl{s}}(S')$, the natural
 map $\pi_1^\bullet(S,\ovl{s})\to
 \pi_1^{\bullet\bullet}(S,\ovl{s})/\Stab^{\bullet\bullet}_{S',\ovl{s}'}$
 is surjective. We may replace $S'$ by the connected component of
 $s'$. According to the previous theorem,
 $\pi_1^{\bullet\bullet}(S,\ovl{s})/\Stab^{\bullet\bullet}_{S',\ovl{s}'}$
 and $\pi_1^{\bullet}(S,\ovl{s})/\Stab^{\bullet}_{S',\ovl{s}'}$ are then
 identified to
 $F^\bullet_{\ovl{s}}(S')=F^{\bullet\bullet}_{\ovl{s}}(S')$, whence the
 first assertion. The second assertion is clear, since
 $\;\pi_1^{\mathrm{alg}}(S,\ovl{s})$ is profinite.
 
 To see that the image of
 $\;\pi_1^\bullet(S,\ovl{s})\to\pi_1^{\bullet\bullet}(S,\ovl{s})$ is
 neither strict nor surjective in general, it suffices to take $S$ such that
 $\pi_1^{\mathrm{top}}(S,s)$ is infinite (\eg a Tate curve). Then
 $\pi_1^{\mathrm{et}}(S,\ovl{s})\to\pi_1^{\mathrm{alg}}(S,\ovl{s})$ is
 neither surjective nor strict. This follows from the commutative square 
 of topological groups
 \begin{equation*}
 \begin{CD}
  \pi_1^{\mathrm{et}}(S,\ovl{s}) @>>> \pi_1^{\mathrm{alg}}(S,\ovl{t})\\
  @VVV @VVV\\
  \pi_1^{\mathrm{top}}(S,s) @>>> \widehat{\pi_1^{\mathrm{top}}(S,t)}
 \end{CD}
 \end{equation*} 
 where the roof denotes the profinite completion; the coimage (\resp
 image) of the ``south-east map''
 $\pi_1^{\mathrm{top}} (S,\ovl{s}) \to  \widehat{\pi_1^{\mathrm{top}}(S,
 t)}$ is infinite discrete (\resp compact).
\end{proof}

\medskip In this context, the following general lemma is useful:

\begin{lem}\label{lem::1.4.9} Let $G, H$ be separated topological groups
 such that any neighborhood of $1$ contains an open 
 subgroup, and let $f: G\to H$ be a continuous homomorphism.
 \begin{enumerate}
  \item If the induced functor 
	$f^\ast:\,\text{$H$-$\Sets$}\to G$-$\Sets$ is essentially
	surjective, then $f$ is injective.
  \item If $f^\ast:\,H$-$\Sets \to G$-$\Sets$ is an equivalence, then $f$ is
	injective, strict and has dense image.
  \item If $f^\ast:\,H$-$\Sets\to G$-$\Sets$ is an equivalence and $G=
	\limproj G/U$ (the limit running over open subgroups of $G$),
	then $f$ is an isomorphism.
  \item If $G$ is complete and if the normal open subgroups are cofinal
	among the open subgroups $U$, then $G=\limproj G/U$.
 \end{enumerate}
\end{lem}
  
\begin{proof}
 (i) otherwise, let $g\neq 1$ be in the kernel of $f$. Because $G$ is
 separated, there exists a neighborhood $U$ of $1$ not containing $g$,
 and by hypothesis, one may assume that $U$ is an open subgroup of
 $G$. By assumption the (connected) $G$-set $G/U$ corresponds via
 $f^\ast$ to a (connected) $H$-set $X$. Since $f(g)=1$, $g$ acts 
 trivially on $f^\ast(X)$, in contradiction with the fact that $g\notin U$.

 (ii) To show that $f$ has dense image, it suffices to show that for $h\in
 H$ and any neighborhood $V$ of $1$, $f(G).h \cap V\neq \emptyset$. By
 hypothesis, one may assume that $V$ is an open 
 subgroup of $H$. Thus one has to show that $H=G.V$ for 
 any open subgroup $V$ of $H$, \ie $G$ acts transitively on $H/V$. If the
 action is not transitive, $f^{\ast}(H/V)$ can be 
 written as a disjoint union of $G$-sets, which contradicts the
 the assumption that $f^{\ast}$ is an equivalence. 

 \noindent It remains to show that $f$ is strict. But we have seen that
 $f^\ast$ induces an equivalence on transitive sets. 
 Therefore $f^{-1}$ induces a bijection between open subgroups of $H$
 and open subgroups of $G$, whence the result.

 (iii) Indeed one has a commutative diagram
 \begin{equation*}
 \begin{CD}
  G @>\subset>> H\\
  @V\wr VV @VVV \\
 \limproj G/U @>\sim>> \limproj H/V
 \end{CD}
 \end{equation*}
 and the right vertical map is injective since $H$ is separated.

 (iv) follows from \cite[III.7.2, prop. 2]{bourbaki74:_topol}.
\end{proof}

\begin{para}\label{para::1.4.10} Let $f:T\to S$ be any morphism of connected
 $K$-manifolds. We have the natural functor 

\begin{equation*}
 \Cov^{\mathrm{et}}_S\to \Cov^{\mathrm{et}}_T:\;\; S'\mapsto S'\times_S T.
\end{equation*}
Let us {\it assume} that it induces a functor 

\begin{equation*}
 f^{\bullet\ast}: \Cov^\bullet_S\to \Cov^\bullet_T.
\end{equation*}
Since
$F^\ast_{S,\ovl{s}}=F^\ast_{T,\ovl{t}}\circ f^{\bullet\ast}$ if
$\ovl{s}=f(\ovl{t})$, this gives rise to a continuous homomorphism

\begin{equation*}
f^\bullet_\ast :\pi_1^\bullet(T,\,\ovl{t})\to
 \pi_1^\bullet(S,\,\ovl{s}=f(\ovl{t})).
\end{equation*}
\end{para}

\begin{cor}\label{cor::1.4.11} $f^\bullet_\ast$ is an isomorphism of
 topological groups if and only if $f^{\bullet\ast}$ 
 is an equivalence of categories. 
\end{cor}

\begin{proof}
 If $f^\bullet_\ast$ is an isomorphism, it induces an equivalence 
 $\pi_1^\bullet(S,\,\ovl{s})$-$\Sets\cong
 \pi_1^\bullet(T,\,\ovl{t})$-$\Sets$, and it follows from \ref{thm::1.4.5} that
 $f^{\bullet\ast}$ is an equivalence of categories. The other
 implication is immediate.
\end{proof}

\begin{cor}\label{cor::1.4.12} Assume that $f$ is an \'etale covering in
 $\Cov^\bullet_S$, thus corresponding to the 
 $\pi_1^\bullet(S,\,\ovl{s})$-orbit
 $\pi_1^\bullet(S,\,\ovl{s})/\Stab^\bullet_{T,\ovl{t}}$ through the 
 dictionary of \ref{thm::1.4.5}.
 \begin{enumerate}
  \renewcommand{\theenumi}{\alph{enumi}}
%  \renewcommand{\labelenumi}{$(\theenumi)$}
  \item Then $\Stab^\bullet_{T,\ovl{t}}$ is the closure of the image of
	$f^\bullet_\ast$.  In particular, if $f$ is a Galois covering,
	its Galois group is
	$\;\pi_1^\bullet(S,\,\ovl{s})/(\Im\,f^\bullet_\ast)^-$.
  \item Assume moreover that for any covering $T'\to T$ in
	$\Cov^\bullet_T$, the composition with $f$ gives rise to a
	covering $T'\to S$ in $\Cov^\bullet_S$. Then $f^\bullet_\ast$
	is open and injective (hence $\;\pi_1^\bullet(T,\,\ovl{t})$ may
	be identified with the open and closed subgroup
	$\Stab^\bullet_{T,\ovl{t}}$ of $\pi_1^\bullet(S,\,\ovl{s})$).  
 \end{enumerate}
\end{cor}
\begin{proof}
 (a) Looking at the natural identifications 
 \begin{equation*}
 \pi_1^\bullet(S,\,\ovl{s})/\Stab^\bullet_{T,\ovl{t}}\cong
 F^\bullet_{\ovl{s}}(T)\cong F^\bullet_{\ovl{t},\ovl{t}}(T\times_S \,T)
 \end{equation*}
 of $\pi_1^\bullet(S,\,\ovl{s})$-orbits. The connected component of
 $T\times_S \,T$ containing $(\ovl{t},\ovl{t})$ is the 
 diagonal $T\subset T\times_S \,T$, therefore the
 $\pi_1^\bullet(T,\,\ovl{t})$-orbit is reduced to the point
 $(\ovl{t},\ovl{t})$. This shows that $\Im f^\bullet_\ast$ lies in
 $\Stab^\bullet_{T,\ovl{t}}$.

 To show that it is dense, we have to show that for any connected $S'$ in
 $\Cov^\bullet_S$ and any geometric point $\ovl{s}'$ of $S'$ above
 $\ovl{s}$, the map $\pi_1^\bullet(T,\,\ovl{t})\to
 \pi_1^\bullet(S,\,\ovl{s})/\Stab^\bullet_{T,\ovl{t}}$ induced by
 $f^\bullet_\ast $ is surjective. This follows from \ref{thm::1.4.5} as
 in \ref{cor::1.4.8}.

 (b) Under the extra assumption, one has functors $\Cov^\bullet_T\to
 \Cov^\bullet_S\to \Cov^\bullet_T: T'/T\mapsto T'/S ;\;S'/S\mapsto
 (S'\times_S \,T)/T$.

 \noindent For any open subgroup $H\subset  \pi_1^\bullet(T,\,\ovl{t})$
 (corresponding as in \ref{thm::1.4.5} to a pointed connected covering
 $(T',\ovl{t}')\;$), $\pi_1^\bullet(T,\,\ovl{t})/H\cong
 F^\bullet_{\ovl{t}' }(T' )$ is a subset of $
 F^\bullet_{\ovl{t}',\ovl{t}}(T'\times_S\,T)\cong 
 F^\bullet_{\ovl{s}}(T')\cong
 \pi_1^\bullet(S,\,\ovl{s})/\Stab^\bullet_{T',\ovl{t}'}$.
 In other words, $f^\bullet_\ast $ induces an open injective homomorphism 
 \begin{equation*}
 \pi_1^\bullet(T,\,\ovl{t})/H \hookrightarrow
 \pi_1^\bullet(S,\,\ovl{s})/f^\bullet_\ast(H).
 \end{equation*}
 Indeed, $f^\bullet_\ast(H)=\Stab^\bullet_{T',\ovl{t}'} $ is open in
 $\pi_1^\bullet(S,\,\ovl{s}) $: such groups correspond to connected
 coverings of $S$ which dominate $T$ (\ie which factorize over $f$).
 Since such coverings are cofinal in 
 $\Cov^\bullet_S$, we can pass to the limit, using \ref{lem::1.4.2}: 

 \begin{equation*}
 \pi_1^\bullet(T,\,\ovl{t}) = \limproj \pi_1^\bullet(T,\,\ovl{t})/H
 \hookrightarrow \limproj \pi_1^\bullet(S,\,\ovl{s}) /f^\bullet_\ast(H)
 \cong \pi_1^\bullet(S,\,\ovl{s})
 \end{equation*}
 and conclude that $f^\bullet_\ast$ is an open injective homomorphism.
\end{proof}

\subsection{Example: Ramero's locally algebraic \'etale fundamental
group.}\label{sub::1.5}

\begin{para}\label{para::1.5.1}
 When $S$ is countable at infinity, L. Ramero has introduced an
 interesting category $\Cov^{\mathrm{loc.alg}}_S$ of \'etale coverings, called
 `locally algebraic', which generalize the logarithmic covering of
 $\A^1$, and which are intimately connected with the local study of
 meromorphic differential equations around singularities \cite[4]{ramero98}.

 An \'etale covering is locally algebraic if the preimage of every {\it
 compact subdomain} $W\subset S$ is a disjoint union $\coprod W_i$ of
 subdomains $W_i$ which are finite \'etale over $W$. Since $S$ is (by
 assumption) a union of an increasing countable family of compact
 subdomains $S_j,\; j\geq 1$, it suffices to check the condition for 
 $W=S_j,\forall j.$ 

 \noindent This category is stable under taking connected components,
 fiber products (over $S$) and quotients. Moreover, it is stable by
 disjoint union and `composition'. Assuming that
 $s\in S_1$, one has a natural homomorphism  
 \begin{equation*}
 \liminj\pi_1^{\mathrm{loc.alg}}(S_j,\,\ovl{s}) = \liminj
  \pi_1^{\mathrm{alg}}(S_j,\,\ovl{s})\to \pi_1^{\mathrm{loc.alg}}(S,\,\ovl{s}).
 \end{equation*}
\end{para}

\begin{pro}\label{pro::1.5.2} This is an isomorphism:
 $\pi_1^{\mathrm{loc.alg}}(S,\,\ovl{s})\cong \liminj
 \pi_1^{\mathrm{alg}}(S_j,\,\ovl{s})$.
\end{pro}
   
\begin{proof}
 It is clear that an automorphism of $F^{\mathrm{loc.alg}}_{S,\ovl{s}}$
 amounts to a compatible family of automorphisms of
 $F^{\mathrm{alg}}_{S_j,\ovl{s}}$, hence the homomorphism is
 bijective. Moreover, the induced map $ \liminj
 \pi_1^{\mathrm{alg}}(S_j,\,\ovl{s})$-sets $\to
 \pi_1^{\mathrm{loc.alg}}(S,\,\ovl{s})$-sets is an equivalence: both
 sides are equivalent to $\Cov^{\mathrm{loc.alg}}_S$, \cf \ref{thm::1.4.5} and
 \cite[4.1.6]{ramero98}. Therefore open subgroups correspond to stabilizers as
 usual, and the bijection $\pi_1^{\mathrm{loc.alg}}(S,\,\ovl{s})\cong
 \liminj \pi_1^{\mathrm{alg}}(S_j,\,\ovl{s})$ is bicontinuous.
\end{proof}

\medskip\noindent It is not known whether
$\pi_1^{\mathrm{loc.alg}}(S,\,\ovl{s})$ is a pro-discrete group (\cf in
this direction {\it loc.~cit.} 4.3). 

\medskip\noindent {\it Example.} The logarithm of any Lubin-Tate formal
group over a $p$-adic field gives rise to a locally algebraic 
\'etale covering of the analytic affine line, \cf {\it loc.~cit.} 6.1.1.

On the other hand, the Gross-Hopkins covering of $\P^1$ is not locally
algebraic, although there is a finite affinoid 
cover $(V_i)$ of $\P^1$ such that the restriction of the covering to any
$V_i$ is a disjoint union of finite \'etale coverings.

\newpage
\section{Temperate fundamental groups.}\label{sec::2}

\begin{abst}
We introduce the notion of temperate \'etale
covering. Such coverings are essentially built from (possibly infinite)
topological coverings of finite \'etale coverings; they are classified
by temperate fundamental groups, which seem to be the right
non-archimedean equivalents of fundamental groups of complex
manifolds. These groups are not discrete, but nevertheless often possess
many infinite discrete quotients.
\end{abst}
   

\subsection{Temperate \'etale coverings and temperate fundamental
groups.}\label{sub::2.1}

\medskip In the archimedean context, the categories
$\Cov^{\mathrm{et}}_S$ and $\Cov^{\mathrm{top}}_S$ coincide; in the
non-archimedean situation, in contrast, $\Cov^{\mathrm{top}}_S$ is ``too
small'' (for instance, it consists of disjoint unions of copies of $S$
if $S$ has dimension one and tree-like reduction, \eg if $S$ is an
annulus), while  $\Cov^{\mathrm{et}}_S$ is ``too big'' (for $S=\P^1$, it
contains many non-trivial objects).

We introduce an intermediate category, which seems to be a closer
analogue of the category of unramified coverings in the 
archimedean context.
 
\begin{dfn}\label{def::2.1.1} An \'etale covering $S'\to S$ is {\rm
 temperate} if it is a quotient\footnote{in the sense of
 \ref{lem::1.2.7}.} of a 
 composite \'etale covering $T'\to T\to S$, where $T'\to T$ is a
 (possibly infinite) topological covering, and $T\to S$ is a finite
 \'etale covering.
\end{dfn}

Hence we have the commutative square 
\begin{equation*}
 \begin{matrix}
  &&T'&&\\
  & \swarrow &&\searrow &\\
  T&&&&S'\\
  &\searrow && \swarrow &\\
  &&S&&
 \end{matrix}
\end{equation*}
with $T'\in \Cov^{\mathrm{top}}_{T},\;T\in \Cov^{\mathrm{alg}}_{S},
S'\in \Cov^{\mathrm{et}}_S$. 

\medskip\noindent {\it Remark}. The notion of temperate covering is
stable under base change as in \ref{para::1.4.10}. 
 
\begin{lem}\label{lem::2.1.2} Let $T\to S$ be a finite Galois \'etale
 covering, and let $\tilde T\to T$ be the universal topological covering
 of $T$. Then the composite $\tilde T \to S$ is a Galois temperate
 covering. 
\end{lem}

\begin{proof}
 From \ref{lem::1.2.4}, we know that it is an \'etale covering, automatically
 temperate. It remains to show that it is Galois. By the universal
 property of $\tilde T$, any automorphism of $T$ lifts to $\tilde T$,
 hence there is an exact sequence of discrete groups:
 \begin{equation*}
 1\to \bigl(\pi_1^{\mathrm{top}}(T,t) \cong\bigr)\; \Aut_T \tilde T \to
 \Aut_S \tilde T \to \Aut_S\,  T\to 1. 
 \end{equation*}
 Since $\tilde T/(\Aut_T \tilde T) =T$ and $T/(\Aut_S T)=S,$ it
 follows that $ \tilde T/(\Aut_S \tilde T)=S$.
\end{proof}


We see from the lemma that in the definition of a temperate
covering, we may assume that $T'\to S$ and $T'\to T$ are Galois.  

\begin{rems}
 (i) It is {\it not true} that any quotient of a tower $T'\to T\to S$
 as in definition \ref{def::2.1.1} is again of this form. As a
 counter-example, let us consider a Tate elliptic curve $E$ over
 $K=\C_p$, and set $T= E\setminus E[2]$ (2-torsion points
 removed). Then $\tilde T $ is $\G_m^{\mathrm{an}}$ deprived from
 the countable subgroup $\pm (\sqrt q)^\Z$, and $T\cong \tilde
 T/q^\Z $. The inversion $x\mapsto 1/x$ on $\tilde T$ induces the
 inversion on $E$, and $\G_m/\{\id, \inv\}=\A^1 \;$ (via $x\mapsto
 x+1/x),\;E/\{\id,\inv\}=\P^1$. Let us consider the infinite
 temperate covering 

 \begin{equation*}
  S'=\tilde{T}/\{\id, \inv\} \to S= T/\{\id, \inv\}
   =\P^1\setminus\{0,1,\lambda,\infty\}.
 \end{equation*}
 Then $S'$ is not a topological covering of any finite \'etale
 covering $S''$ of $S$. Indeed, otherwise $\tilde T$ would be the
 universal topological covering of some connected component $T'$
 of $T\times_S S''$. The tower $\tilde T\to T' \to T$ would then
 extend to a tower of topological coverings $\G_m^{\mathrm{an}}\to
 E'\to E$, where $E'$ is another Tate curve isogenous to $E$ (we
 use the fact that the topological fundamental group of $T'$ and
 $E$ is the same). Then $S''\subset E'/\{\id, \inv\}\cong
 \P^1$. But any open subset in $\P^1_{\C_p}$ is simply-connected,
 whence $\tilde T= S''$, a contradiction.
 
 \noindent This example also shows that {\it a temperate covering may
 become a topological covering after finite \'etale base change, without
 being itself a topological covering}.

 (ii) The notion of temperate covering is not local on $S$. In fact,
 the Gross-Hopkins \'etale covering of
 $(\P^1_{\C_p})^{\mathrm{an}}$ is not temperate, although there is
 a finite open cover $(U_i)$ of $(\P^1_{\C_p})^{\mathrm{an}}$ such
 that it becomes a disjoint union of finite \'etale coverings over
 each $U_i$ (by compactness).
\end{rems}

\begin{para}\label{para::2.1.3} We denote the category of temperate
 (\'etale) coverings of $S$ by $\Cov^{\mathrm{temp}}_S$. 
 From the fact that topological and finite \'etale coverings
 respectively are stable under taking connected components and fiber
 products, it is not difficult to see that $\Cov^{\mathrm{temp}}_S$ is
 stable under taking connected components and fiber products (over
 $S$). By definition, it is stable under taking quotients. 
 
\medskip According to the results of 1.4, there is thus defined a
 separated totally discontinuous group, the {\it temperate fundamental group}
 $\pi_1^{\mathrm{temp}}(S,\ovl{s})$, which does not depend on the
 geometric point $\ovl{s}$ up to isomorphism. Moreover:
\end{para}

\begin{lem}\label{lem::2.1.4} $\pi_1^{\mathrm{temp}}(S,\ovl{s})$ is a
 pro-discrete group (hence complete, 
\cf \cite[III. 57.7]{bourbaki74:_topol}).  A basis of open neighborhoods of
$1$ is given by normal closed subgroups $H$ such that
 $\pi_1^{\mathrm{temp}}(S,\ovl{s})/H$ is the Galois group of the
 universal topological covering of some finite \'etale Galois covering of $S$.
\end{lem}
\begin{proof}
 This follows from \ref{cor::1.4.7}. and \ref{lem::2.1.2} (connected
 temperate coverings admit a Galois closure).
\end{proof}

 \medskip 
 The inclusion $\; \Cov^{\mathrm{top}}_S \hookrightarrow
 \Cov^{\mathrm{temp}}_S$ gives rise to a {\it surjective} homomorphism 
 \begin{equation*}
  \pi_1^{\mathrm{temp}}(S,\ovl{s}) \to \pi_1^{\mathrm{top}}(S,s) 
 \end{equation*} 
 (\cf \ref{cor::1.4.8}, taking
 into account the discreteness of $\pi_1^{\mathrm{top}}(S,s) $), which
 shows that $\pi_1^{\mathrm{temp}}(S,\ovl{s})$, although not a discrete 
 group, has many {\it infinite discrete quotients} in general; we say
 that $\pi_1^{\mathrm{temp}}(S,\ovl{s})$ is ``lacunary''.

 \medskip On the other hand, the inclusion $\; \Cov^{\mathrm{alg}}_S
 \hookrightarrow  \Cov^{\mathrm{temp}}_S$ gives rise to a 
 continuous homomorphism with dense image
 \begin{equation*}
  \pi_1^{\mathrm{temp}}(S,\ovl{s})\to \pi_1^{\mathrm{alg}}(S,\ovl{s}),
 \end{equation*}
 and $ \pi_1^{\mathrm{alg}}(S,\ovl{s}) $ may be identified with the
 profinite completion of $\pi_1^{\mathrm{temp}}(S,\ovl{s})$. This shows
 that $\pi_1^{\mathrm{temp}}(S,\ovl{s})$ is not ``too small'' (as opposed
 to $\pi_1^{\mathrm{top}}(S,s) $), and the following two propositions
 show that it is not ``too big'' as well (as opposed to
 $\pi_1^{\mathrm{et}}(S,\ovl{s}) $): 
 
\begin{pro}\label{pro::2.1.5} If $\dim\,S=1$, the homomorphism
 $\pi_1^{\mathrm{temp}}(S,\ovl{s})\to \pi_1^{\mathrm{alg}}(S,\ovl{s}) $
 is injective. 
\end{pro}

\begin{proof}
 By \ref{lem::2.1.2} and \ref{lem::2.1.4}, it suffices to see that if
 $T\to S$ is a finite 
 Galois \'etale covering, the Galois group of the temperate covering
 $\tilde T\to S$ is residually finite. But this group contains
 $\pi_1^{\mathrm{top}}(T,t)$ as a normal subgroup of finite index. The
 point is that in dimension one, topological fundamental groups are free,
 hence residually finite.
\end{proof}

\medskip The lacunary subgroup
$\pi_1^{\mathrm{temp}}(\P^1_{\C_p}\setminus \{0,1,\infty\})$ of the
profinite free group on two generators
$\pi_1^{\mathrm{alg}}(\P^1_{\C_p}\setminus \{0,1,\infty\})$ is
especially interesting and mysterious, and will be the 
object of our investigations in later sections\footnote{recall, on the
other hand, that $\pi_1^{\mathrm{top}}(\P^1_{\C_p}\setminus
\{0,1,\infty\})$ is trivial.}.

\begin{pro}\label{pro::2.1.6} If $S$ is algebraic and $K$ of
 characteristic zero, the Galois group of any Galois temperate 
 covering of $S$ is finitely generated. Moreover, 
 if $K$ has only countably many finite 
 extensions in a fixed algebraic closure $\ovl{K}$ (\eg if $K$ is a 
 $p$-adic field), then 
 $\pi_1^{\mathrm{temp}}(S,\ovl{s})$ and all its quotients by closed
 normal subgroups are `polish', \ie metrizable of countable type and
 complete.
\end{pro}

\begin{proof}
 Let us consider the first assertion. By \ref{lem::2.1.2}, it suffices to
 treat the case of the universal topological covering of a finite Galois
 \'etale covering.  By the already mentioned Gabber-L\"utkebohmert
 theorem, that finite \'etale covering is algebraic. The Galois group
 admits as a normal subgroup of finite index the topological fundamental
 group of that finite \'etale covering, which is finitely generated 
 (\cf \ref{para::1.1.2}); whence the result.

 \noindent Since $\pi_1^{\mathrm{temp}}(S,\ovl{s})$ is separated, its
 metrizability is equivalent to the existence of a countable fundamental
 system of neighborhood of $1$ \cite[IX.3.1 prop. 1]{bourbaki74:_topol}.
 By \ref{lem::2.1.4}, it 
 suffices to see that there are countably many finite \'etale coverings
 of $S$, 
% which follows from the topological finite generation of the
% algebraic fundamental group.
 which follows, by d\'evissage, from the assumption that
 $\Gal(\ovl{K}/K)$ has countably many open subgroups, and from the
 topological finite generation of the algebraic fundamental group over
 $\ovl{K}$.
 This also shows that
 $\pi_1^{\mathrm{temp}}(S,\ovl{s})$ is not only a pro-discrete group
 (which already implies that it is complete), but even a countable
 inverse limit of discrete countable groups. It follows that the
 metrizable group $\pi_1^{\mathrm{temp}}(S,\ovl{s})$ admits a 
 countable dense subset, hence its topology has a countable basis 
 \cite[IX.2.8 prop. 12]{bourbaki74:_topol}.

 \noindent It follows that any closed subgroup of
 $\pi_1^{\mathrm{temp}}(S,\ovl{s})$ and any quotient of
 $\pi_1^{\mathrm{temp}}(S,\ovl{s})$ by a 
 closed normal subgroup is metrizable and complete 
 \cite[IX.3.1 prop. 4]{bourbaki74:_topol}, 
 and of countable type (for a quotient: because it 
 has a countable dense subset).
\end{proof}

Notice however that the first assertion does not imply that
$\pi_1^{\mathrm{temp}}(S,\ovl{s})$ is topologically finitely generated. 
On the other hand, we shall show below that it is not locally compact in
general (\ref{pro::2.3.12}).

\begin{pro}\label{pro::2.1.7} Let $f: T\to S$ be an \'etale finite Galois
 covering of a connected $K$-manifold with group $G$, and $\ovl{t}$ be a
 geometric point of $T$. Then there is an exact sequence 
 \begin{equation*}
  1\to \pi_1^{\mathrm{temp}}(T,\ovl{t})\to
 \pi_1^{\mathrm{temp}}(S,f(\ovl{t}))\to G\to 1. 
 \end{equation*}
\end{pro}

\begin{proof}
 It is clear from the definition \ref{def::2.1.1} that any temperate covering
 of $T$ gives rise, by composition with $f$, to a 
 temperate covering of $S$. Thus the proposition follows from
 \ref{cor::1.4.12}.b.
\end{proof}

\medskip\noindent {\it Remark.} We shall see below (\ref{para::2.3.3})
that on the 
other hand, an \'etale covering of $S$ which is a finite \'etale
covering of a topological covering is not necessarily temperate.

\medskip The following result is an important criterion to recognize
topological, \resp temperate, \resp locally algebraic coverings: 

\begin{thm}\label{thm::2.1.8} Let $S'$ be a Galois \'etale covering of
 $S$ with group $G$. Then
 \begin{enumerate}
  \renewcommand{\theenumi}{\alph{enumi}}
%  \renewcommand{\labelenumi}{$(\theenumi)$}
  \item if $G$ is torsion-free,
	$S'$ is a topological covering of $S$;
  \item if $G$ is virtually torsion-free, $S'$ is a temperate covering
	of $S$; more precisely, $S'$ is a topological covering of some
	finite \'etale Galois covering of $S$;
  \item assume that $S$ is countable at infinity; then $S'$ is a locally
	algebraic covering of $S$ if and only if $G$ is the union of an
	increasing sequence of finite subgroups.
 \end{enumerate}
\end{thm}

\begin{proof}
 (a) (\cf also \cite{efgonas}) Let
 $\rho:\;\pi_1^{\mathrm{et}}(S,\ovl{s}) \to G$ be the 
 continuous homomorphism corresponding $S'\to S$. Let $\ovl{t}$ 
 be another geometric point of $S$, and let $\ovl{\gamma}$ be an \'etale
 path from $\ovl{t}$ to $\ovl{s}$.  Under the composed continuous homomorphism
 \begin{equation*}
  \pi_1^{\mathrm{et}}(t,\ovl{t})\to \pi_1^{\mathrm{et}}(S,\ovl{t})
   \xrightarrow{\mathrm{ad}(\ovl{\gamma})}\pi_1^{\mathrm{et}}(S,\ovl{s})\to G  
 \end{equation*}
 the compact (profinite) group $\pi_1^{\mathrm{et}}(t,\ovl{t})\cong
 \Gal(\mathscr{H}(t)^{sep}/\mathscr{H}(t))$ is sent to 
 $1$ if $G$ is torsion-free. This means that for any $\ovl{t}'\in
 F_{\ovl{t}}(S')$,  
 $\mathscr{H}(t')=\mathscr{H}(t)$. Let $U$ be an open neighborhood of $t$
 (which we may assume to be contractible thanks to Berkovich's theorem),
 such that $S'\times_S U\cong \coprod V_j$, with $V_j$ finite \'etale
 over $U$. Since $\mathscr{H}(t')=\mathscr{H}(t)$ for any point $t'$ in
 $V_j$ above $t$, it follows from \ref{para::1.2.1} that
 $V_j$ is a topological covering of
 the contractible open set $U$, hence a disjoint union of finitely many
 copies of $U$. This holds 
 for any point $t\in S$, therefore $S'\to S$ is a topological covering. 

 (b) If $G$ is virtually torsion-free, let $N$ be a torsion-free normal
 subgroup of finite index, and let $S'\to T\to S$ 
 be the corresponding intermediate Galois covering. By $a)$, $S'\to T$ is
 a topological covering, hence $S'\to S$ is a temperate \'etale covering.

 (c) Assume that $S'$ corresponds to a principal homogeneous space under
 $G_S$, classified by an element of $\mathrm{H}^1(S,G)$. Let us write
 $S=\bigcup S_j$, with $S_j$ compact. Assume that $G=\bigcup G_k$, with
 $G_k$ finite. Then $\mathrm{H}^1(S_j,G)= \liminj \mathrm{H}^1(S_j, G_k)$, therefore the
 restriction of $S'$ to any $S_j$ is a disjoint union of finite \'etale
 coverings of $S_j$, hence 
 $S'\to S$ is locally algebraic. Conversely, if $S'\to S$ is Galois
 locally algebraic, its group $G$ is a discrete image of $\liminj
 \pi_1^{\mathrm{alg}}(S_j,\ovl{s})$, each
 $G_j:=\Im(\pi_1^{\mathrm{alg}}(S_j,\ovl{s}))$ is finite, and $G=\bigcup
 G_j$.
\end{proof}  

\begin{cor}\label{cor::2.1.9} In dimension one, $
 \pi_1^{\mathrm{temp}}(S,\ovl{s})$ is the pro-VTFD-completion of
 $\pi_1^{\mathrm{et}}(S,\ovl{s})$. 
\end{cor}

By pro-VTFD-completion of a topological group $\Gamma$, we mean the
inverse limit of its virtually torsion-free discrete quotients. This is
a filtered limit (if $\Gamma/H'$ and $\Gamma/H''$ are VTFD, so is
$\Gamma/(H'\cap H'')\;$). Pro-VTFD-completions are functorial (due to
the fact that any subgroup of a VTF group is VTF).

\begin{proof}
 By \ref{lem::2.1.2}, the coverings of the form $\tilde T\to T \to S$, with $T$
 finite \'etale Galois over $S$, are cofinal among temperate
 coverings. In dimension one, we have seen that the Galois groups of such
 Galois coverings are virtually torsion-free, whence the result (using
 \ref{para::1.3.2}).
\end{proof}

\medskip\noindent {\it Examples.} (i) It follows from \ref{cor::2.1.9}
that the logarithmic (locally algebraic) covering
\begin{equation*}
\log:\; \matheur{D}_{\C_p}(1,1^-)\longrightarrow  \A^1_{\C_p} 
\end{equation*}
of the affine line is {\it not }temperate, since its Galois group
$\Q_p/\Z_p$ has no non-trivial virtually torsion-free 
quotient. This also follows from \ref{pro::2.1.5}, since
$\pi_1^{\mathrm{alg}}(\A^1_{\C_p})=0$. 

\medskip\noindent {\it Remark.} It is easy to draw from \ref{thm::2.1.8} that a
connected \'etale covering which is both temperate 
and locally algebraic is (finite) algebraic.

\medskip \noindent (ii) It also follows from \ref{thm::2.1.8} that there is no
\'etale Galois covering of the punctured disk with Galois group
$\Z$. This fact has noteworthy consequences on the local monodromy of
$p$-adic differential equations (\cf next sections).

{\small
\medskip\noindent {\it Remark.} In \cite{poaptg}, we have limited
ourselves to temperate coverings for which the tower $T'\to T\to S$ 
occuring in definition \ref{def::2.1.1} is Galois with VTF group. The
fundamental attached to the category of such coverings is 
then the pro-VTFD-completion of $\pi_1^{\mathrm{et}}(S,\ovl{s})$. The
definition \ref{def::2.1.1} seems however more natural (and is much more 
general in higher dimension). In \cite{poaptg}, another topological
group is introduced, namely the coimage of 
$\pi_1^{\mathrm{et}}(S,\ovl{s})\to \pi_1^{\mathrm{top}}(S,\ovl{s})\times
\pi_1^{\mathrm{alg}}(S,\ovl{s})$, denoted by
$\pi_1^{\mathrm{red}}(S,\ovl{s})$. In dimension one, there is a continuous
injective homomorphism with dense image $\pi_1^{\mathrm{red}}(S,\ovl{s})\to
\pi_1^{\mathrm{temp}}(S,\ovl{s})$, but it not clear whether this is an
isomorphism (the question is the openness of 
$\pi_1^{\mathrm{et}}(S,\ovl{s})\to 
\pi_1^{\mathrm{temp}}(S,\ovl{s})$).
}

% \medskip The following result is a kind of non-archimedean version of a
% theorem of Grauert-Remmert in complex geometry 
% (\cf \eg \cite[exp. XII, 5.4]{sga1}):
%
% \medskip\noindent
% There is a useful extension of the result in the case of mild
% singularities. Namely:
%
%\begin{cor}\label{cor::1.1.4} Let $\ovl{S}$ be a (paracompact strictly) 
% \emph{normal} connected $K$-analytic space, and let $Z$ be a
% Zariski-closed reduced analytic subset such that $S:=\ovl{S}\setminus
% Z$ is (rigid-analytically) smooth and non-empty. Assume that for any
% point $s$ of $\ovl{S}$, there is a finite (possibly ramified) covering 
% from a $K$-manifold onto a neighborhood of $s$.  Then any 
% topological covering of $S$ extends uniquely to a topological covering 
% of $\ovl{S}$.
%\end{cor}
%
%\begin{proof}
% Again, the problem is the existence of extensions, and this is a local
% problem. The problem is again  one extends first the pull-back
%\end{proof}


\begin{thm}\label{thm::2.1.10} Assume that $K$ is algebraically closed
 of characteristic zero. Let $\ovl{S}$ be  
 a (paracompact strictly) {\rm normal} connected $K$-analytic space
 $\ovl{S}$, and let $Z$ be a closed analytic subset such 
 that $S:=\ovl{S}\setminus Z$ is smooth.
 Assume that for any point $s$ of $\ovl{S}$, there is a covering (in the
 sense of \ref{def::1.2.3}) from a $K$-manifold 
 onto a neighborhood $U(s)$ of $s$ which is \'etale above $U(s)\cap S$. 

 \noindent
 Then any temperate \'etale covering $S'\to S $ extends to
 a (possibly ramified) covering  
 $\ovl{S}'\to \ovl{S}$, in the sense of \ref{def::1.2.3}: any point $x$
 has an open neighborhood $U(x)$ such that $S'\times_S U(x)$ 
 is a disjoint union of finite (ramified) coverings $V_j$ of
 $U(x)$. Moreover, one may assume that the $V_j$ are normal; 
 such an extension
 $\ovl{S}'\to \ovl{S}$ is then unique. 

 Furthermore, if $Z$ has codimension $\geq 2$ and $\ovl{S}$
 is smooth, any temperate
 \'etale covering $S'\to S $ extends to a unique temperate covering
 $\ovl{S}'\to \ovl{S}$.
\end{thm}

\begin{proof}
 We start with the statement about unicity. Since it is local on
 $\ovl{S}$, it is enough to establish it in the case of finite coverings
 (by definition \ref{def::1.2.3}). Let $\ovl{S}'$ and $\ovl{S}''$ be two
 connected 
 finite coverings of $\ovl{S}$ which coincide over $S$. The diagonal
 $\Delta$ is a connected component of $S'\times_S S''$. Its closure
 $\ovl{\Delta} $ in the normalization of $\ovl{S}'\times_{\ovl{S}}
 \ovl{S}''$ is again a connected component (using \cite[3.3.16]{spasalc}).
 The projections $\ovl{\Delta} \to \ovl{S}',\;\ovl{\Delta} \to
 \ovl{S}''$ are finite, and their restrictions to $\Delta$ are
 isomorphisms; hence they are themselves isomorphisms if $\ovl{S}'$ and
 $\ovl{S}''$ are normal.

 Let us turn to the existence of the extension (with the normality
 property). By unicity, this is a local question; we may and shall
 assume that $Z$ is a hypersurface, and that 
 $\ovl{S}$ is (strictly) affinoid, and that there exists a finite
 covering $\ovl{T}\to \ovl{S}$ by a 
 rigid-analytically smooth affinoid $\ovl{T}$, which restricts to a finite
 \'etale covering $T\to S$. 

 Let us first show that any quotient of a temperate covering $S'\to S$
 which has the extension property inherits this extension property: by
 unicity, the question is again local on $\ovl{S}$; by definition of
 coverings, this reduces the question to the case when $S'\to S$ is
 finite. By glueing (\cf \cite[1.3.3]{ecfnas}), it is again enough to treat
 the affinoid (normal) case. The argument is similar to the above one
 for unicity and left to the reader.
 
 This remark allows us to assume that $S'\to S$ is a topological covering of a
 finite \'etale covering of $S$. 
 Let us first treat the case of a finite \'etale covering. In this
 case, the extension property is proved in Gabber-L\"utkebohmert 
 \cite{repfapf}.
\end{proof}

The crucial case is the case of a punctured disk $S=\matheur{D}^{\ast}$
which is worth stating separately:

\begin{lem}[Gabber, L\"utkebohmert]\label{lem::2.1.11} Any finite
 \'etale covering of $\matheur{D}^\ast$ becomes a finite disjoint union
 of Kummer coverings on a smaller concentric punctured disk.
\end{lem}
\noindent (Here the assumption that $K$ is of characteristic zero is
essential). From there, one draws the extension property, in the product
case, when $\ovl{S}$ is a polydisk and $Z$ a union of axes (\cf {\it
loc.~cit.}, 3.3). By Kiehl's theorem on the existence of tubular
neighborhoods \cite{kiehl67:_rham_kohom_mannig_koerp}, the statement
also holds if $Z$ is a strict 
normal crossing divisor in $\ovl{S}$. In the general case, one can
invoke embedded resolution for $Z$, \cf \cite{schoutens99:_embed} (note that 
\cite[\S 4]{repfapf} proposes an alternative strategy which uses embedded
resolution only in dimension $2$).
 
It remains to deal at last with the case of a topological covering
$S'\to S$.  In order to avoid problems of singularities, we consider a finite 
etale covering $T\to S$ as above, $T$ being a rigid-analytically 
smooth affinoid.  It will be more convenient to replace $S'$ by 
$S'\times_S T$ (of which $S'$ is a quotient) so that the new $S'$ is 
a topological covering of $T$ 
itself a finite \'etale covering of $S$.
By \ref{pro::1.1.3}, $S'$ extends to a topological covering $\ovl{S}'
\to\ovl{T}$, and composition with $\ovl{T}\to \ovl{S}$ provides the 
desired extension of $S'/S$ to a covering $\ovl{S}'/\ovl{S}$.
 
The last statement of the theorem follows from the Nagata-Zariski purity
theorem (which becomes easy in the case of a finite morphism of
one-dimensional affinoids).

\medskip\noindent {\it Remark.} This extension property of temperate
coverings does not hold for arbitrary \'etale coverings. 

\noindent Indeed, one can show that the logarithmic covering of the
affine line does not extend to a (ramified) covering 
of the projective line.

\medskip The previous theorem is completed by the following result,
which is proven as in the complex case \cite[exp. XII, 5.3]{sga1}:

\begin{pro}\label{pro::2.1.12} Let $\ovl{S}$ be a normal analytic space,
 and let $Z$ be a closed analytic subset such that $S:=\ovl{S}\setminus
 Z$ is dense in $\ovl{S}$. Then the functor which attaches to any finite
 normal covering of $\ovl{S}$ its restriction to $S$ is fully faithful.
\end{pro}

It follows that in \ref{thm::2.1.10}, the extension of coverings from temperate
coverings of $S$ to `normal' coverings of $\ovl{S}$ `is' a fully
faithful functor.

%\newpage
\subsection{Interlude: tangential base-points.}\label{sub::2.2}

\begin{para}\label{para::2.2.1} Instead of basing fundamental groups at
 geometric points, it is sometimes more convenient\footnote{for instance
 in the study of Galois actions on fundamental groups.} to base them at 
 tangent vectors at points at infinity. 

 P. Deligne has explained the construction for affine algebraic curves in 
 characteristic zero, for topological, algebraic (= profinite) and other
 kinds of fundamental groups \cite{lgfdldpmtp}. In the profinite case, these
 tangential base-points amount to genuine geometric generic points,
 corresponding to embeddings of the function field of the curve into the
 field of Puiseux series at the chosen point at infinity (with respect
 to a fixed local parameter $z$ at that point; however, the embedding
 depends only on the tangent vector $\partial/\partial z$).

 If we try to imitate the construction in our non-archimedean setting, we
 face at once the following difficulty: let us assume for simplicity
 that the curve is the affine line; since there is no way to extend the
 Gauss norm on $K(z)$ to the field of Puiseux series with arbitrary coefficients in $K$, one does not get geometric (generic) points in the sense of \ref{para::1.2.2}.  We
 do not know how to overcome this difficulty for (analytic) \'etale coverings in
 general (and we doubt that this is possible). For temperate \'etale
 coverings, however, we can circumvent the difficulty as follows.
\end{para}

\begin{para}\label{para::2.2.2} We assume that $K$ is algebraically
 closed of characteristic $0$. Let $\ovl{S}$ be a $K$-manifold of
 dimension one, and let $S$ be the complement of a {\it classical} point
 $0$ in $\ovl{S}$ (so that $0$ has a basis of open neighborhoods
 consisting in disks centered at $0$). As above, it follows from
 \ref{lem::2.1.11} 
 that {\it the restriction of any temperate covering of $S$ to a
 punctured disk around $0$ of sufficiently small radius splits into a
 disjoint union of Kummer coverings of bounded degree}. If we replace
 $\ovl{S}$ by the its tangent space at $0$ and $0$ by the zero tangent
 vector, the same construction yields an equivalence of categories, with
 an obvious inverse given by `extension'. We thus get a {\it
 specialization functor}:
 \begin{equation*}
 \Cov^{\mathrm{temp}}_S\to  \Cov^{\mathrm{temp}}_{T_{\ovl{S},0}^0}.
 \end{equation*}         
 Any geometric point $\vec t$ of $T_{\ovl{S},0}^0$, 
 \eg any non-zero tangent vector at $0$, then gives rise by composition
 to a fiber functor 
 \begin{equation*}
  F^{\mathrm{temp}}_{\vec t}\;:\; \Cov^{\mathrm{temp}}_S\to  \Sets.
 \end{equation*}         
\end{para}

\begin{pro}\label{pro::2.2.3} This fiber functor is isomorphic to
 $F^{\mathrm{temp}}_{\ovl{t}}$ for any geometric point $\ovl{t}$.
\end{pro}
\begin{proof}
 We follow the line of proof of \ref{lem::1.3.3}. We may assume that
 $\ovl{S}$ is affinoid (hence compact). Let us consider finite open covers
 $\ovl{\mathcal{U}}=(\ovl{U}_i)_{i\geq 1}$ of $S$ {\it such that
 $\ovl{U}_1$ is a disk centered at $0$ and $\;0 \notin \ovl{U}_i\;$ 
 for $\;i>1$.}

 Let us denote by $ \mathcal{U}=(U_i)$ its trace on $S$, and by
 $\Cov^{\mathrm{temp}}_{S,\mathcal{U}}$ the full subcategory of 
 $\Cov^{\mathrm{temp}}_{S}$ of temperate coverings $f:S'\to S$ such that
 $f^{-1}U_i= \coprod V_{ij}$ and $V_{ij}$ is \'etale finite over 
 $U_i$ via $f$. 
 
 The point is that {\it by \ref{thm::2.1.10} (or \ref{para::2.2.2})}, we
 have
 $\displaystyle
 \Cov^{\mathrm{temp}}_{S}=\underset{\mathcal{U}}{\liminj} 
 \Cov^{\mathrm{temp}}_{S,\mathcal{U}}, $ where $\mathcal{U}$ runs among
 the open covers of $S$ satisfying the above properties. Also
 \begin{equation*}
  \Isom(F^{\mathrm{temp}}_{\vec t}, F^{\mathrm{temp}}_{\ovl{t}})
   =\underset{\mathcal{U}}{\limproj} \Isom(F^{\mathrm{temp}}_{\vec
   t,\mathcal{U}}, F^{\mathrm{temp}}_{\ovl{t},\mathcal{U}})
 \end{equation*}
 where $F^{\mathrm{temp}}_{\vec t,\mathcal{U}},
 F^{\mathrm{temp}}_{\ovl{t},\mathcal{U}}$ denote the restrictions of 
 $F^{\mathrm{temp}}_{\vec t }, F^{\mathrm{temp}}_{\ovl{t} }$ to
 $\Cov^{\mathrm{temp}}_{S,\mathcal{U}}$.

 To show that $\Isom(F^{\mathrm{temp}}_{\vec t},
 F^{\mathrm{temp}}_{\ovl{t}}) \neq \emptyset$, we construct 
 non-empty compact subsets $K_\mathcal{U}$ of
 $\Isom(F^{\mathrm{temp}}_{\vec t,\mathcal{U}},
 F^{\mathrm{temp}}_{\ovl{t},\mathcal{U}})$ (topologized 
 as before) such that if $\mathcal{U}'$ refines $\mathcal{U}$,
 $K_{\mathcal{U}'}$ maps to $K_\mathcal{U} $ under the natural map 
 $\Isom(F^{\mathrm{temp}}_{\vec t,\mathcal{U}'},
 F^{\mathrm{temp}}_{\ovl{t},\mathcal{U}'})\to
 \Isom(F^{\mathrm{temp}}_{\vec t,\mathcal{U}},
 F^{\mathrm{temp}}_{\ovl{t},\mathcal{U}})$; it will follow that
 $\Isom(F^{\mathrm{temp}}_{\vec t}, F^{\mathrm{temp}}_{\ovl{t}}) \supset
 \limproj_{\mathcal{U}}\, K_\mathcal{U}\neq 
 \emptyset$. 

 In fact, it suffices to do so for a cofinal system of
 $\mathcal{U}$'s. We now use a continuous embedding 
 $I=[0,1]\hookrightarrow \ovl{S}$ with $0\mapsto 0\in \ovl{S},\;1\mapsto t$. We
 consider open covers $\mathcal{U}  =(U_1,\ldots, U_n)$ as above, and
 such that
 \begin{equation*}
  U_1\cap I = 0,r_1[,\;\ldots, U_i\cap I=]r_i,t_i[,\;\ldots, U_m\cap
   I=]r_m,1],\;U_{m+1}\cap I=\emptyset,\ldots
 \end{equation*}
 for some points $r_i<t_{i-1}<r_{i+1}<t_i$ in $I \;(i=2,\ldots, m-1)$ and
 some $ m\leq n$.
 
 Given $\mathcal{U}$, one chooses $s_i$ in $U_i\cap U_{i+1}\cap I$ (which
 amounts to a real number between $r_{i+1}$ 
 and $t_i$), and a geometric point $\ovl{s}_i$ above $s_i\; (i=1,\ldots,
 m-1)$; one completes this collection by setting 
 $\ovl{s}_0=\vec t,\; \ovl{s}_m=\ovl{t}$. 

 Let us look at the diagram 
 \begin{equation*}
  \Cov^{\mathrm{temp}}_{S,\mathcal{U}}
   \xrightarrow{\sigma_i}
   \Ind\Cov^{\mathrm{alg}}_{\mathcal{U}_{i+1}}
   \begin{array}{c}
    \xrightarrow{F^{\mathrm{alg}}_{\ovl{s}_i, U_{i+1}}}\\
    \xrightarrow{F^{\mathrm{alg}}_{\ovl{s}_{i+1}, U_{i+1}}}\\
   \end{array}
   \longrightarrow \Sets.
 \end{equation*}               
 One has $F^{\mathrm{alg}}_{\ovl{s}_j, U_{i+1}}\circ \sigma_i =
 F^{\mathrm{temp}}_{\ovl{s}_i, \mathcal{U} },\;\;j=i,i+1$, whence
 continuous maps 
 \begin{equation*}
  \Isom(F^{\mathrm{alg}}_{\ovl{s}_i,
   U_{i+1}},F^{\mathrm{alg}}_{\ovl{s}_{i+1}, 
   U_{i+1}})\to \Isom(F^{\mathrm{temp}}_{\ovl{s}_i,
   \mathcal{U}},F^{\mathrm{temp}}_{\ovl{s}_{i+1}, \mathcal{U}}).
 \end{equation*} 
 But $\Isom(F^{\mathrm{alg}}_{\ovl{s}_i,
 U_{i+1}},F^{\mathrm{alg}}_{\ovl{s}_{i+1}, U_{i+1}})$ is a non-empty
 compact set for any $i=0,\ldots, m$: for $i=0\;$ ($s_0=\vec t$), one
 uses the analogue of \ref{pro::2.2.3} for the category of finite \'etale
 coverings, which follows from Deligne's discussion, {\it loc.~cit} (or
 even directly from \cite[V]{sga1}, interpreting
 $F^{\mathrm{alg}}_{\vec t, U_{i+1}}$ as a fiber functor at a geometric
 generic point). The rest of the proof goes exactly as in \ref{lem::1.3.3}.
\end{proof}

\begin{para}\label{para::2.2.4} This proposition allows to construct
 fundamental groups $\pi_1^\bullet(S,\,\vec t\,)$ based at a tangential
base-point $\vec t\,$ for any full subcategory $\Cov^\bullet_S\to
\Cov^{\mathrm{temp}}_S$ stable under taking connected components, fiber
products and quotients.  The theory of 1.4 extends without change to this
variant, as far as one restricts to curves over an algebraically closed
field of characteristic zero. For $\;\Cov^{\mathrm{top}}_S\to
\Cov^{\mathrm{temp}}_S$, the interpretation of tangential base-points is
not as transparent as in the complex-analytic context.
\end{para}

%\newpage
\subsection{Description of some temperate fundamental groups.}\label{sub::2.3}

 In this subsection, we assume that $K$ is {\it
 algebraically closed of characteristic $0$} (typically $K=\C_p$).

\begin{para}\label{para::2.3.1} It follows from \ref{pro::2.1.5}
 (together with \ref{para::1.4.1}) 
 that $\P^1$ and $\A^1$ have trivial temperate fundamental
 groups\footnote{in contrast, recall that $\A^1$ has non-trivial
 locally-algebraic fundamental group, and that $\P^1$ has trivial
 locally-algebraic fundamental group but non-trivial \'etale fundamental
 group.}.

Let us next examine $\G_m=\P^1\setminus\{0,\infty\}$. We have
$\pi_1^{\mathrm{top}}(\G_m,1)=0$; on the other hand, any connected finite
\'etale covering is Kummer $\G_m\to \G_m,\; x\mapsto x^n$. It follows that
\begin{equation*}
\pi_1^{\mathrm{temp}}(\G_m,1)=\pi_1^{\mathrm{alg}}(\G_m,1)=\widehat{\Z}(1):=
 \underset{n}{\limproj}\,\mu_n \cong \widehat{\Z}=\prod_\ell \Z_\ell.
\end{equation*}

\noindent Similarly, $\displaystyle\pi_1^{\mathrm{temp}}(\G_m^n,1)\cong
 \widehat{\Z}^n$. 
\end{para}

\begin{para}\label{para::2.3.2} {\it Elliptic curves.} Let us first consider
 the case of an {\it elliptic curve with good reduction} $S\;$ (so that
 $\pi_1^{\mathrm{top}}(S,1)=\{1\}$).  Then any finite
 \'etale covering is given by an isogeny
 $S'\to S$ and $S'$ has good reduction. Hence any temperate covering is
 an isogeny, and 
 \begin{equation*}
  \pi_1^{\mathrm{temp}}(S,1)=\pi_1^{\mathrm{alg}}(S,1)\cong \widehat{\Z}\times
   \widehat{\Z}.
 \end{equation*} 
 Let now $S=\G_m/q^\Z$ be instead a {\it Tate elliptic curve}. Then the
 homomorphisms 
 \begin{equation*}
  \pi_1^{\mathrm{temp}}(\G_m,1)\to \pi_1^{\mathrm{temp}}(S,1)\to q^\Z
 =\pi_1^{\mathrm{top}}(S,1)  
 \end{equation*}
 actually give rise to a canonical exact sequence
 (\ref{cor::1.4.12}. (b) may be applied here)
 \begin{equation*}
  1\to \widehat{\Z}(1)\to  \pi_1^{\mathrm{temp}}(S,1)\to  \Z \to 1,
\end{equation*}
 and since $ \pi_1^{\mathrm{temp}}(S,1)\subset
 \pi_1^{\mathrm{alg}}(S,1)\cong \widehat{\Z}\times \widehat{\Z}$ is
 abelian, to a non-canonical isomorphism 
 \begin{equation*}
  \pi_1^{\mathrm{temp}}(S,1)\cong \widehat{\Z}\times \Z.
 \end{equation*} 
 This is the simplest manifestation of the lacunary character of
 temperate fundamental groups.

 \medskip\noindent (The case of a Tate curve deprived from its origin is
 more difficult).
\end{para}

\begin{para}
 \label{para::2.3.3} {\it Mumford curves.} Let now $S$ be a Mumford
 curve of genus $g>1$:

 \noindent $S=\Omega/\Gamma,\;
 \;\Omega\cong \tilde S,\;$ and $\pi_1^{\mathrm{top}}(S,s)\cong \Gamma$,
 a Schottky group, free on $g$ generators.

 \noindent We have homomorphisms 
 \begin{equation*}
 \pi_1^{\mathrm{temp}}(\Omega,\bar{\tilde s})\to
  \pi_1^{\mathrm{temp}}(S,\ovl{s})\to  \pi_1^{\mathrm{top}}(S,s)
 \end{equation*}
 which give rise to an exact sequence  
 \begin{equation*}
 1\to H \to \pi_1^{\mathrm{temp}}(S,\ovl{s})\to 
 \pi_1^{\mathrm{top}}(S,s) \to 1 
 \end{equation*}
 with $H:= \Im(\pi_1^{\mathrm{temp}}(\Omega,\bar{\tilde
 s})\to \pi_1^{\mathrm{temp}}(S,\ovl{s}))^-\subset
 \pi_1^{\mathrm{temp}}(S,\ovl{s}) $. As an abstract group,
 $\pi_1^{\mathrm{temp}}(S,\ovl{s})$ is thus the semidirect product
 $H.\Gamma$.
 
 \noindent Passing to the profinite completion, we get an exact sequence
 \begin{equation*}
  \widehat{H} \to \pi_1^{\mathrm{alg}}(S,\ovl{s})\to 
 \widehat{\Gamma} \to 1, 
 \end{equation*}
 and the homomorphism
 $\pi_1^{\mathrm{alg}}(\Omega,\bar{\tilde s})\to
 \pi_1^{\mathrm{alg}}(S,\ovl{s})$ factors through 
 $\widehat{H}$. 

 \medskip This homomorphism is very far from being injective. The point
 is that the \'etale covering of $S$ obtained from an arbitrary finite
 Galois \'etale covering $\Omega'\to \Omega$ is not Galois (or, what
 amounts to the same by \ref{thm::2.1.8}, not temperate). The finite
 Galois \'etale coverings $\Omega'\to \Omega$ which give rise to a
 Galois covering of $S$ are called ``equivariant'' in \cite[1.3]{ecoamc}; 
 by \ref{thm::2.1.8}, they 
 correspond to a topological Galois covering of some Galois \'etale
 covering $S'\to S$.

 Moreover, there are examples \cite[2.7]{ecoamc} (one may choose for $\Omega$
 the Drinfeld half-plane) where the degrees of $\Omega'/
 \Omega$ and $S'/ S$ coincide, but $\Omega'$ is not the universal
 topological covering of $S'$ ; it follows that $\Omega'$ 
 is not simply-connected (in fact, its topological fundamental group is
 free of infinite type), and that the injection 
 $\pi_1^{\mathrm{temp}}(\Omega,\bar{\tilde s})\to
 \pi_1^{\mathrm{alg}}(\Omega,\bar{\tilde s})$ {\it is not a bijection in
 general}.

 Let us return to the non-injectivity of
 $\pi_1^{\mathrm{alg}}(\Omega,\bar{\tilde s})\to
 \pi_1^{\mathrm{alg}}(S,\ovl{s})$. Actually, the 
 following stronger result holds: {\it $\pi_1^{\mathrm{alg}}(S,\ovl{s})$
 is metrizable, while $\pi_1^{\mathrm{alg}}(\Omega,\bar{\tilde s})$ is 
 not}.  

 \noindent Recall that a profinite group is metrizable if and only if the
 set of its open subgroups is countable, 
 \cf \cite[1.3]{serre94:_cohom}. This is the case for
 $\pi_1^{\mathrm{alg}}(S,\ovl{s})$, which is even topologically finitely
 generated, but 
 not for $\pi_1^{\mathrm{alg}}(\Omega,\bar{\tilde s})$. In fact, one can
 show that for any 
 $n\geq 1$, the set of open subgroups $H\subset
 \pi_1^{\mathrm{alg}}(\Omega,\bar{\tilde s})$ endowed with an isomorphism
 $\pi_1^{\mathrm{alg}}(\Omega,\bar{\tilde s})/H\to \Z/n\Z$ is already
 uncountable. Indeed, this set may be identified with
 $\mathcal{O}(\Omega)^\times/(\mathcal{O}(\Omega)^\times)^n$, or else
 with the set of $\Z/n\Z$-valued 
 currents\footnote{\ie $\Z/n\Z$-valued functions on the edges of the
 tree, such that the sum of its values at the edges starting from an 
 arbitrary vertex is zero.} on the tree associated to
 $\Omega$ (\cf \cite[I.8.9, V.2.3]{garea}). Since $g\geq 2$, this 
 tree has infinitely many vertices connected to more than 
 two edges, and it follows that the $\Z/n\Z$-valued currents do not form
 a countable set.  

 A similar picture holds for Mumford curves deprived from a
 finite set of points.
\end{para}

\begin{para}\label{para::2.3.4} {\it Affinoid curves with good reduction.} In
 this case, one has 
 \begin{equation*}
  \pi_1^{\mathrm{temp}}(S,\ovl{s})=\pi_1^{\mathrm{alg}}(S,\ovl{s}) .
 \end{equation*} 
 We may assume that $s$ is the maximal point of $S$ corresponding to the
 sup-norm. The assertion then follows from De Jong's observation 
 \cite[7.5]{efgonas} that the map
 $\pi^{\mathrm{et}}_1(s,\ovl{s})=\pi^{\mathrm{alg}}_1(s,\ovl{s})\to
 \pi^{\mathrm{alg}}_1(S,\ovl{s})$ is already surjective. We shall only
 need the case of the closed unit disk. The proof goes 
 as follows (assuming that $S\subset \P^1$ as in {\it loc.~cit.}).

 One has to show that any finite \'etale covering $f:\;S'\to S$ of
 degree $n$ such that $S'_s= \{t_1,\ldots, t_n\}$ splits.  The
 assumption implies that $f$ is a local isomorphism around each $t_i$
 (\cf \ref{para::1.2.1}). Hence one can find an affinoid neighborhood $V$ of $s$ such
 that $S'_{| V}= W_1\coprod \ldots \coprod W_n$, and $W_i\cong
 V$. Because $V$ contains the maximal point $s$ which is the boundary of
 $S$ in $\P^1$ (in the sense of \cite[3.1]{staagonf}, we get an admissible
 covering (for the Grothendieck topology) of $\P^1= S \cup
 (\P^1\setminus(S\setminus V))$ by closed analytic subdomains. One then
 glues $S'$ and $n$ copies of $\P^1\setminus(S\setminus V)$ together via
 the isomorphisms $W_i\to V$ in order to get a finite \'etale covering of
 $\P^1$ which restricts to $S'$ over $S$. Such a covering is trivial.
 
 \medskip Concerning the algebraic fundamental group of the closed unit
 disk over $\C_p$, one has the following result, 
 essentially due to M. Raynaud (not used in the sequel):
\end{para}

\begin{thm}\label{thm::2.3.5} $\quad$ A finite group is a quotient of the
 profinite group $\;$ 
 $\pi_1^{\mathrm{alg}}\bigl(\matheur{D}_{\C_p}(0,1^+),0\bigr)$
 if and only if it is generated by its $p$-Sylow subgroups.
\end{thm}

 The necessity of the condition means that any finite Galois
 \'etale covering of $\matheur{D}_{\C_p}(0,1^+)$ of order prime to $p$
 is trivial, which is proven in \cite[6.3.3]{ecfnas} and
 \cite[2.11]{repfapf}.  The 
 converse is much harder and follows from Raynaud's solution of the
 Abhyankar conjecture on finite \'etale coverings of $\A^1_{\ovl{\F}_p}$,
 \cf \cite{rdldaecpecd}. Of course, we may replace $0$ by any geometric
point or tangential base-point (\cf \ref{pro::2.2.3}).

\medskip The standard example of a non-trivial finite Galois \'etale
 covering of $\matheur{D}(0,1^+)$ is the Artin-Schreier covering
 $\matheur{D}(0,1^+) \xrightarrow{z\mapsto z^p-z} \matheur{D}(0,1^+)$
 (with group $\Z/p$); it splits on any smaller concentric disk.
 A less standard 
 example, with non-solvable Galois group $A_5$, is given by $\matheur{D}(1,1^+)
 \setminus \matheur{D}(0,1^-) \xrightarrow{z\mapsto z^3-z^{-2}}
 \matheur{D}(0,1^+)$ for $p=3$.
 
\begin{para}\label{para::2.3.6} {\it Punctured disks.} Let us fix $r\in  {|
 K^\times|}$, and set $\; \matheur{D}_r=\matheur{D}_K(0,r^+),\;\;
 \matheur{D}_r^\ast=D_r\setminus \{0\}$.  We let $\ovl{s}$ be any
 geometric point or tangential base-point of $\matheur{D}^\ast$.
\end{para}

\begin{pro}\label{pro::2.3.7} There are canonical isomorphisms 
\begin{equation*}
\pi_1^{\mathrm{temp}}(\matheur{D}_r^\ast,\ovl{s})=
 \pi_1^{\mathrm{alg}}(\matheur{D}_r^\ast,\ovl{s})=
 \pi_1^{\mathrm{alg}}(\matheur{D}_r,\ovl{s})\times \widehat{\Z}(1).
\end{equation*}
\end{pro}

\begin{proof}
 By homothety, we may assume that $r=1$ (and drop the index
 $r$). These canonical isomorphisms are induced by the 
 obvious fully faithful functors (\cf \ref{cor::1.4.8}):
 \begin{equation*}
 \Cov^{\mathrm{kum}}_{\matheur{D}^\ast} \hookrightarrow
 \Cov^{\mathrm{alg}}_{\matheur{D}^\ast},\;\; 
 \Cov^{\mathrm{alg}}_{\matheur{D}} \hookrightarrow 
 \Cov^{\mathrm{alg}}_{\matheur{D}^\ast}\hookrightarrow
 \Cov^{\mathrm{temp}}_{\matheur{D}^\ast} 
 \end{equation*}
 where $\Cov^{\mathrm{kum}}_{\matheur{D}^\ast}$ denotes the category of
 finite disjoint unions of Kummer coverings of $\matheur{D}^\ast$; it is
 stable under taking connected 
 components, fiber products (the fiber product of a Kummer covering of
 degree $m$ and a Kummer covering of degree $n$ is a finite disjoint
 union of Kummer coverings of degree $\mathrm{lcm}(m,n)$) and quotients. The
 fundamental group $\pi_1^{\mathrm{kum}}(\matheur{D}^\ast,\ovl{s})$ is
 clearly $\Z(1)$.

 Assume for a while that $s$ is the maximal point of $\matheur{D}$ (the
 statement does not depend on the choice of $s$). Then as we have seen in
 \ref{para::2.3.4}, 
 the homomorphism $\pi_1^{\mathrm{temp}}(s,\ovl{s})\to
 \pi_1^{\mathrm{alg}}(\matheur{D},\ovl{s})$ is surjective. Since it
 factors through 
 $\pi_1^{\mathrm{temp}}(\matheur{D}^\ast,\ovl{s})$, it follows that
 $\pi_1^{\mathrm{temp}}(\matheur{D}^\ast,\ovl{s})\to\pi_1^{\mathrm{temp}}
 (\matheur{D},\ovl{s}) 
 = \pi_1^{\mathrm{alg}}(\matheur{D} ,\ovl{s})$ is surjective.
 
 Let us now take for $s$ a tangential base-point at $0$. Let us consider 
 the maps $\matheur{D}_r\to \matheur{D}$ (for $r\in {|K^\times|},\;r<1$) 
 and the corresponding homomorphism (\cf \ref{para::1.4.10}):
 \begin{equation*}
 \pi_1^{\mathrm{temp}}(\matheur{D}_r^\ast,\ovl{s})\to
 \pi_1^{\mathrm{temp}}(\matheur{D}^\ast,\ovl{s}).
 \end{equation*}
 It follows from \ref{lem::2.1.11} that 
 \begin{equation*}
 \underset{r\to 0}\limproj\,
  \Im(\pi_1^{\mathrm{temp}}(\matheur{D}_r^\ast,\ovl{s})\to
  \pi_1^{\mathrm{temp}}(\matheur{D}^\ast,\ovl{s})) = \Z(1),
 \end{equation*}
 the composition 

 \begin{equation*}
 \Z(1)\to \pi_1^{\mathrm{temp}}(\matheur{D}^\ast,\ovl{s})\to\Z(1)
 \end{equation*}
 being identity. It remains to show that the sequence 

 \begin{equation*}
 1\to \Z(1)\to \pi_1^{\mathrm{temp}}(\matheur{D}^\ast,\ovl{s}) \to
 \pi_1^{\mathrm{temp}}(\matheur{D} ,\ovl{s})=
 \pi_1^{\mathrm{alg}}(\matheur{D} ,\ovl{s})\to 1 
 \end{equation*}
 is exact. The problem is only in the middle; since
 $\pi_1^{\mathrm{temp}}(\matheur{D}^\ast,\ovl{s})$ is a pro-discrete group, it
 amounts to the following fact, which is a very special case of
 \ref{lem::2.1.11}: a 
 temperate Galois covering of $\matheur{D}^\ast$ extends to
 $\matheur{D}$ if and only if its 
 restriction to a sufficiently small punctured disk centered at $0$ splits.
\end{proof}

\medskip
On the other hand, the consideration of the logarithmic \'etale covering
restricted to a punctured disk centered at infinity shows that
$\pi_1^{\mathrm{loc.alg}}(\matheur{D}_r^\ast,\ovl{s})\neq
\pi_1^{\mathrm{alg}}(\matheur{D}_r^\ast,\ovl{s})$, and {\it a fortiori},
\begin{equation*}
 \pi_1^{\mathrm{et}}(\matheur{D}_r^\ast,\ovl{s})\neq
 \pi_1^{\mathrm{temp}}(\matheur{D}_r^\ast,\ovl{s}).
\end{equation*}
  
\begin{para}\label{para::2.3.8} {\it $\A^1$ minus a few points.}
 Let $S=\A^1\setminus\{\zeta_1,\ldots,\zeta_m\}$. For any $i=1,\ldots,
 m$, let $\vec t_i$ be a tangential base-point at $\zeta_i$, and let
 $\alpha_i$ be a temperate path from $\vec t_i$ to a fixed geometric
 point $\ovl{s} $ (or tangential base-point) of $ S$. Then we have a
 canonical homomorphism corresponding to Kummer coverings induced by
 $z-\zeta_i\mapsto (z-\zeta_i)^n$ (\cf \ref{para::2.3.6})

 \begin{equation*}
  \widehat{\Z}(1)\hookrightarrow \pi_1^{\mathrm{temp}}(S,\vec
  t_i)\to\widehat{\Z}(1)
 \end{equation*}
 whose composition is identity. 
 By composition with $\mathrm{ad}(\alpha_i)$, we get a monomorphism 
 \begin{equation*}
  \widehat{\Z}(1)\hookrightarrow \pi_1^{\mathrm{temp}}(S,\ovl{s}) 
 \end{equation*}
 which depends of course on $\alpha_i$.  Let us denote by $\gamma_i$ the
 image in $\pi_1^{\mathrm{temp}}(S,\ovl{s})$ of a topological generator
 of $\widehat{\Z}(1)\cong \widehat{\Z}$. The `local monodromies' $\gamma_i$ are
 not at all canonical, but the {\it smallest closed normal subgroup}
 $\an{\gamma_i}^-$ of $\pi_1^{\mathrm{temp}}(S,\ovl{s}) $ which contains
 them is independent of all choices. In fact we have\footnote{this is
 the correct version of cor. 5.3 in \cite{poaptg}, which was wrongly
 stated.}: 
\end{para}

\begin{pro}\label{pro::2.3.9}
 $\an{\gamma_i}^-=
 \pi_1^{\mathrm{temp}}(\A^1\setminus\{\zeta_1,\ldots, \zeta_m\},\ovl{s})$.
\end{pro}

\begin{proof}
 This amounts to saying that a temperate (Galois, if one wishes)
 covering splits if its restriction to sufficiently small punctured disks
 around any of the $\zeta_i$ splits. But the latter condition implies
 that the covering extends to a temperate covering of $\A^1$ 
 (\cf \ref{thm::2.1.10}), which automatically splits (\cf \ref{para::2.3.1}).
\end{proof}

\medskip By analogy with the complex situation, \ref{pro::2.3.9}
suggests two questions:

\begin{que}\label{que::2.3.10} Is it possible to choose the $\gamma_i$'s
 in such a way that they are topological generators of
 $\pi_1^{\mathrm{temp}}(S,\ovl{s})$ ? 
\end{que}

Following the above construction, one can introduce also introduce a
local monodromy $\gamma_\infty$ at $\infty$. 

\begin{que}\label{que::2.3.11} Is it possible to choose the $\gamma_i$'s
 and $\gamma_\infty$ in such a way that their product (in suitable
 order) is identity?
\end{que}

We do not know the answer to \ref{que::2.3.10}, but shall answer 
\ref{que::2.3.11} in \ref{sec::6}, when we have more tools at disposal.

For $m\geq 2$, $\pi_1^{\mathrm{temp}}(S,\ovl{s})$ is indeed a rather
complicated topological group, far from being discrete and from being compact:

\begin{pro}\label{pro::2.3.12} For $m\geq 2$,
 $\pi_1^{\mathrm{temp}}(\A^1\setminus\{\zeta_1,\ldots,
 \zeta_m\},\ovl{s}) $ is not locally compact.
\end{pro}

\begin{proof}
 Assume that $\pi_1^{\mathrm{temp}}(S,\ovl{s}) $ is
 locally compact. Then because $\pi_1^{\mathrm{temp}}(S,\ovl{s})$ is a
 prodiscrete group, there would exist a {\it compact} (necessarily
 totally discontinuous, hence profinite) open normal subgroup $H$. Let
 $S'\to S$ be the corresponding Galois temperate covering with group
 $G=\pi_1^{\mathrm{temp}}(S,\ovl{s})/H$. Up to replacing $H$ by a small
 subgroup, we may and shall assume that $S'=\tilde T$ is the topological
 universal covering of a finite Galois \'etale covering $T$ of $S$. Via
 \ref{pro::2.1.7}, we then have an exact sequence
 \begin{equation*}
 1\to H\to \pi_1^{\mathrm{temp}}(T,\ovl{t})\to \pi_1^{\mathrm{top}}(T,
 \ovl{t} )\to 1. 
 \end{equation*}
 Let $T'$ be any finite \'etale covering of $T$. Then the universal
 topological covering $\tilde T'$ is a temperate covering of $T$ which
 dominates $S'=\tilde T$. Due to the compactness of $H$, the map $\tilde 
 T' \to \tilde T$ must be finite. 

 We now choose $T'$ in order to contradict this property. Let $\ovl{X}$
 be a projective smooth analytic curve over $K$ such that $\rk
 \pi_1^{\mathrm{top}}(\ovl{X}) > \rk \pi_1^{\mathrm{top}}(T)$. Assume
 that $\ovl{X}$ is endowed with a finite morphism to $\P^1$ which is
 unramified above $S$ (since $m\geq 2$, this just amounts to saying that
 $\ovl{X}$ is defined over $\ovl{\Q}$, according to a well-known theorem
 of Belyi). We set $X=\ovl{X}\times_{\P^1} S \subset \ovl{X}$, and choose
 for $T'$ a connected component of $ X\times_S T$. Note that since $T$ is
 Galois over $S$, $T'$ is Galois over $X$. Let $Y$ be a connected
 component of $X\times_S \tilde T$ lying above $X$. Then $\tilde T' $ is
 a topological covering of $Y$, and is finite over $Y$ (if and) only if 
 it coincides with $Y$.  We then have a commutative square of Galois 
 temperate coverings
 \begin{equation*}
  \begin{matrix}
   && Y = \tilde T'&&\\
   &\swarrow &&\searrow &\\
   \tilde X &&&& T' \\
   &\searrow &&
   \swarrow &\\
   &&X&&
  \end{matrix}
 \end{equation*}
 Now $\Aut_X Y$ admits the free group $\Aut_{T'} Y$ as a normal subgroup
 of finite index, and admits the free group $\Aut_X \tilde X$ as a
 quotient. Since $\rank \Aut_{T'} Y \leq \rank \pi_1^{\mathrm{top}}(T)<
 \rank \pi_1^{\mathrm{top}}(\ovl{X})=\rank\,\pi_1^{\mathrm{top}}(X)$, this is
 impossible. Therefore $\pi_1^{\mathrm{temp}}(S) $ is not locally compact.  
\end{proof}

 \medskip\noindent {\it Remark.} In \cite{andre:_une_gal}, temperate fundamental
 groups of algebraic $p$-adic manifolds will be used to give a geometric 
 description of the local Galois group $\Gal(\ovl{\Q}_p/\Q_p)$. 
 
\newpage

\section{Local and global monodromy of $p$-adic differential
 equations.}\label{sec::3}

\markboth{\thechapter. $p$-ADIC ORBIFOLDS AND MONODROMY.}%
{\thesection. LOCAL AND GLOBAL MONODROMY.}

\begin{abst}
We briefly review and compare the theory of singularities of ordinary
linear differential equations in the formal, complex and $p$-adic
contexts. We then outline the Christol-Mebkhout theory of $p$-adic
slopes of differential modules over annuli. There is a close relation
with Galois representations of local fields of characteristic $p$,
expressed by Crew's $p$-adic local monodromy conjecture, which has been
recently proved \cite{andre01:_filtr_hasse_arf}. Complications occuring in the $p$-adic
context reflect the fact that the \'etale fundamental groups of annuli
are non-abelian.  \hfill\break We then define and study the
non-archimedean Riemann-Hilbert functor, which attaches a vector bundle
with integrable connection to any continuous representation with
discrete coimage of the \'etale fundamental group of a non-archimedean
manifold.  Connections in the image are characterized by the fact that
the \'etale sheaf of germs of solutions is locally constant.
\end{abst}


\subsection{Introduction.}\label{sub::3.1}

\begin{para}\label{para::3.1.1} Let us briefly review some well-known facts
 about singularities of ordinary linear differential equations. For
 concreteness, we consider a differential equation   

 \begin{equation}
  \alpha_\mu(z) \partial^\mu y + \cdots + \alpha_0 (z)
   y =0,\;\; \partial=z\frac{d}{dz}, \tag{$\ast$}
 \end{equation}
 with polynomial coefficients $\alpha_i(z)$, and concentrate on
 phenomena at infinity.
\end{para}

\begin{para}\label{para::3.1.2}  We begin the discussion with the formal
 setting. One defines the formal Newton polygon at $\infty$,
 $\widehat{NP}_\infty$ to be the convex hull of $\{(X\leq i,Y\geq \deg_z
 \alpha_0 -\deg_z \alpha_i )\}$ in the real 
 plane. Its `height'\footnote{the 
 height is the difference between the $Y$-coordinate of the highest 
 vertex and the $Y$-coordinate of the lowest vertex.}
 is called the (formal) irregularity of $(\ast)$ at
 $\infty$. The equation is called regular at $\infty$ if the
 irregularity is $0$, \ie if the only finite slope is $0$.

\medskip It is known that $(\ast)$ has a basis of `formal' solutions of
 the form 

\begin{equation*}
 \hat y=\sum {\hat u_i}\Bigl(\frac{1}{z}\Bigr) z^{-{\hat e}_i}(\log^{k_i}
  z) e^{P_i(z^{1/\mu !})},
\end{equation*}
 where the $P_i$'s are polynomials, and the ${\hat u_i}$ are just formal
 power series in $\frac{1}{z}$. The degree of the `ramified polynomials'
 $P_i(z^{1/\mu !}) $ actually corresponds to the slopes of $\widehat 
{NP}_\infty$. The ${\hat e}_i$'s are called the formal or Turrittin
 exponents; they can be computed algebraically.

 \medskip
 \noindent {\it Example.} Let us illustrate these notions with the {\it
 Bessel equation}
 \begin{equation*}
  (B)_{\nu}:\;\;\;\;(\partial^2  + (z^2  - \nu^2))y=0 
 \end{equation*}
 A basis of `formal solutions' at $\infty$ is given by
 \begin{equation*}
  \hat y_\pm= z^{-1/2}e^{\pm\sqrt{-1}z}.
   {}_2F_0\Bigl(\frac{\nu +1}{2},\frac{-\nu +1}{2}; \frac{\mp
   \sqrt{-1}}{2z}\Bigr), \quad {\hat e}_i=1/2,
 \end{equation*}
 and $\widehat{NP}_\infty$ is
 
%%% PICTURE 3a
 \begin{figure}[h]
  \begin{picture}(100,100)(0,0)
   \put(0,0){\line(1,0){40}}
   \put(40,0){\line(1,1){40}}
   \put(40,0){\makebox(0,0)[c]{{\small $\bullet$}}}
   \put(60,0){\makebox(0,0)[c]{{\small $\bullet$}}}
   \put(80,0){\makebox(0,0)[c]{{\small $\bullet$}}}
   \put(80,20){\makebox(0,0)[c]{{\small $\bullet$}}}
   \put(80,40){\makebox(0,0)[c]{{\small $\bullet$}}}
   \put(80,40){\line(0,1){40}}
  \end{picture}
  \caption{}
 \end{figure}
\end{para}

\begin{para}\label{para::3.1.3} {\it Complex-analytic setting}.  Solutions at
 any non-singular point $s$ (\ie a point in $\C^\times$ which is not a
 root of $\alpha_\mu$) converge up to the next singularity.

 Let ${\mathcal{C}}= \mathcal{C}(]r',r[) =\{z, \;r'<|z| < r\}$ be an
 annulus which contains no singularity. Using the fact the
 $\pi_1(\mathcal{C},s)\cong \Z$, it is classical to deduce that $(\ast)$
 has a basis of `analytic' solutions of the form

 \begin{equation*}
  y=\sum {u_i}\Bigl(\frac{1}{z}\Bigr) z^{-{e}_i} (\log^{k_i} z),
 \end{equation*}
 where the ${u_i}$ are Laurent series in $\frac{1}{z}$
 which converge in $\mathcal{C}$.
 The ${ e}_i$'s are called the analytic exponents.\footnote{in terms of
 differential systems, this amounts to the fact that a system
 $\partial Y = A Y$ on $\mathcal{C}$ has a fundamental solution matrix
 $Y$ of the form 
 $U.z^{-E}$, where $U$ is analytic in $\mathcal{C}$ and $E$ is constant (Fuchs
 normal form).}

 \medskip
 The relation between the $y$'s (analytic) and
 the $\hat y$'s (formal) is not straightforward. If $r=\infty$, the $\hat
 y$'s occur as asymptotic expansions of the $y$'s on suitable sectors,
 and the theory of multisummability/acceleration provides a kind of
 canonical inverse process (again on suitable sectors). If $\infty$ is a
 regular singularity, the $\hat u_i$ converge so that with $u_i=\hat u_i,
 e_i=\hat e_i$, the solutions $y$ and $\hat y$ coincide.

 \medskip
 \noindent {\it Example.} For the Bessel equation (say with non-integral
 parameter $\nu$ to simplify), one has a basis of solutions on
 $\mathcal{C}=\C^\times$ given by
  
\begin{equation*}
 y_\pm= z^{\pm\nu}\sum_0^\infty \;\frac{\;(-1)^n z^{2n }}{4^n n! (\pm\nu
  +1)_n}. 
\end{equation*}
 Moreover $e_i = \pm \nu$ (analytic exponents at $\infty = -$ analytic
 exponents at $0$); here we are `lucky', there are just two
 singularities, one of them is regular. In general, the computation of
 analytic exponents involves infinite determinants. Even in the
 so-called fuchsian case where there are only regular singularities, the
 computation of analytic exponents on an annulus which encloses several
 regular singularities is transcendental; e.g. the exponents of the
 differential equation of order two with solution
 ${}_2F_{1}(\frac{1}{2},\frac{1}{2},1;z^2)$, on an 
 annulus enclosing $0$ and $1$ but not $-1$ and
 $\infty$, are transcendental numbers:
 $\frac{\log(3\pm 2\sqrt 2)}{2\pi i}$\footnote{An independent
 solution is given by $i. \,{}_2F_{1}(\frac{1}{2}, \frac{1}{2}, 1;
 1-z^2)$, and in this basis, the local monodromy matrix 
 around $1$ (\resp $0$) is
 $\bigl(\begin{smallmatrix}1&0\\ -2&1\end{smallmatrix}\bigr)$
 (\resp $\bigl(\begin{smallmatrix}1&4\\0&1\end{smallmatrix}\bigr)$); the
 eigenvalues of the 
 monodromy around $\mathcal{C}$ are the eigenvalues of 
 $\bigl(\begin{smallmatrix}1&0\\-2&1\end{smallmatrix}\bigr)
  \bigl(\begin{smallmatrix}1&4\\0&1\end{smallmatrix}\bigr)
  =\bigl(\begin{smallmatrix}-7&4\\-2&1\end{smallmatrix}\bigr)$, 
 \ie $3\pm 2\sqrt 2$.}.
\end{para}

\begin{para}\label{para::3.1.4}  {\it $p$-adic-analytic setting}.  The first
 problem one encounters when dealing with $p$-adic differential
 equations (\cf I.\ref{sub-r-and-unit-root}) is that the solutions at
 non-singular points have small radii of convergence in general, and do
 not converge up to the next singularity.
 
 \medskip
 \noindent {\it Example.}
 $\frac{d}{dz}y=y,\;  e^z=\sum \frac{z^n}{n!} $ converges only in
 $\matheur{D}(0,|\pi|^-)$, $\pi^{p-1}=-p$. 

 \medskip
 This problem leads one to normalize somehow the situation by postulating 

 \medskip\noindent {\it Dwork's condition: the radius of convergence of
 solutions of $(\ast)$ at a generic point\footnote{a generic point of
 modulus $\rho$ in the sense of Dwork is a point $t_\rho$ in a complete
 extension $\Omega$ of $K$ with $| t_\rho|=\rho$ and such that the disk
 $\matheur{D}_\Omega(t_\rho,\rho^-)$ does not contain any element of $K$.}
 $\;t_1$ of modulus $1$ is $\geq 1$}.

 \medskip
 \noindent {\it Example.} One `normalizes' Bessel's equation by the
 change of variable $ z\mapsto 2\pi\sqrt{-1} z$. The normalized 
 Bessel differential operator is then   
 \begin{equation*}
  B_{\nu,\pi} = \partial^2  - (4\pi^2 z^2 +\nu^2)  
 \end{equation*}
 which is solvable in the $p$-adic generic disk whenever $\nu \in \Z_p$.
\end{para}

\begin{para}\label{para::3.1.5} Let us now discuss problems at $p$-adic
 singularities, concentrating again on phenomena at infinity.

 \medskip\noindent $\bullet$ The first one is the problem of `small
 divisors': if some formal exponent $\hat e_i\in \Z_p$ (or rather some
 difference $\hat e_i -\hat e_j$) is a Liouville number, \ie very
 closely approximated by infinitely many rational integers, the formal
 power series $\hat u_i$ occuring in the formal solutions $\hat y$ may
 diverge, even in the regular case. Otherwise, the $\hat u_i$ 
 converge, even in the irregular case \cite{baldassarri82:_differ}.

 \medskip
 \noindent {\it Example.} if $\nu \in \Z_p$ is well-approximable by
 negative integers, $\sum_0^\infty \frac{(-1)^n z^{2n}}{4^n n! (\nu
 +1)_n}$ may diverge, although $0$ is a regular singularity of the
 Bessel equation.

\medskip\noindent 
 $\bullet$ Even if Liouville numbers do not show up (\eg if everything
 is defined over $\ovl{\Q}$) so that the $\hat u_i$'s converge, these
 series may not converge as far as one might expect (\eg in a unit disk
 under the Dwork condition). This subtle phenomenon occurs in the
 normalized Bessel case $B_{\nu,\pi}$ for $p=2$ (${}_2F_0(\frac{\nu
 +1}{2},\frac{-\nu +1}{2}; \frac{\pm 1}{4\pi z})$ do not converge up
 to the unit circle): there is an annulus $\mathcal{C}(]1,r[)$ such that
 neither $y_\pm$ nor $\hat y_\pm$ provide any information about the
 solutions of $B_{\nu,\pi}$ in $\mathcal{C}(]1,r[)$.

 In particular, the question arises whether there is a basis of solutions 

 \begin{equation*}
  \sum{u_i}\Bigl(\frac{1}{z}\Bigr) z^{-{e}_i}(\log^{k_i} z),
 \end{equation*}
 where the ${u_i}$ are Laurent series in $\frac{1}{z}$ which converge in
 such an annulus $\mathcal{C}(]1,r[)$. More generally, concerning the
 existence of a `$p$-adic Fuchs normal form': 

 \medskip\noindent 
 $\bullet$ Does there exist a basis of solutions of $(\ast)$ of the form 
 \begin{equation*}
  y=\sum {u_i}\Bigl(\frac{1}{z}\Bigr) z^{-{e}_i} (\log^{k_i} z),
 \end{equation*}
 where the ${u_i}$ are Laurent series which converge in
 a given annulus $\mathcal{C}(]r',r[)$?

 \medskip\noindent A necessary condition has been put forward by
 P. Robba: {\it for any $\rho\in ]r',r[$, the radius of convergence of
 the solutions at the generic point $t_\rho$ of modulus $\rho$ is $\geq
 \rho$} (in the variable $z$, not $1/z$).

 The problems are: to which extent is the converse true? What can be
 said if Robba's condition is not fulfilled?

 \medskip\noindent These questions have recently become almost complete
 answers thanks to the Christol-Mebkhout theory of $p$-adic slopes of
 differential modules over annuli.  
\end{para}

\subsection{Local monodromy of $p$-adic differential equations. Outline
 of Christol-Mebkhout theory.}\label{sub::3.2}

\begin{para}\label{para::3.2.1} Let $K$ be a complete subfield of
 $\C_p$. We assume for simplicity that the restriction of the $p$-adic 
valuation to $K$ is discrete. 

The Robba $K$-algebra at `infinity' is 
 \begin{equation*}
\mathcal{R}=\mathcal{R}_{z,K}= \{K\text{-analytic functions on some }
 \mathcal{C}(]1,r[)\}.
\end{equation*}
 Endowed with $\partial = zd/dz$, this is a differential (integral)
 $K$-algebra. Although it is not simple as a differential algebra (there
 are many non-trivial differential ideals), it has the following useful
 property, which often allows one to replace $\mathcal{R}$ by its
 fraction field $Q(\mathcal{R})$ (\cf \cite{andre:_repres_bessel}.): 

\begin{quote}
 {\it any $\mathcal{R}[\partial]$-module $M$ which is of finite
 presentation over $\mathcal{R}$ is free over $\mathcal{R}$.} 
\end{quote}
 In the sequel, we simply refer to such an $M$ as a
 {\it differential module} over $\mathcal{R}$.
\end{para}

\begin{para}\label{para::3.2.2} For any $\rho >1$ sufficiently close to $1$,
 let us denote by $R(M,\rho)$ the infinimum of $\rho$ and the maximal
 radius $R$ such that $M$ is solvable (= has a basis of analytic
 solutions) in $\matheur{D}(t_\rho, R^-)$.

 One says that  $M$ is  {\it solvable} (at the inner boundary), if  
 $\displaystyle\lim_{\rho \to 1^+}R(M,\rho) = 1$ 

 \noindent(this is essentially Dwork's condition, except that $R(M,1)$
 itself is not defined.)
\end{para}

\begin{para}\label{para::3.2.3} A useful criterion (which goes back the
 Dwork's early studies) for $M$ being solvable is the existence 
 of a Frobenius structure. 

 Let $\sigma$ be an automorphism of $K$ inducing $t\mapsto
 t^{p^f}$ in characteristic $p$, and let $\varphi_f$ be the
 $\sigma$-linear endomorphism of $\mathcal{R}$ defined by $z\mapsto
 z^{p^f}$.

 One says that $M$ {\it has a Frobenius structure} if $\varphi_f^\ast
 M\cong M $. In what follows, the particular choice of such an
 isomorphism is irrelevant.

 \medskip\noindent
 {\it Example.} Let
 $M_{\nu,\pi}=\mathcal{R}/\mathcal{R}B_{\nu,\pi}$ be the differential
 module over $\mathcal{R}$ attached 
 to the normalized Bessel operator. It turns out that, up to isomorphism,
 $M_{\nu,\pi}$ does not depend on $\nu\in \Z_p$ 
 \cite{andre:_repres_bessel}. Choosing $\nu=0$ for instance, the existence of a
 Frobenius structure is known by the interpretation of $M_{\nu,\pi}$ in 
 terms of Kloosterman sums \cite{bfapfota} (in fact, for any $\nu \in \Q\cap
 \Z_p$, there is a relation to twisted Kloosterman sums, and a Frobenius
 structure has a geometric interpretation, \cf \cite{berthelot84:_cohom_dwork}).
\end{para}

\begin{para}\label{para::3.2.4} We know turn to $p$-adic slopes\footnote{which
 should not be confused with the so-called slopes of Frobenius which
 measure the $p$-adic size of a fixed Frobenius structure (when it
 exists); a unit-root $F$-crystal corresponding to the case of Frobenius
 slope zero.}, as defined by G. Christol and Z. Mebkhout.

\begin{itemize}
 \item A solvable $M$ is of {\it greatest slope} $\lambda \in
       [0,\infty[$ if $R(M,\rho) = \rho^{1-\lambda}$ for $\rho $
       close to $1$.
 \item $M$ is {\it purely of slope} $\lambda$ if any non-zero 
       solution at $t_\rho$ has radius of convergence $\rho^{1-\lambda}$ for
       $\rho $ close to $1$.
\end{itemize}
\end{para}

\begin{thm}\label{thm::3.2.5} Any solvable $M$ has a decreasing
 filtration $M_{>\lambda}$ such that $Gr_\lambda\; M$ is free and 
 purely of slope $\lambda$.  
\end{thm}
\par \noindent \cf \cite{sltdldedp3}. Following Christol-Mebkhout, one
defines the {\it $p$-adic Newton polygon} 
(at infinity) $NP_\infty$: the horizontal length of the segment of slope
$\lambda$ is  $rk \;Gr_\lambda\; M $.  In particular, the greatest slope
of $M$ in the above sense is the greatest slope of $NP_\infty$ in the
usual sense. The height of $NP_\infty$ is called the $p$-adic
irregularity of $M$. 

A solvable module $M$ is called {\it tame} (or $p$-adically regular) if
the only finite slope of $NP_\infty$ is zero. 
This amounts to say that $M$ satisfies the Robba condition of
\ref{para::3.1.5}: $R(M,\rho)=\rho$ when $\rho\to 1^+$.

\medskip One of the main results in the Christol-Mebkhout theory is
(\cite{sltdldedp3}): 
 
\begin{thm}\label{thm::3.2.6}
 $NP_\infty$ has integral vertices. 
\end{thm}
 
\par \noindent  A useful complement in the computation of $NP_\infty$
is given by (\cite{andre:_repres_bessel}):

\begin{lem}\label{lem::3.2.7} Assume $M$ comes from a differential module
 over the subring of $(K\cap \ovl{\Q})((\frac{1}{x}))$ consisting of
 convergent series in  $\mathcal{C}(]1,\infty[)$. Then the $p$-adic
 greatest slope $\lambda$ is $\leq$ than the greatest formal slope
 $\lambda_\infty$.
\end{lem}

 \noindent
 {\it Example.} It can be shown (using \ref{thm::3.2.8} below) that the
 Bessel module $M_{\nu,\pi}$ is not tame; \ref{thm::3.2.6} and
 \ref{lem::3.2.7} then leave only two possibilities for 
 $NP_\infty$, and it turns out that both occur  (\cf
 \cite{andre:_repres_bessel}).

%%% PICTURE 3b
 \begin{figure}[h]
  \begin{picture}(300,125)(0,-25)
   \put(0,0){\line(1,0){40}}
   \put(40,0){\line(2,1){40}}
   \put(40,0){\makebox(0,0)[c]{{\small $\bullet$}}}
   \put(60,0){\makebox(0,0)[c]{{\small $\bullet$}}}
   \put(80,0){\makebox(0,0)[c]{{\small $\bullet$}}}
   \put(80,20){\makebox(0,0)[c]{{\small $\bullet$}}}
   \put(80,20){\line(0,1){60}}
   \put(55,20){\makebox(0,0)[r]{{\footnotesize $\lambda=1/2$}}}
   \put(40,-20){\makebox(0,0)[c]{$p=2$}}

   \put(200,0){\line(1,0){40}}
   \put(240,0){\line(1,1){40}}
   \put(240,0){\makebox(0,0)[c]{{\small $\bullet$}}}
   \put(260,0){\makebox(0,0)[c]{{\small $\bullet$}}}
   \put(280,0){\makebox(0,0)[c]{{\small $\bullet$}}}
   \put(280,20){\makebox(0,0)[c]{{\small $\bullet$}}}
   \put(280,40){\makebox(0,0)[c]{{\small $\bullet$}}}
   \put(280,40){\line(0,1){40}}
   \put(255,25){\makebox(0,0)[r]{{\footnotesize $\lambda=1$}}}
   \put(240,-20){\makebox(0,0)[c]{$p>2$}}
  \end{picture}
  \caption{}
 \end{figure}
 \newpage
 \noindent
 As for the converse of Robba's criterion (\cf \ref{para::3.1.5}), one 
 has (\cf \cite{sltdldedp2}): 
\begin{thm}[$p$-adic local monodromy theorem, tame
 case]\label{thm::3.2.8}
 Assume $M$ has a Frobenius structure\footnote{there is a more 
 general version without Frobenius structure, in which one encounters 
 some trouble with the problem of small divisors. The existence of a
 Frobenius structure simplifies the statement, if not the proof, since
 it implies a priori that if $p$-adic exponents exist, they are
 rational, hence non-Liouville.}, and is purely of slope zero. Then 
 $M$ is an iterated extension of rank one differential modules of type
 $\partial-\alpha\;$, $\alpha\in \Q\cap \Z_p$ (called the $p$-{\rm adic
 exponents} of $M$).
 
 With respect to the conclusion about the rationality of the $p$-adic
 exponents, this result may be seen as a $p$-adic analogue of
 Grothendieck's $\ell$-adic local monodromy theorem for tame $\ell$-adic
 coefficients.
\end{thm}
The result, if not its proof, is closely related to the fact that a tame
(= of degree prime to $p$) Galois \'etale finite 
covering of any annulus $\mathcal{C}$ is Kummer (the tame algebraic
fundamental group of $\mathcal{C}$ is $\widehat{\Z}'(1)$, where the 
dash means as usual that the $p$-component is omitted,
\cf \cite{ecfnas}). The relationship will appear more clearly below.

\begin{para}\label{para::3.2.9} Let us return to our original differential
 equation $(\ast)$. Let us assume that the only singularities at finite
 distance are in $\matheur{D}(0,1^+)$.  Let us 
 also assume that $(\ast)$ satisfies Dwork's condition (so that the
 associated differential module over $\mathcal{R}$ is solvable), and that
 the formal exponents at $\infty$ are rational. Then the $p$-adic
 behavior of $(\ast)$ in a small punctured disk $\matheur{D}(\infty,
 \epsilon)^\ast$ centered at $\infty$ is dictated by the formal behavior
 (\cf \ref{para::3.1.5}), while the $p$-adic behavior in a thin annulus
 $\mathcal{C}(]1,r[)$ is predicted by Christol-Mebkhout theory. However,
 in the intermediate region $\mathcal{C}([r,1/\epsilon])$, nothing seems
 to be known.

 We shall return to this problem in 3.5. 
\end{para}

\subsection{Local monodromy of $p$-adic differential equations and
  Galois representations.}\label{sub::3.3}

\begin{para}\label{para::3.3.1} The differential sub-$K$-algebra 
 \begin{equation*}
  \mathcal{E}^{\dag} = \{\text{bounded $K$-analytic functions on some }
   \mathcal{C}(]1,r[)\}
 \end{equation*} 
 of bounded functions in $\mathcal{R}$ is actually a field. With the
 $p$-adic sup-norm, it is even a henselian field, with residue field
 $k((\frac{1}{z}))$ (where $k$ stands for the residue field of $K$).
 Hence any Galois extension $k((\frac{1}{z'}))/k((\frac{1}{z}))$ gives
 rise to an unramified Galois extension
 $\mathcal{E}^\dagger_{z',K'}/\mathcal{E}^\dagger_{z,K}$ (and, further,
 to a finite integral extension
 $\mathcal{R}_{z',K'}/\mathcal{R}_{z,K}$).
\end{para}

\begin{para}\label{para::3.3.2} In I.\ref{para-unit-root}, we have
 reviewed Katz' functor 
 from representations of fundamental groups of algebraic curves in
 characteristic $p$ to unit-root $F$-(iso)crystals. This construction has
 a local analogue, and one has the following theorem of N. Tsuzuki 
 \cite{tsuzuki98:_finit_f}:
\end{para}

\begin{thm}\label{thm::3.3.3} There is a natural equivalence of
 categories between the category of continuous $Q(W(k))$-linear 
 representations of
 $\Gal(k((\frac{1}{z}))^{\mathrm{sep}}/k((\frac{1}{z})))$ with 
 {\rm finite inertia} and the category of overconvergent unit-root
 $F$-isocrystals on $\Spec\;k((\frac{1}{z}))$.
\end{thm}

\par\noindent In particular, disregarding the precise Frobenius
structure, one gets a functor
\begin{equation*}
 \left\{\begin{matrix}
    \text{$p$-adic representations of $k((\frac{1}{z}))$}\\
    \text{with finite inertia}
   \end{matrix}
   \right\}
   \to
   \left\{\begin{matrix}
      \text{differential modules over $\mathcal{E}^\dagger$}\\
      \text{with Frobenius structure}
     \end{matrix}
     \right\}
\end{equation*}
hence, loosing even more structure, a functor

\begin{align*}
 &\bigl\{\begin{matrix}
	 \text{finite $p$-adic representations of $k((\frac{1}{z}))$}
	\end{matrix}\bigr\}\\
 &\qquad\qquad\qquad \to
 \left\{
 \begin{matrix}
  \text{differential modules over $ \mathcal{R}_{z,K}$ with Frobenius
  }\\
  \text{structure, which become trivial over some 
  }\\
  \text{finite extension $\mathcal{R}_{z',K'}$ of $\mathcal{R}_{z,K}$ as above}
 \end{matrix}
 \right\}
\end{align*}

\begin{dfn}\label{def::3.3.4} (R. Crew) A differential module $M$ over
$\mathcal{R}$ is {\rm quasi-unipotent} if there is a Galois extension
$k((\frac{1}{z'}))/k((\frac{1}{z}))$ such that after tensoring with the
corresponding finite extension $\mathcal{R}_{z',K'}$ of
$\mathcal{R}_{z,K} $, $M$ becomes an iterated extension of trivial
differential modules over $\mathcal{R}_{z',K'}$.

 There is an
 extensive work by Crew, S. Matsuda and Tsuzuki relating Galois
 representation theory and ramification theory of $k((\frac{1}{z})))$ and
 Christol-Mebkhout theory in the case of quasi-unipotent differential
 modules, \cf \cite{ceiatsc}.  In particular, they 
 establish a dictionary between Swan conductor and $p$-adic irregularity
 (in the above sense), break decomposition (in the sense of Katz) and
 $p$-adic slope decomposition... 
\end{dfn}

 \medskip\noindent {\it Examples.} (i) A {\it tame} module $M$
 with Frobenius structure is quasi-unipotent by \ref{thm::3.2.8}; indeed, it
 suffices to take a Kummer extension
 $k((\frac{1}{z^{1/n}}))/k((\frac{1}{z}))$. 

 (ii) $\mathcal{R}/\mathcal{R}(\partial-\pi z)$ is
 quasi-unipotent: take an Artin-Schreier extension defined by
 $z'-z^{\prime p}=z$.

 (iii) $p$-adic Bessel $M_{\nu,\pi}$ for odd $p$: a cyclic
 extension of $k((\frac{1}{z}))$ of degree $2p$ of 
 suffices. Indeed, it suffices to deal with the case $\nu = 1/2$ (\cf
\ref{para::3.2.3}). Then $(\frac{1}{z})^{\pm 1/2}.e^{\pm 2\pi z}$ is a
basis of solutions, and one reduces to the previous two examples.

\begin{coj}[Crew's $p$-adic local monodromy conjecture]\label{coj::3.3.5}
 Any differential module $M$ over $\mathcal{R}$, with Frobenius structure, is
 quasi-unipotent.
\end{coj}

This conjecture, which can be seen as a $p$-adic analogue of
Grothendieck's $\ell$-adic local monodromy theorem for wild $\ell$-adic
coefficients, and as a {\it non-abelian} extension of
\ref{thm::3.2.8} to the wild case, has now been solved
\cite{andre01:_filtr_hasse_arf}.
The proof relies heavily 
on theorems \ref{thm::3.2.5}, \ref{thm::3.2.6}, \ref{thm::3.2.8}.

\medskip\noindent A crucial test \cite{andre:_repres_bessel}, which is
at the origin 
of the solution, is the case of the {\it diadic} normalized Bessel
module $M_{\nu,2}$ (this test was suggested by Mebkhout as a potential
counterexample). This example provides a typical non-abelian instance of
\ref{coj::3.3.5}:

\begin{thm}\label{thm::3.3.6} For $K=\Q_4$, the normalized Bessel
 differential $\mathcal{E}^\dagger_{x,K}$-module $M$ is 
 quasi-unipotent.
\end{thm}

\medskip More precisely, let $E$ be the supersingular elliptic curve
 over $\F_4$ with affine equation 
  $Y^2+Y=X^3+\omega,\;$ ($\omega^3=1, \omega\neq 1$); then the covering
 $E\to E/\Aut(E)\cong \P^1\;$ is wildly totally ramified at the origin
 of $E$, and defines an unramified Galois extension of
 $\mathcal{E}^\dagger_{x,K}$ of degree $24$, with group $\Aut(E)\cong
 SL_2(\Z/3\Z)$, which contains all solutions of $M$. The underlying
 Galois representation of $G_{\F_4((1/z))}$ corresponds to the action of
 $\Aut(E)$ on $\mathrm{H}^1_{\mathrm{cris}}(E)$.


\subsection{A non-archimedean Riemann-Hilbert equivalence.}\label{sub::3.4} 

\begin{para}\label{para::3.4.1} Let us return to a general ordinary linear
 differential equation with polynomial or analytic coefficients 

 \begin{equation}
  \alpha_\mu(z) \partial^\mu y + \cdots + \alpha_0 (z) y = 0
  \;\text{ on $S\subset \A^1$, with $\partial=z\frac{d}{dz}$}.\tag{$\ast$}
 \end{equation}

 In the $p$-adic case, contrary to the complex situation, there
 does {\it not} exist, in general, any \'etale covering $S'\to S$ such
 that $(\ast)$ has a full set of solutions in $\mathcal{O}_{S'}$.
 
 \medskip\noindent {\it Example}: $\partial^2 y =0, \;S \supset
 \mathcal{C}(]1-\epsilon, 1[)$; there is no ``$p$-adic Riemann surface 
 of the logarithm'' on $S$.
 
 \medskip However, there are interesting global situations where this
 does happen. 

 \medskip\noindent {\it Example}: $\partial y - zy =0 $ on $S=\A^1$ :
 $S'=\matheur{D}(1,1^-)\xrightarrow{\log} S$, the infinite
 Galois \'etale covering with group $\mu_{p^\infty}\cong \Q_p/\Z_p$,
 trivializes the equation. 
 
 \medskip\noindent The object of this subsection is to investigate this
 phenomenon in full generality.
\end{para}

\begin{para}\label{para::3.4.2} Let $K$ be a complete non-archimedean
 field. We note that by definition, $K$-manifolds are `good' in the
 sense of \cite[1.2.6]{ecfnas}, \ie any point has an affinoid
 neighborhood. This property simplifies the theory of coherent sheaves
 (in the dictionary between rigid-analytic spaces and Berkovich analytic
 spaces, the theories of coherent sheaves and vector bundles correspond,
 \cf \cite[1.3.4, 1.6]{ecfnas})\footnote{Actually there are a number of
 technicalities involved in the analytic \'etale topology. First of all,
 since we restrict our attention to strictly analytic spaces in order to
 have the dictionary with rigid analytic varieties at disposal,
 Berkovich's definition of the \'etale site has to be slightly modified:
 this point is discussed in detail in R. Huber's book
 \cite[8.3]{ecoravaas}. Translated into rigid geometry, one gets the
 `partially proper \'etale' site of {\it loc.~cit.}. The associated
 sheaves correspond to the `overconvergent/conservative/constructible
 \'etale sheaves' of Van 
 der Put, Schneider, and De Jong \cite{jong96:_etale}, \cf
 \cite[8.2.12]{ecoravaas}. 
 We don't insist further on these subtleties.}.

 On the other hand, each coherent sheaf $M$ on $S$ gives rise to an
 \'etale sheaf $ M_{\mathrm{et}} $ on $S$ \cite[4.1.2]{ecfnas}.
 Moreover, one has \'etale descent for coherent sheaves: descent data are
 effective (\cf proof of \cite[4.1.10]{ecfnas}, and also
 \cite[3.2.3]{jong96:_etale} in the rigid setting). We shall use this only for
 locally free coherent sheaves (vector bundles) of rank $n$. These are
 in bijection with $\mathrm{H}^1_{\mathrm{et}}(S, GL_n)$.
  
\medskip In the sequel, $S$ will be a connected $K$-manifold, and $\ovl{s}$ a
geometric point of $S$. We assume that the point $s$ is $K$-rational,
 \ie is a classical point with residue field $K$.
\end{para}

\begin{para}\label{para::3.4.3} A continuous $K$-linear representation $\rho$
 of a topological group $\Gamma$ will be called {\it discrete} if its
 coimage $\Gamma/\ker \rho$ is discrete, or equivalently, if $\rho$
 factors through a discrete group.

 An {\it \'etale local system} on $S$ is a locally constant \'etale sheaf
 of $K$-vector spaces on $S$.
\end{para}

\begin{pro}\label{pro::3.4.4} There is a natural equivalence of categories   
\begin{equation*}
\{\text{discrete }\pi_1^{\mathrm{et}}(S,\ovl{s})\text{-representations }\}
 \xrightarrow{\sim}
 \{\text{\'etale local systems on $S$}\}.
\end{equation*}
\end{pro}

\begin{proof}
 Let $\rho $ be a discrete representation in a $K$-vector space
 $V$. This amounts to a continuous homomorphism $\pi_1^{\mathrm{et}}(S,
 \ovl{s}) \to \Aut_K V$, where $\Aut_K V$ denotes the discrete group of
 $K$-linear automorphisms of $V$.  On the other hand, any \'etale local
 system $\mathcal{V}$ on $S$ gives rise in the usual way to an \'etale 
 \v{C}ech cocycle with values in $\Aut_K V$, with $V=\mathcal{V}_s$, and
 its isomorphism class gives a well-defined element of 
 $\mathrm{H}^1_{\mathrm{et}}(S, \Aut_K V)$. Such a cohomology class may also be
 interpreted as usual as an $\Aut_K V$-principal homogeneous space
 $\mathcal{F}$ (as an \'etale sheaf) over $S$ up to isomorphism. What has
 to be shown is that such a principal homogeneous space is representable
 by an \'etale covering space. By lemma \ref{lem::1.2.6}. (ii), this is a local
 question on $S$. We thus may assume that $S$ is affinoid, and that for
 some finite \'etale covering $S'\to S$, $\mathcal{F}_{S'}$ is
 representable by an \'etale covering of $S'$. In that case, the
 representability of $\mathcal{F}$ follows from the descent theory of
 schemes.

 Starting from a representation $V$ of $\pi_1^{\mathrm{et}}(S,\ovl{s})$, we 
 have thus attached the $\Aut_K V$-principal homogeneous space and 
 etale sheaf $\FF$; the looked for \'etale local system is the 
 associated etale sheaf $\VV$ of $K$-vector spaces, with respect to 
 the canonical representation $V$ of $\Aut_K V$ (the pull-back of this 
 etale local system on $\FF$ is the constant sheaf with value $V$).
 
 Conversely, from an etale local system $\VV$, one attaches the $\Aut_K 
 V$-principal homogeneous space $\FF$ viewed as an \'etale covering 
 space (with $V=\VV_{\bar s}$); the structure of
 $\pi_1^{\mathrm{et}}(S,\ovl{s})$-set on its fiber at $\ovl{s}$ comes from
 a continuous homomorphism 
 $\pi_1^{\mathrm{et}}(S,\ovl{s})\to \Aut_K V$ (the cohomology class of
 $\FF$ in $\mathrm{H}^1(\pi_1^{\mathrm{et}}(S,\ovl{s}), \Aut_K V)$ is
 nothing but the latter homomorphism up to conjugation).
\end{proof}

In the sequel of this section, we assume that $K$ is of {\it
characteristic zero}.

\begin{para}\label{para::3.4.5} Let $M$ be a vector bundle (= locally free
 $\mathcal{O}_S$-module of finite rank) with a connection: $\nabla:
 M\to \Omega^1_S \otimes M$ (\ie a $K$-linear map satisfying Leibniz
 rule). This can also be interpreted $\mathcal{O}_S$-linearly, in terms
 of first-order jets, as a section $M\to \mathcal{P}^1(M)$ of the natural
 projection $\mathcal{P}^1(M)\to M$.

 \noindent The vector bundles with connection (\resp with integrable
 connection) form a tannakian category over $K$. The unit for $\otimes$
 is the trivial connection  $d: \mathcal{O}_S\to  \Omega^1_S $.

\medskip Because $s$ is a classical point with residue field $K$ (of
 characteristic zero), the stalk at $s$ of the sheaf $M^\nabla=\ker
 \nabla$ of germs of horizontal sections is a $K$-vector space of rank
 the rank of $M$, which can be identified with the fiber $M_{(s)}$ (note
 that since $S$ has a $K$-rational point, 
 $(\mathcal{O}_{S_{\mathrm{et}}})^{d_{S_{\mathrm{et}}}}=K$).

 \medskip To $M$ (\resp $\mathcal{P}^1(M)$), one attaches an \'etale sheaf
 $M_{\mathrm{et}}$ (\resp $\mathcal{P}^1(M)_{\mathrm{et}}$) --- locally free
 $\mathcal{O}_{S_{\mathrm{et}}}$-modules. Thus $\nabla$ also correspond
 to a section of the natural projection $\mathcal{P}^1(M)_{\mathrm{et}}\to 
 M_{\mathrm{et}}$, or equivalently, to a $K$-linear map
 $\nabla_{\mathrm{et}}: \;M_{\mathrm{et}}\to  \Omega^1_{S_{\mathrm{et}}}
 \otimes M_{\mathrm{et}}$. The kernel of the latter map is an \'etale
 sheaf of $K$-vector spaces denoted by $M_{\mathrm{et}}^\nabla$ (the
 \'etale sheaf of germs of horizontal sections). Let $t$ be point of
 $S$. The fiber of $M_{\mathrm{et}}^\nabla$ at $t$ (\cf
 \cite[4.2]{ecfnas} is a $\ovl{K}$-vector space endowed with a discrete
 $\Gal(\ovl{\mathscr{H}(t)}/{\mathscr{H}(t)})$-action. If
 $M_{\mathrm{et}}^\nabla$ is locally constant, the dimension of this
 vector space is independent of $t$; taking $t=s$, we find that it is
 the rank of $M$.  
\end{para}

 \medskip\noindent{\it Remark.} Any classical point admits a neighborhood
 over which $M^\nabla$ is constant of rank equal 
 to the rank of $M$ (by the non-archimedean analogue of Cauchy's theorem
 on solutions of analytic differential equations). 
 The constancy of $M_{\mathrm{et}}^\nabla$ imposes a constraint only at
 non-classical points of the Berkovich space $S$.  

\begin{thm}\label{thm::3.4.6} There is a natural equivalence of tannakian
 categories (the {\rm non-archimedean \'etale Riemann-Hilbert functor}):
 \begin{equation*}
  \left\{
   \begin{matrix}
    \text{discrete finite dimensional}\\
    \pi_1^{\mathrm{et}}(S, \ovl{s})\text{-representations}
   \end{matrix}
   \right\}
   \xrightarrow{RH^{\mathrm{et}}}
   \left\{
    \begin{matrix}
     \text{vector bundles with integrable}\\
     \text{connection $(M,\nabla)$ on $S$ such that}\\
     \text{$M_{\mathrm{et}}^\nabla$ is an \'etale local system}
    \end{matrix}
   \right\}.
 \end{equation*}
 In this correspondence, the subspace of 
 $\pi_1^{\mathrm{et}}(S, \ovl{s})$-invariants corresponds to the
 space of global sections of $M^\nabla$.
\end{thm}

\begin{proof}
 Using the previous proposition and \'etale descent of coherent
 modules, it suffices to show that the natural maps 
 \begin{equation*}
  \mathcal{V}\mapsto (\mathcal{V}\otimes_K \mathcal{O}_{S_{\mathrm{et}}},
   \id\otimes d_{S_{\mathrm{et}}}),\;\;\;(M,\nabla)\mapsto
   M_{\mathrm{et}}^\nabla,
 \end{equation*}
 induce quasi-inverse functors
 \begin{equation*}
  \left\{
   \begin{matrix}
    \text{locally constant \'etale sheaves}\\
    \text{of fin. dim. $K$-vector spaces}
   \end{matrix}
  \right\}
   \rightleftharpoons
   \left\{
    \begin{matrix}
     \text{vector bundles with integrable}\\
     \text{ connection $(M,\nabla)$ such that}\\
     \text{$M_{\mathrm{et}}^\nabla$ is an \'etale local system}
    \end{matrix}
   \right\}
 \end{equation*}
 It is clear that these maps are functorial (and the first one is clearly
 compatible with the tensor 
 structures). The point is to show that they are quasi-inverse. 
 
 On the one hand, it is clear that $\id\otimes d_{S_{\mathrm{et}}}$ is an
 integrable connection, and that $(\mathcal{V}\otimes_K
 \mathcal{O}_{S_{\mathrm{et}}} )^{\id\otimes
 d_{S_{\mathrm{et}}}}=\mathcal{V}$ (taking into account that
 $(\mathcal{O}_{S_{\mathrm{et}}} )^{d_{S_{\mathrm{et}}}}=K$).
 We have to show that if $(M,\nabla)  $ is a vector bundle with integrable
 connection such that $M_{\mathrm{et}}^\nabla$ is locally constant, then
 the natural horizontal map

 \begin{equation*}
 \phi: M_{\mathrm{et}}^{\nabla}\otimes_K \mathcal{O}_{S_{\mathrm{et}}}\to
 M_{\mathrm{et}}.
 \end{equation*}
 is an isomorphism. By the remark preceding the
 theorem, the stalks of $M_{\mathrm{et}}^\nabla$ are $\ovl{K}$-vector
 spaces of constant dimension, which is the rank of $M$. It follows that 
 $M_{\mathrm{et}}^\nabla\otimes_K \mathcal{O}_{S_{\mathrm{et}}}$ is the
 \'etale sheaf attached to a vector bundle of the same 
 rank as $M$. It is a general well-known fact about integrable
 connections\footnote{proven by passing to the completion around
 classical points.} that $\phi$ is a monomorphism of vector bundles, and
 that $\Im\phi$ is locally a direct summand of
 $M_{\mathrm{et}}$ (because the cokernel inherits an integrable
 connection, hence is locally free). To show that this it is an 
 epimorphism, it thus suffices to check that
 $M_{\mathrm{et}}^\nabla\otimes_K \mathcal{O}_{S_{\mathrm{et}}}$ and
 $M_{\mathrm{et}}$ have the same rank, which is the case by assumption.
\end{proof}

\begin{para}\label{para::3.4.7} The $ \pi_1^{\mathrm{et}}(S, \ovl{s})
 $-representation attached to $(M,\nabla)$ is called the {\it
mono\-dromy representation}. Note that the underlying vector space is
canonically identified with the fiber $M_s$. The image of the monodromy
representation is called the {\it monodromy group}. It may be
non-discrete (while the coimage is discrete by assumption); this
phenomenon is not specific to the non-archimedean situation, and is
indeed familiar in the complex case.
\end{para}

\begin{para}\label{para::3.4.8} The \'etale Riemann-Hilbert functor unifies
 several earlier constructions. Indeed, if the monodromy representation
 factors through $\pi_1^{\mathrm{top}}(S,s)$, we recover the
 `topological Riemann-Hilbert functor' considered in
 I.\ref{sub-connection}, whose image consists of connections with
 locally constant sheaves of solutions.

\noindent It also includes some cases of overconvergent unit-root
isocrystals with finite monodromy already considered in
I.\ref{sub-chiaroscuro} (in that case, the monodromy factors through a
discrete, hence finite, 
representation of $\pi_1^{\mathrm{alg}}(S, \ovl{s})$).

 Of course, connections in the image of the \'etale Riemann-Hilbert
 functor are still rather special\footnote{for instance, it does not
 account for the differential equation of the logarithm, as we 
 have said in \ref{para::3.4.1}; to see this, one can notice that otherwise the
 monodromy group would be contained in the additive 
 group of $K$, hence torsion-free, but this implies that the \'etale
 covering which trivializes the connection is a 
 topological covering (\ref{thm::2.1.8}), hence is trivial.}, but it
 seems difficult to extend its scope without invoking ``relative
 Fontaine rings''.
\end{para}

\begin{rems}
  (i) Whether $M_{\mathrm{et}}^\nabla$ is an
	\'etale local system or not can be checked locally on $S$. In 
	particular, the property for a connection to be in the essential image
	of $RH^{\mathrm{et}}$ is local. 

  (ii) The \'etale fundamental group $\pi_1^{\mathrm{et}}(\P^1_{\C_p},
	\ovl{s})$, although highly non-trivial, does not have any non-trivial
	discrete finite-dimensional representation. Indeed, 
	the corresponding connection on
	$\P^1_{\C_p}$ would be algebrizable, and necessarily trivial.
\end{rems} 


\subsection{Examples and applications.}\label{sub::3.5} 

\begin{para}\label{para::3.5.1} For any full subcategory $\Cov^\bullet_S$ of
 $\Cov^{\mathrm{et}}_S$ which is stable under 
 taking connected components, fiber products  and quotients, we have seen
 that there is a natural continuous homomorphism 
 \begin{equation*}
  \pi_1^{\mathrm{et}}(S, \ovl{s})\to \pi_1^\bullet(S, \ovl{s}) 
 \end{equation*}
 with dense image (\ref{cor::1.4.8}). Any discrete representation
 of $ \pi_1^\bullet(S, \ovl{s})$ gives rise to a 
 discrete representation of $ \pi_1^{\mathrm{et}}(S, \ovl{s})$ with the
 same coimage. It follows that the \'etale Riemann-Hilbert 
 functor induces a fully faithful functor 
 \begin{equation*}
  \left\{
   \begin{matrix}
    \text{discrete fin. dim.}\\
    \text{$\pi_1^\bullet(S,\ovl{s})$-representations}
   \end{matrix}
   \right\}
   \xrightarrow{RH^\bullet}
   \left\{
    \begin{matrix}
     \text{vector bundles with integrable}\\
     \text{connection $(M,\nabla)$ on $S$}
    \end{matrix}
   \right\}.
 \end{equation*}
 Note that the monodromy group defined in \ref{para::3.4.7} coincides
 with the image 
 of $\pi_1^\bullet(S, \ovl{s})$.  We shall discuss the avatars
 $RH^\bullet$ in the cases of $\Cov^{\mathrm{top}}_S$,
 $\Cov^{\mathrm{alg}}_S$, $\Cov^{\mathrm{loc.alg}}_S$,
 $\Cov^{\mathrm{temp}}_S$.
\end{para}

\begin{para}\label{para::3.5.2} The functor $RH^{\mathrm{top}}$ was studied in
 I.\ref{sub-connection} (in the case $K=\C_p$). Its essential image consists 
 of connections which satisfy `Cauchy's theorem' at every (possibly
 non-classical) point of $S$, or equivalently, which have a full set of
 multivalued analytic solutions, or else, such that 
 $M^\nabla$ is a locally constant sheaf on the topological space $S$.
\end{para}

\begin{para}\label{para::3.5.3} The essential image of the functor
 $RH^{\mathrm{alg}}$ consists of objects $(M,\nabla)$ such that
 $M_{\mathrm{et}}^\nabla$ becomes constant on some finite \'etale
 covering of $S$.

 Let us assume for instance that $S$ is affinoid with good reduction
 $S^0$, and that $K$ is a finite extension of
 $\widehat{\Q_p^{\mathrm{ur}}}$ as in 
 I.\ref{sub-r-and-unit-root}. Then among the finite representations of $
 \pi_1^{\mathrm{alg}}(S, \ovl{s}) $, one can distinguish those which
 factor through $ \pi_1^{\mathrm{alg}}(S_0, \ovl{s}_0) $. The
 corresponding connections underlie overconvergent unit-root isocrystals
 with finite global monodromy. We have seen a standard example in
 I.\ref{exa-dwork-exp}, given by the differential equation
 $\frac{d}{dz} y = -\pi y$ 
 on $S=\matheur{D}(0,1^+)$, which is trivialized over an Artin-Schreier covering.
 
 Let us now consider a case of bad reduction, namely the case of an
 annulus $S_r=\mathcal{C}(]1,r[),\;r\in | K|, 
 r>1$ over any discretely valued $p$-adic field $K$. Let $\mathcal{R}$ be
 the Robba ring at infinity. Then $RH^{\mathrm{alg}}$ induces a 
 functor
 \begin{equation*}
  \{\text{finite representations of $\pi_1^{\mathrm{alg}}(S_r,
   \ovl{s})$}\}
   \to \{\text{$\mathcal{R}$-differential modules}\}.
 \end{equation*}
 In relation to the local monodromy theorem \ref{coj::3.3.5}, it would be very
 interesting to understand the essential image of this functor as $r\to
1^+$. Note however that there is a difficulty with base-points in
passing to the limit $r\to 1^+$ (the problem disappears if one considers
only abelian representations).

Nevertheless, any $\mathcal{R}$-differential module $M$ extends to a
vector bundle $(M_r,\nabla_r)$ with connection on some $S_r$. Let us
assume that $M$ is semisimple and has a Frobenius structure.

If $M$ is tame, then \ref{thm::3.2.8} ensures that $M$ is in the image
of the above functor
\begin{equation*}
 \{\text{finite $\pi_1^{\mathrm{alg}}(S_r, \ovl{s})$-representations}
  \}\to \{\text{$\mathcal{R}$-differential modules}\}
\end{equation*}
for $r$ close enough to
$1$. The local monodromy theorem \ref{coj::3.3.5} implies the same if
$M$ is wild, 
although the monodromy representation might be non-abelian in this case,
as \ref{thm::3.3.6} shows.
\end{para}

\begin{para}\label{para::3.5.4} The essential image of the functor
 $RH^{\mathrm{loc.alg}}$ consists of objects $(M,\nabla)$ such that for every
 compact subdomain $W\subset S$, $(M_{\mathrm{et}}^\nabla)_{| W}$ becomes
 constant on some finite \'etale covering of $W$.
 Although not explicitly introduced in \cite{ramero98}, it is clear that
 these connections played an important heuristic role in that paper,
 especially those coming from abelian representations of the locally
 algebraic fundamental group (1.5).

 The basic example is $(\mathcal{O}_{\A^1}, d-1)$ over $\C_p$ (with
 solution $e^z$). It is trivialized by the logarithmic \'etale covering
 of ${\A^1}$.
 
 Similarly, for any ramified polynomial $P(z^{1/m})$, the differential
 equation with basic solution $e^{P(z^{1/m})}$ is trivialized over a
 locally algebraic \'etale covering of a sufficiently small punctured
 disk $\matheur{D}_{\C_p}(\infty,\epsilon)^\ast$ around $\infty$ (the
 logarithmic 
 covering of a Kummer covering). By \ref{para::3.1.5}, we see that any
 differential 
 equation with polynomial coefficients --- or more generally meromorphic
 singularity at $\infty\;$ --- with {\it rational} formal exponents at
 $\infty$, and with no logarithm in the formal solutions $\hat y$ (\cf
 \ref{para::3.1.2}), corresponds to an object in the image of the
 locally algebraic Riemann-Hilbert functor
 \begin{equation*}
  \left\{
   \begin{matrix}
    \text{finite dimensional}\\
    \text{$\pi_1^{\mathrm{loc.alg}}(\matheur{D}(\infty,\epsilon)^\ast,
    \ovl{s})$-representations}
   \end{matrix}
  \right\}
  \to
  \{\text{$\mathcal{O}(\matheur{D}(\infty,\epsilon)^\ast)$-differential
  modules}\}
 \end{equation*}
 for $\epsilon$ small enough.

 There is a certain analogy between the logarithmic \'etale
 covering of $\matheur{D}_{\C_p}(\infty,\epsilon)^\ast$ and the {\it
 real blow-up} of $\matheur{D}_{\C}(\infty,\epsilon)^\ast$ which plays
 such an important role in 
 the study of irregular singularities of complex differential
 equations. In both cases, the point at infinity becomes a limiting
 circle\footnote{L. Ramero points out that the change from analytic to
 \'etale topology is essential in the $p$-adic case, while in the
 complex case, the real blow-up corresponds to completion for a
 different uniform structure of the punctured disk with no change of the
 topology.}.

\medskip Let us return to a general vector bundle $(M,\nabla)$ with
 connection on $\matheur{D}(\infty,1^-)^\ast$ which extends to a vector bundle
 with meromorphic connection at $\infty$. We assume that it is has a
 Frobenius structure near the unit circle, and that the formal exponents
 at infinity are rational. Regarding problem \ref{para::3.2.9}, a
 natural question is:
\end{para}

\begin{que}\label{que::3.5.5} Under which further conditions, if any,
 does there exist a locally algebraic covering of
 $\matheur{D}(\infty,1^-)^\ast$ 
 over which the connection becomes an iterated extension of trivial
 connections?
\end{que}
At any rate, in the cases where this holds, this would let Ramero's
locally algebraic fundamental group play the same role, 
over $\C_p$, as Ramis' wild fundamental group (over $\C$) in the local
study of meromorphic irregular singularities.

\begin{para}\label{para::3.5.6} The essential image of the functor
 $RH^{\mathrm{temp}}$ consists of objects $(M,\nabla)$ such that 
 $M_{\mathrm{et}}^\nabla$ becomes constant on some temperate \'etale
 covering of $S$.  This implies that there is a finite \'etale Galois
covering $f:\;T\to S$ such that the pull-back
$(M',\nabla')=f^\ast(M,\nabla)$ has a full set of multivalued analytic
solutions, \ie $M^{\prime \nabla'}$ is locally constant. Let $G$ be the group
$\Aut_S\,T$, and set $n=| G| $.
 
Assume that there are $n$ distinct $K$-rational points $t_i$ in $T$
 lying above $s$. Then the monodromy group admits the following {\it
 concrete interpretation}: 

\noindent let $\alpha_{ij}$ be a Berkovich path (up to homotopy)
connecting $t_i$ to $t_j$. Then $f_\ast(\alpha_{ij})\in
\pi_1^{\mathrm{temp}}(S,\ovl{s})$. When $i,j$ and $\alpha_{ij}$ vary,
they generate the image $\Aut_S\,\tilde T$ of
$\pi_1^{\mathrm{temp}}(S,\ovl{s})$ (note that, varying $j$, they
generate the finite quotient $G$). It follows that the monodromy group
consists of the images of the $f_\ast(\alpha_{ij})$'s in $GL(M_{(s)})$
by the monodromy representation.

Therefore, except for the occurrence of the finite \'etale covering, the 
picture is quite similar to the complex situation.
 

\end{para}

\begin{pro}\label{pro::3.5.7} Let $S=\A^1\setminus\{\zeta_1,\ldots,
 \zeta_m\}$. Assume\footnote{this very strong assumption is rarely met
 in practice.} that the residue characteristic $p$ of $K$ does not
 divide $n$. Then:
 \begin{enumerate}
  \renewcommand{\theenumi}{\alph{enumi}}
  \item for any open or closed disk $\matheur{D}\subset S$, the restriction of
	$\nabla$ to $\matheur{D}$ is trivial (\ie has a full set of solutions); in
	particular, $\nabla$ satisfies the Robba condition;
  \item for any open or closed annulus
	$\mathcal{C}\subset S$, the pull-back of $\nabla$ to the Kummer
	covering of $\mathcal{C}$ of degree $n$ is trivial.
 \end{enumerate}

 In particular,
 $p$-adic exponents in the sense of Christol-Mebkhout exist (in any
 annulus which does not contain any $\zeta_i$), and they are rational
 numbers with denominator $n$ (or a divisor of $n$).
\end{pro}

\begin{proof}
 One is immediately reduced to considering closed disks and
 annuli. We know the pull-back of $\nabla$ to $\tilde T$ is trivial. Let
 $\matheur{D}'$ (\resp $\mathcal{C}'$) be a connected component of $\matheur{D}\times_S
 \tilde T$ (\resp $\mathcal{C}\times_S \tilde T$). Because $n$ is assumed
 to be prime to $p$, the covering $\matheur{D}'\to \matheur{D}$ is trivial (\resp is a Kummer
 covering of degree dividing $n$), \cf \cite[6.3.3, 6.3.5]{ecfnas} (taking
 into account the fact that disks and annuli are simply-connected). The
 assertions follow.
\end{proof}

\medskip\noindent {\it Example.} (see also the example at the end of
\ref{lem::2.1.2}). Let $K$ be a $p$-adic field for any odd $p$. 
Let $T=E\setminus E[2]$ be a Tate curve deprived from its four
$2$-torsion points (assumed to be $K$-rational), written in 
Legendre form: 
\begin{equation*}
 y^2=x(x-1)(x-\lambda),\quad |\lambda(\lambda -1)|\neq 1.
\end{equation*}
Let $f$ be the projection from $T$ to $S=\A^1\setminus
\{0,1,\lambda \}$.
Here $\tilde T$ is of the form $\G_m\setminus \pm(\sqrt q)^\Z$, and
$T\cong \tilde T/ q^\Z $.

The sum of the canonical one-dimensional $K$-linear representation of
$q^\Z$ ($q$ acting by multiplication by a factor $q$) and its dual
extends to a representation of the semi-direct product of $q^\Z$ and
$\Z/2\Z$ (which is the Galois group of the covering $\tilde T\to S$). On
can choose the local monodromies at $0,\lambda, 1,\infty$ respectively
in such a way that they correspond to the matrices
\begin{equation*}
 \begin{pmatrix}0&1\\1&0\end{pmatrix},
 \begin{pmatrix}0&1\\1&0\end{pmatrix},
  \begin{pmatrix}0&q\\q^{-1}&0\end{pmatrix},
  \begin{pmatrix}0& q^{-1}\\q&0\end{pmatrix}
\end{equation*}
respectively. It is
then easy to check that the monodromy around an annulus surrounding $0$
and $\lambda$ alone ($|\lambda|$ being assumed $<1$ for convenience) is
trivial: the $p$-adic exponents are zero.
 
\begin{para}\label{para::3.5.8} Much more interesting examples arise in the
 context of $p$-adic period mappings. We place ourselves in the 
 setting of  II.\ref{sub:7.2}. Assuming the congruence subgroup sufficiently
 small, the picture was the following: 

\begin{center}
 \begin{picture}(150,90)(0,-5)
  \put(10,60){\makebox(0,0)[c]{$\til{S}$}}
  \put(25,60){\vector(1,0){20}}
  \put(25,45){\vector(1,-1){20}}
  \put(60,60){\makebox(0,0)[c]{$\mathcal{M}$}}
  \put(60,45){\vector(0,-1){20}}
  \put(60,10){\makebox(0,0)[c]{$S$}}
  \put(115,60){\makebox(0,0)[c]{$\mathcal{D}$}}
  \put(75,60){\vector(1,0){25}}
  \put(88,65){\makebox(0,0)[c]{${}^\mathcal{P}$}}
  \put(115,45){\vector(0,-1){20}}
  \put(75,10){\vector(1,0){20}}
  \put(115,10){\makebox(0,0)[c]{$\mathcal{D}/\Gamma$}}
 \end{picture}
\end{center}
where $S$ was a suitable connected `tube-like' domain
 in the Shimura variety $\Sh_{\C_p}$ ($\mathcal{M}^0$ was a deformation
 space for $p$-divisible groups quasi-isogenous to a fixed one, and
 $\mathcal{D}$ was the period domain).

 The period mapping turned out to be expressable in terms of quotients of
 solutions of the Gauss-Manin connection (on $S$), 
 and the Gauss-Manin connection was trivializable on $\mathcal{M}^0 $,
 hence on $\tilde S$ (II.\ref{cor:7.2.6}).

 \medskip Let us now abandon the assumption that the congruence subgroup
 is sufficiently small. Then a similar picture holds 
 true: $\Gamma$ is now only virtually torsion-free, and the covering
 $\mathcal{M}^0\to S $ is now ramified. By removing the 
 singularities, one obtains a {\it temperate \'etale covering}, which
 trivializes the Gauss-Manin 
 connection (using II.\ref{cor:7.2.6} after passing to a suitable
 finite \'etale covering). It follows that the 
 {\it Gauss-Manin connection (restricted to $S$ deprived from the
 singularities) is in the image of the temperate 
 Riemann-Hilbert functor}. 

 The situation is therefore similar the the complex one, except that $S$
 is only a small domain in the Shimura variety in 
 general. The most favorable situation occurs in the case of global
 uniformization (II.\ref{sub:7.3}), \ie when $S$ is a component of $\Sh$.
\end{para}

\begin{para}\label{para::3.5.9} Let us now give a criterion for a connection
 in the image of $RH^{\mathrm{et}}$ to be in the image of $RH^{\mathrm{temp}}$.
\end{para}

\begin{pro}\label{pro::3.5.10} Let $\rho:
\pi_1^{\mathrm{et}}(S,\ovl{s}) \to GL(V)$ be a finite-dimensional
 discrete representation (corresponding to an \'etale local system 
 $\mathcal{V}$). Let us consider the following conditions:
 \begin{enumerate}
  \renewcommand{\theenumi}{\alph{enumi}}
  \item $\rho$ factors through a discrete representation of
	$\pi_1^{\mathrm{temp}}(S,\ovl{s})$,
  \item $\Im \rho$ is virtually torsion-free,
  \item $\Im \rho$ is finitely generated.
 \end{enumerate}
 Then {\rm (c)} $\Rightarrow$ {\rm (b)} $\Rightarrow$ {\rm (a)}. If moreover
 $S$ is the 
 analytification of a smooth algebraic $K$-variety, 
 then {\rm (a)} $\Rightarrow$ {\rm (c)}.
\end{pro}

\begin{proof}
 (c) $\Rightarrow$ (b): this is Selberg's lemma: any finitely
 generated subgroup of $GL(V)$ is virtually torsion-free.

 (b) $\Rightarrow$ (a) follows from \ref{thm::2.1.8}.~(b). 

 (a) $\Rightarrow$ (c): under (a), $\Coim\,\rho$ is the Galois group of a
 Galois temperate covering of $S$. If $S$ is 
 algebraic, $\Coim\,\rho$ must be finitely generated by \ref{pro::2.1.6}. 
\end{proof}

\begin{para}\label{para::3.5.11} We finish the section by showing that the
 connections which lie in the image of the temperate 
Riemann-Hilbert functor are algebrizable (this is not automatic in the
 non-archimedean situation), and even fuchsian: 
\end{para}

\begin{thm}\label{thm::3.5.12} Assume that $S$ is the analytification of
 a smooth algebraic $K$-variety $S^{\mathrm{alg}}$. Let 
 $(M,\nabla)$ be a vector bundle with integrable connection on $S$ which
 comes from a discrete representation of 
 $\pi_1^{\mathrm{temp}}(S,\ovl{s})$. Then $(M,\nabla)$ is the
 analytification of a (unique) algebraic vector 
 bundle with integrable connection on $S^{\mathrm{alg}}$, which has
 regular singularities and exponents in $\Q/\Z$ at any 
divisorial valuation of $K(S)$.  
\end{thm}

\begin{proof}
 By assumption, there is a finite Galois \'etale covering $S'\to
 S$ (with group $G$) such that $(M,\nabla)$ becomes 
 trivial on the topological universal covering of $S'$. The pull-back
 $(M,\nabla)_{S'}$ corresponds to a (unique) 
 representation of 
 $\pi_1^{\mathrm{top}}(S', s')$ (for any point $s'$ above $s$). Let
 $\ovl{S}'$ be a smooth algebraic compactification of $S'$ 
 (Hironaka). We may replace $K$ by the completion of its algebraic
 closure. By \ref{pro::1.1.3}, the homomorphism 
 $\pi_1^{\mathrm{top}}(S', s')\to\pi_1^{\mathrm{top}}(\ovl{S}', s')$ is
 an isomorphism, and it follows that $(M,\nabla)_{S'}$ extends 
 (uniquely) to a vector bundle with integrable connection on 
 $\ovl{S}'$. By GAGA (K\"opf), the latter is the analytification of a
 (unique) algebraic vector 
 bundle with integrable connection on $\ovl{S}^{\prime \mathrm{alg}}$. It is
 now easy to conclude by Galois descent from $S^{\prime \mathrm{alg}}$ to 
 $S^{\mathrm{alg}}$ (we refer to 
 \cite{andre01:_de_rham} for the details on regular singularities,
 exponents, and GAGA functor for connections).
\end{proof}

\begin{para}\label{para::3.5.13} Besides discrete representations of
 $\pi_1^{\mathrm{temp}}(S,\ovl{s})$, it is also useful to 
 consider {\it projective discrete representations}, \ie discrete
 representations with value in $PGL(V)$.  They give rise to {\it
 projective connections} over $S^{\mathrm{alg}}$, at least when
 $S^{\mathrm{alg}}$ is an {\it affine curve}. 

 Indeed, replacing $S^{\mathrm{alg}}$ by a finite Galois \'etale
 covering $S^{\prime \mathrm{alg}}$ with Galois group $G$, we may
 replace $\pi_1^{\mathrm{temp}}(S')$ by $\pi_1^{\mathrm{top}}(S')$, 
 which is a free group, for which the lifting property is immediate. This
 provides a connection $\nabla'$ over $S^{\prime \mathrm{alg}}$,
 and the associated projective connection is endowed with a $G$-action
 compatible with the $G$-action on $S^{\prime \mathrm{alg}}$, hence
 descends to a projective connection on $S^{\mathrm{alg}}$.
\end{para}

One has more: 
\begin{pro}\label{pro::3.5.14} Then any projective discrete representation
 $\ovl{\rho}:$ $\pi_1^{\mathrm{temp}}(S,\ovl{s})$ $\to$ $PSL(V)$ lifts to
 a genuine discrete representation
 $\rho:\,\pi_1^{\mathrm{temp}}(S,\ovl{s})\to SL(V)$, after replacing $K$
 by a finite extension.
\end{pro}

\begin{proof}
 We have seen that $\ovl{\rho}$ corresponds to a projective connection on
 $S^{\mathrm{alg}}$. Now it is well known that any projective connection
 lifts to a genuine connection after replacing $K$ by a finite extension
 (because $\mathrm{H}^2(S^{\mathrm{alg}}_{\ovl{K}, et},
 \Z/2\Z)=0$, \cf \cite[2.2.2.1]{katz90:_expon}). It suffices to check that the
 pull-back $\nabla''$ of this connection on $S^{\prime \mathrm{alg}}$
 comes from a discrete representation
 $\rho:\,\pi_1^{\mathrm{temp}}(S',\ovl{s}')\to SL(V)$. But we have on the
 other hand our connection $\nabla'$ which is projectively equivalent to
 $\nabla''$, 
 and which we may assume to come from a discrete representation 
 $\pi_1^{\mathrm{temp}}(S',\ovl{s}')\to SL(V)$. Let $S''$ be a finite
 \'etale covering of $S'$ such that the pull-back of any class in
 $\mathrm{H}^1(S^{\mathrm{alg}}_{\ovl{K}, 
 et}, \Z/2\Z)$ vanishes in $\mathrm{H}^1(S^{\prime \prime\,
 \mathrm{alg}}_{\ovl{K}, et}, 
 \Z/2\Z)$. Since $\nabla'$ and $\nabla''$ are projectively equivalent,
 their pull-backs on $S''$ are isomorphic 
 ({\it loc.~cit.}). It follows that $\nabla''$ comes from a discrete
 representation $\pi_1^{\mathrm{temp}}(S',\ovl{s}')\to SL(V)$.        
\end{proof}

\medskip We leave it as an exercise to translate \ref{pro::3.5.10} in
the context of projective connections.

 
\newpage
\section{Non-archimedean orbifolds, uniformizing differential equations
and period mappings.}\label{sec::4}

\markboth{\thechapter. $p$-ADIC ORBIFOLDS AND MONODROMY.}%
{\thesection. NON-ARCHIMEDEAN ORBIFOLDS.}

\begin{abst}
We review the complex case, and then define and study non-archimedean
orbifolds, their associated orbifold fundamental groups, and
uniformizing differential equations.
\end{abst}


\subsection{Review of complex orbifolds.}\label{sub::4.1}

\begin{para}\label{para::4.1.1} There are several variants of the notion 
 of complex orbifold. We follow \cite{calip} and see orbifolds 
 as being locally quotients of a complex manifold by a finite group 
 acting faithfully. More precisely, a {\it complex orbifold} 
 $\mathcal{S}=(\ovl{S}, (Z_i,e_i))$ consists in 
 \begin{itemize}
  \renewcommand{\labelitemi}{\normalfont\bfseries\textendash}
  \item a normal analytic space $\ovl{S}$,
  \item a locally finite collection of irreducible (Weil) divisors
	$(Z_i)$,
  \item a positive integer\footnote{in other variants, one also allows
	$e_i=\infty$. This leads to `parabolic structures' which have
	indeed non-archimedean counterparts in positive characteristic
	(cusps of Drinfeld modular curves...). But we shall be mostly
	interested in non-archimedean orbifolds in characteristic zero,
	where such `parabolic structures' cannot occur. } $e_i$
	attached to each $Z_i$.
 \end{itemize}


\medskip\noindent It is assumed that $\ovl{S}$ is covered by so-called
 {\it orbifold charts}, \ie morphisms  $W 
\xrightarrow{\phi} V \subset \ovl{S} $ where 
 \begin{itemize}
  \renewcommand{\labelitemi}{\normalfont\bfseries\textendash}
  \item $W$ is a complex manifold,
  \item $\phi$ is a (ramified) covering in the sense of \ref{def::1.2.3}
	(\ie $V$ is covered by open subsets $U$ such that $\coprod V_j =
	\phi^{-1}U \to U$ and the restriction of $\phi$ to every $V_j$
	is finite),
  \item except for a closed subset $\Sigma$ of codimension $2$ in $V$,
	$\phi$ is \'etale above $V\setminus (V\cap\bigcup Z_i)$ (hence
	induces a topological covering of $V\setminus (V\cap\bigcup
	Z_i)$), and is {\it ramified with index $e_i$ above
	$Z_i$}.\footnote{it is easily seen that this definition is equivalent
	to that of \cite{calip}, using {\it loc.~cit.} 14.2.}
 \end{itemize}

\medskip Note that if $s$ is a non-singular point of $\ovl{S}$ not
 contained in $\bigcup Z_i$, then any smooth open neighborhood of $s$ in
 $\ovl{S}$ disjoint from $\bigcup Z_i$ provides an orbifold chart around
 $s$. Hence we may and shall assume that $\Sigma $ is contained in
 $\bigcup Z_i\cup \ovl{S}^{\mathrm{sing}}$ (where $\ovl{S}^{\mathrm{sing}}$ denotes the
 set of singular points of $\ovl{S}$).
 
 \noindent We shall often write
 $Z=\bigcup Z_i\cup
 \ovl{S}^{\mathrm{sing}}$, and
 $S=\ovl{S}\setminus Z$ (which is a manifold).  

\medskip\noindent {\it Remark.} By shrinking the $W$'s, one can assume
 that they are open subsets of $\C^n$; this matches with the conventional
 use of the word `chart'. However, we do not impose here this condition
 which would have no non-archimedean counterpart.


\end{para}

\begin{para}\label{para::4.1.2} A {\it morphism of orbifolds}
 $f:\mathcal{S}'\to \mathcal{S}$ is a morphism $\ovl{S}'\to \ovl{S}$, such
 that $\ovl{S}'$ is covered by orbifold charts  $\phi': W'\to V'$ with
 the property that $f\circ \phi'$ factors through some orbifold 
 chart $\phi:W\to V$ of $\mathcal{S}$: $f\circ \phi'=\phi\circ f'$ 
(for some morphism $f'\colon W'\to W$).

 \medskip\noindent {\it Example.}  Let $(\ovl{S}, (Z_i,e_i)_{i\in I})$ be
 an orbifold, and let $J\subset I$ be the subset of indices for which
 $e_i>1$. Then the identity of $\ovl{S}$ induces an isomorphism
 $(\ovl{S}, (Z_i,e_i)_{i\in I})\cong (\ovl{S}, (Z_i,e_i)_{i\in J})$. In
 other words, one can `forget' the divisors with index $e_i=1$.
 
 \noindent Considering complex manifolds as orbifolds with an empty
 collection of divisors ($S=\ovl{S}$), one can embed the 
 category of complex manifolds as a full subcategory of the category of
 complex orbifolds.

 \noindent It will be convenient to attach to any irreducible divisor
 $D\subset \ovl{S}$ an index $e_D\in \Z_{>0}$ which is $e_i$ if
 $D=Z_i$ and $1$ if $D$ is none of the $Z_i$'s.


\end{para}

\begin{para}\label{para::4.1.3} A morphism $f$ is called an {\it
 orbifold-covering} if it is a covering in the sense of \ref{def::1.2.3}
 and for any 
 orbifold chart $\phi': W'\to V'$ of $\mathcal{S}'$, $f\circ\phi': W'\to
 V=f(V')$ is an orbifold chart of $\mathcal{S}\;$ (in
 loose words: any lifting $f'$ as in \ref{para::4.1.2} is 
 \'etale).  
 
 \medskip\noindent The restriction of $f$ above $S$ is then an \'etale
 covering: indeed, this is clear except above
 a closed subset of codimension $\geq 2$ in $S$, thus above the whole of
 $S$ by purity since $S$ is smooth.
 
 \medskip\noindent {\it Example.} $(\A^1, (0,e'))
 \xrightarrow{z\mapsto z^n} (\A^1, (0,e)) $ is
 a morphism if and only if $e| ne'$, and is an orbifold-covering if
 and only if 
 $e=ne'$.   
  
\medskip \noindent {\it Remark.} Let $f:\mathcal{S}'\to \mathcal{S}$ be
 an orbifold-covering. As we just saw, we may assume 
 that all the indices $e'_i,e_j$ are $>1$. Then $f(Z')\subset Z$ and
 $f^{-1}(S)\subset S'$.
\end{para}

\begin{para}\label{para::4.1.4} Following Thurston, an orbifold is said to be
 good or {\it uniformizable} if it admits a {\it global} orbifold chart
 (\ie with $V=\ovl{S}$).  In that case, there exists a {\it universal} 
 global orbifold chart, called the {\it universal orbifold-covering} (it
 is indeed an orbifold-covering in the above sense); its underlying space
 is a simply-connected complex manifold if $\ovl{S}$ is connected.
 
 \medskip\noindent {\it Example.} By a classical result due to Fox,
 $(\P^1, (\zeta_i,e_i)_{i=0,\ldots, n})$ is uniformizable if and only if
 $n\geq 2$ or ($n=1$ and $e_0=e_1$). Under these conditions, 
 there exist algebraic global charts which are finite coverings of $\P^1$. 
\end{para}

\begin{para}\label{para::4.1.5} It is proved in {\it loc.~cit.} that orbifold
 charts are locally unique. Apart from a subset of codimension $2$,
 they are given in local coordinates around $\phi^{-1}(Z_i)$ by:
 $(z_1,z_2,\ldots, z_n)\mapsto (z_1^{e_i},z_2,\ldots, z_n)$.
 
 \noindent It follows that there exist orbifold charts of the form
 $W\to V=W/G$ ({\it loc.~cit.} 14.6), where $W$ is an
 open subset of $\C^n$ stable under a finite group $G$ of linear
 transformations.
\end{para}

\begin{para}\label{para::4.1.6} Let $\mathcal{S}=(\ovl{S}, (Z_i,e_i))$ be a
 connected orbifold, and let $\Gamma$ be a discrete group acting
 faithfully and properly by automorphisms on $\mathcal{S}$. The {\it orbifold
 quotient} $\mathcal{S}/\Gamma$ is defined as follows: the underlying
 analytic space is $\ovl{S}/\Gamma$. The ramifications divisors are the
 images of the $Z_i$'s, and the irreducible divisors $D$ with
 non-trivial fixer $\Gamma_D$; the
 latter are endowed with the multiplicity $e_D.| \Gamma_D|$.

 Note that the projection $\mathcal{S}\to \mathcal{S}/\Gamma$ is an
 orbifold-covering.
\end{para}

\begin{para}\label{para::4.1.7} One of the main applications of complex
 orbifolds is to represent moduli problems for isomorphism classes of
 objects with finite automorphism group. We have already alluded to this
 in II.\ref{sub:1.1}.  Here we shall content ourselves with examining
 the case of complex elliptic curves.

 \noindent Let $H$ be the subgroup of $SL_3(\R)$ consisting
 of matrices of the form
 $\begin{pmatrix}
 1&x_1&x_2\\ 0&a&b\\ 0&c&d
  \end{pmatrix}$ (this is a
 semi-direct product $SL_2(\R)\ltimes\R^2$).  It acts on the product
 $\mathfrak{h}$ 
 $\times\C$ of the Poincar\'e upper half plane $\mathfrak{h}$ and $\C$ by
 the rule 
 \begin{equation*}
  \begin{pmatrix}
   1&x_1&x_2\\0&a&b\\0&c&d
  \end{pmatrix}.
  (\tau, v)=\Bigl(\frac{a\tau+b}{c\tau +d},\frac{v+x_1\tau+x_2}{c\tau+d}\Bigr).
 \end{equation*}
 As is well-known, elliptic curves $E$ together with a symplectic isomorphism
 $\mathrm{H}_1(E,\Z)\cong \Z^2$ are classified by points of the
 Poincar\'e upper half plane $\mathfrak{h}$, the universal object being
 the quotient of $\mathfrak{h}\times\C$ by the action of the subgroup 
 $1.\Z^2\subset H$ (mapping to $\mathfrak{h}$ through the first projection). 

 Let $\Ker(SL_2(\Z)\to SL_2(\Z/N\Z))$ be the full congruence group of
 level $N$, $\Gamma(N)$ its image in $PGL(2)$, and $\tilde\Gamma(N)$ be
 the subgroup 
 \begin{equation*}
  \Ker(SL_2(\Z)\to SL_2(\Z/N\Z)).\Z^2\subset H.
 \end{equation*}
 For $N\geq 3$, the group $\tilde \Gamma(N)$ is
 torsion-free and acts freely on $\mathfrak{h}\times\C$; elliptic curves 
 together with a level $N$ structure (\cf II.\ref{para:1.1.2}) are
 classified by $\mathfrak{h}/\Gamma(N)$, the universal object being
 $(\mathfrak{h}\times\C)/\tilde \Gamma(N)\to \mathfrak{h}/\Gamma(N)$. 
 
 This still holds for $N=1$ or $2$, provided one interprets these
 quotients as orbifold quotients. For $N=2$, the element 
  $\begin{pmatrix}1&&0\\ &-1&\\ 0&&-1 \end{pmatrix} \in
 \tilde \Gamma(N)$ of order $2$ has fixed divisors in
 $\mathfrak{h}\times\C$: the quotient of
 $\{(\tau,x_1\tau+x_2) \mid x_1,x_2\in (\frac{1}{2}\Z)^2\}$ by $\tilde
 \Gamma(N)$, which is a disjoint union $\coprod_1^4 Z_i$ of four
 `horizontal' curves (horizontal with respect to the projection to 
 $\mathfrak{h}/\Gamma(2)=\P^1\setminus\{0,1,\infty\}$); they have the
 index $e_i=2$ (the order of the fixer).
 
This orbifold is uniformizable: the Legendre elliptic pencil provides a
global orbifold chart. Note however that, as an analytic variety,
$\mathfrak{h}\times\C/\tilde \Gamma(2)$ is a not an elliptic surface but
a $\P^1$-bundle.

The case $N=1$ is similar: $\mathfrak{h}/\Gamma(1)=\A^1$ is the $j$-line. It should be considered as an orbifold, with
divisor $\{j=0 \,{\rm(index \;3)} ,\;j=1728\,{\rm(index \;2)} \}$. In addition to the image of the horizontal divisors
$Z_i$ on $\mathfrak{h}\times\C/\tilde \Gamma(1)$, there are three vertical divisors with index $2$, and two vertical
divisors with index $3$.  


\end{para}

\begin{para}\label{para::4.1.8} More generally, let us consider a Shimura
 variety of PEL type as in II.\ref{para:1.2.4}.  Recall from {\it
 loc.~cit.} the commutative diagram
\begin{equation*}
\begin{CD}
 \tilde S @>\mathcal{P}>> \mathcal{D}\\
 @V\mathcal{Q}VV @VV\mathcal{Q}V\\
 S @>\mathcal{P}>> \mathcal{D}/\Gamma
\end{CD}
\end{equation*}
where $S$ is a component of the Shimura variety, the horizontal maps are
 {\it isomorphisms}, and $\Gamma$ is a congruence subgroup. If $\Gamma$
 has torsion, it is appropriate to consider $S$ as an orbifold quotient 
 (defined by $\mathcal{Q}$ as in \ref{para::4.1.6}). 
\end{para}


\subsection{Orbifold fundamental groups.}\label{sub::4.2}

\begin{para}\label{para::4.2.1} Another important application of orbifolds is
 the Galois theory of {\it ramified} coverings.   

  Let $\mathcal{S}=(\ovl{S}, (Z_i,e_i))$ be a uniformizable connected
orbifold. We denote by $\tilde{\mathcal{S}}$ the universal
orbifold-covering (which is a simply-connected complex manifold,
together with a (ramified) covering morphism $\tilde{\mathcal{S}}\to
\ovl{S}$). We fix a base-point $s$ of $S =\ovl{S}\setminus \,(\bigcup
Z_i\cup\ovl{S}^{\mathrm{sing}})$, and a point $\tilde s$ of $\tilde{\mathcal{S}}$
above $s$. Then the pair $(\tilde{\mathcal{S}},\tilde s)$ is unique up to
unique isomorphism.
\end{para}

\begin{para}\label{para::4.2.2} Each divisor $Z_i$ defines a conjugacy class
of elements $\gamma_i$ in $\pi_1^{\mathrm{top}}(S,s)$: the local
monodromy at $Z_i$ based at $s$.

\medskip The {\it orbifold fundamental group}
$\pi_1^{\mathrm{orb}}(\mathcal{S},s)$ is the quotient of
$\pi_1^{\mathrm{top}}(S,s)$ by the {\it normal subgroup}
$\an{\gamma_i^{e_i}}$ generated by the $\gamma_i^{e_i}$'s.

 \medskip This definition is justified by the following fact
 (\cf \cite[14.10]{calip}): 
 one has a commutative square of functors 
 \begin{equation*}
  \begin{CD}
   \{\text{orbifold-coverings of $\mathcal{S}$}\} @>\subset>>
   \{\text{\'etale (= topological) coverings of $S$}\}\\
   @V\wr VV  @VV\wr V \\
   \{\text{$\pi_1^{\mathrm{orb}}(\mathcal{S},s)$-sets}\}
   @>\subset >>
   \{\text{$\pi_1^{\mathrm{top}}(S,s)$-sets}\}
  \end{CD}
 \end{equation*}
 where the top horizontal functor is the (fully faithful) restriction
 functor, and the vertical 
 functors are equivalences of categories. 

 The orbifold fundamental group may also be identified with the group of
 automorphisms of the (ramified) covering 
 $\tilde{\mathcal{S}}\to \ovl{S}$. Galois orbifold-coverings correspond to
 normal subgroups of $ \pi_1^{\mathrm{orb}}$.

 \medskip\noindent {\it Example.} Let $\mathcal{S}=(\P^1, (\zeta_i,
 e_i)_{i=0,\ldots, n})$ with $n\geq 2$ or ($n=1$ and $e_0=e_1$). Then
 $\pi_1^{\mathrm{orb}}(\mathcal{S},s)$ has generators $\gamma_0,\ldots, \gamma_n$
 and relations $\gamma_0^{e_0}=\ldots =\gamma_n^{e_n}= \Pi \gamma_i = 1$.
 In particular, it depends only, up to isomorphism, on the collection of
 integers $e_i$, and not on the $\zeta_i's$.

 \noindent The universal orbifold-covering is $\P^1, \A^1$ or the open
 unit disk $\matheur{D}$ according to whether $\sum (1-\frac{1}{e_i})$ is $< 2,\;=2,\;$
 or $>2$ respectively. 
 
\medskip\noindent {\it Remark.} The {\it finite} orbifolds coverings are
 classified by a profinite group
 $\pi_1^{\mathrm{alg.orb}}(\mathcal{S},s)$, which 
 is actually the profinite completion of
 $\pi_1^{\mathrm{orb}}(\mathcal{S},s)$.  When $\ovl{S}$ is 
 algebraic, this profinite group has a purely algebraic definition. 


\end{para}

\begin{para}\label{para::4.2.3} Let $\phi:\,\ovl{S}'\to \ovl{S}$ be a global orbifold chart, such that the induced map $S'\to S$ is a
Galois (topological) covering with group $G$. Let $s'$ be a point of $S'$ above $s$. Then one has an exact sequence
\begin{equation*}
1\to \pi_1^{\mathrm{top}}(\ovl{S}',s')\to \pi_1^{\mathrm{orb}}(\mathcal{S},s) \to G \to 1.
\end{equation*}



\end{para}

\subsection{Uniformizing differential equations.}\label{sub::4.3}

\begin{para}\label{para::4.3.1} We follow the exposition in \cite[4,
 5.2]{fde} and \cite[4]{aeoaafg}. We now assume that $\ovl{S}$ is a 
complex smooth projective {\it curve}. Then $Z=\{\zeta_0,\ldots
 ,\zeta_n\}$ is just a finite set of points and $S=\ovl{S}\setminus Z$.

 We assume that $\mathcal{S}$ is uniformizable. By Riemann's
 uniformization theorem, the universal orbifold-covering
 $\tilde{\mathcal{S}}$ is then isomorphic to either
 $\P^1, \A^1$ or the open unit disk $\matheur{D}$.  We fix such an
 isomorphism\footnote{in some cases, it will be more 
 convenient to work with $\mathfrak{h}$ instead of $\matheur{D}$.}; the
 group $\Aut 
 \tilde{\mathcal{S}}\;$ is then a Lie subgroup of $PSL_2(\C)$. The
 orbifold fundamental group $\pi_1^{\mathrm{orb}}(\mathcal{S},s)$ is a discrete
 subgroup of $\Aut\, \tilde{\mathcal{S}}$ acting properly on
 $\tilde{\mathcal{S}}$, and the quotient space is
 $\tilde{\mathcal{S}}/\pi_1^{\mathrm{orb}}(\mathcal{S},s)=\ovl{S}$.      
\end{para}

\begin{para}\label{para::4.3.2} We fix a non-constant rational function $z\in
 \C(\ovl{S})$. We shall assume for simplicity that $z$ has no 
 pole on $S$. 

 \noindent Having identified $\tilde{\mathcal{S}}$ with a domain of
 $\P^1$ allows us to consider the canonical map 
 \begin{equation*}
  \tilde S\to \tilde{\mathcal{S}}\subset \P^1
 \end{equation*}
 as an analytic\footnote{in the case $\tilde{\mathcal{S}}=\P^1$,
 it is only meromorphic a priori. But in that case $\ovl{S}=\P^1$ and
 the covering $\P^1\to \P^1$ is finite. It will be convenient to
 normalize things so that $\infty$ lies above $\infty$ and $z$ is the
 standard coordinate on $\P^1$. With this choice, $\tau$ has no pole.} 
 function $\tau$ on $\tilde S$ (in loose terms, this is the multivalued
 inverse of the projection
 $\tilde{\mathcal{S}}\to
 \ovl{S}=\tilde{\mathcal{S}}/\pi_1^{\mathrm{orb}}(\mathcal{S},s)$,
 restricted to $S$). Notice that the derivation 
 $\frac{d}{dz}$ extends canonically to $\mathcal{O}(\tilde S)$. 
\end{para}

\begin{pro}\label{pro::4.3.3} There is a unique differential equation of
 the form
 \begin{equation*}
  y''- qy=0,\;\;\; q\in \C(\ovl{S}),
 \end{equation*}
 the {\rm uniformizing differential equation}, such that $\tau$ is the
 quotient of two solutions in $\mathcal{O}(\tilde S)$ of this
 differential equation. It is fuchsian and has singularities only at
 $Z=\{\zeta_0,\ldots ,\zeta_n\}$.  The monodromy group pointed at $s$ is
 a discrete finitely generated subgroup of $SL_2(\C)$ whose image in
 $PSL_2(\C)$ is $\pi_1^{\mathrm{orb}}(\mathcal{S},s)$.

 \hfill\break If $\ovl{S}=\P^1$ and $\zeta_n=\infty$, the exponents are
 rational numbers given by the Riemann scheme:
 \begin{equation*}
  \begin{pmatrix}
   \zeta_0 & \ldots & \zeta_{n-1} & \zeta_n=\infty \\
   &&&\\
   \frac{1}{2}(1+\frac{1}{e_0})& \ldots &\frac{1}{2}(1+\frac{1}{e_{n-1}})&
   \frac{-1}{2}(1+\frac{1}{e_{n}}) \\
   \\
   \frac{1}{2}(1-\frac{1}{e_0})& \ldots &\frac{1}{2}(1-\frac{1}{e_{n-1}})&
   \frac{-1}{2}(1-\frac{1}{e_{n}}) \\
  \end{pmatrix}.
 \end{equation*}
\cf {\it loc.~cit.}
\end{pro}

\begin{proof}
 Existence: $q$ is nothing but the Schwarz derivative
 \begin{equation*}
  \{\tau, z\}=
   \Bigl(\frac{d\tau}{dz}\Bigr)^{1/2}\frac{d^2}{dz^2}
   \Bigl[\Bigl(\frac{d\tau}{dz}\Bigr)^{-1/2}\Bigr],
 \end{equation*} 
 which is single-valued on $\ovl{S}$. A basis of solutions of the
 uniformizing differential equation is given by
 $(\frac{d\tau}{dz})^{-1/2}$ and $\tau.(\frac{d\tau}{dz})^{-1/2}$. The
 fact that $q\in \C(\ovl{S})$ follows immediately from the invariance of
 the Schwarz derivative under homographic transformations.

 Unicity: this is a purely algebraic fact: $q$ is determined by
 the projective connection attached to the differential equation $y''- qy=0$. 

 The determination of the Riemann scheme comes from this
 expression and the fact that around $\zeta_{i<n}$ (\resp around
 $\infty$) $\tau$ is given by $z-\zeta_i\mapsto c_i.(z-\zeta_i)^{1/e_i} +
 h.o.t.$ (\resp $1/z\mapsto c_n/z^{e_n} + h.o.t.$).
\end{proof}

\begin{rems}
  (i) The case of three singularities
 will be detailed in the next section. In the case of more than three
 singularities, it is extremely difficult in general to write down
 explicitly the uniformizing differential equations.

 (ii) In \cite{fde}, an extension of the theory of
 uniformizing differential equations to higher dimensions is
 developed. Its practical scope seems limited however to the case when
 $\tilde{\mathcal{S}}$ is a symmetric domain.

 (iii) The projective equivalence class of the uniformizing
 differential equation is a variant of what is called the Schwarz
 structure or canonical indigenous $\P^1$-bundle over $\ovl{S}$, \cf 
 \cite{gunning67:_special_rieman}, \cite{mochizuki96}.
\end{rems}

 
\subsection{Non-archimedean orbifolds.}\label{sub::4.4}

\begin{para}\label{para::4.4.1} We now work over a complete non-archimedean
 field $K$. The definition of orbifolds is parallel to the one 
 which we have selected in the complex situation. Being mainly interested
 in orbifolds which are quotients of 
 manifolds by finite groups, it is natural to replace everywhere \'etale
 (= topological) coverings in the complex case by {\it temperate}
 \'etale coverings in the non-archimedean context.

 \medskip\noindent A {\it $K$-orbifold} $\;\mathcal{S}=(\ovl{S},
 (Z_i,e_i))$ consists in 
 \begin{itemize}
  \renewcommand{\labelitemi}{\normalfont\bfseries\textendash}
  \item a normal $K$-analytic space $\ovl{S}$ (in the sense of
	Berkovich),
  \item a locally finite collection of irreducible divisors $(Z_i)$,
  \item a positive integer $e_i$ attached to each $Z_i$.
 \end{itemize}

 \noindent It is assumed that $\ovl{S}$ is covered by so-called
 {\it orbifold charts}, \ie morphisms
 $W \xrightarrow{\phi} V \subset \ovl{S} $ where 

\begin{itemize}
 \renewcommand{\labelitemi}{\normalfont\bfseries\textendash}
 \item $W$ is a $K$-manifold,
 \item $\phi$ is a (ramified) covering in the sense of \ref{def::1.2.3}
       (\ie $V$ is covered by open subsets $U$ such that 
       $\coprod V_j = \phi^{-1}U \to U$ and the restriction of $\phi$ to
       every $V_j$ is finite),
 \item except for a closed subset $\Sigma$ of codimension $2$ in $V$,
       $\phi$ restricts to a {\it temperate} \'etale covering 
       above $V\setminus (V\cap\bigcup Z_i)$, and is {\it ramified with index
       $e_i$ above $Z_i$}.
\end{itemize}  

 \medskip Here again, we may and shall assume that that $\Sigma $ is
 contained in $\bigcup Z_i\cup \ovl{S}^{\mathrm{sing}}$ and set
 $Z=\bigcup Z_i\cup \ovl{S}^{\mathrm{sing}},\;S=\ovl{S}\setminus Z$
 (which is a $K$-manifold).
\end{para}

\begin{rems}
  (i) At this point, because the notion of orbifold is local, it would
 not change if we had modified the definition of orbifold charts by
 replacing `temperate \'etale' by `finite \'etale' or by `\'etale'.  In
 contrast, we insist that orbifold charts are defined in terms of
 temperate coverings.
 
 (ii) Since we only prescribe the ramification indices and not
 the ramification groups, this notion of orbifold is 
 only reasonable in the tame case, \ie when $\ch(K) \nmid e_i$.  We 
 shall be mostly interested in the characteristic $0$ case.

 (iii) If $\dim \ovl{S} = 1$, then $S=\ovl{S}\setminus \bigcup Z_i$.
\end{rems}

\begin{para}\label{para::4.4.2} A {\it morphism of orbifolds}
 $f:\mathcal{S}'\to \mathcal{S}$ is a morphism $\ovl{S}'\to \ovl{S}$, such
 that $\ovl{S}'$ is covered by orbifold charts  $\phi': W'\to V'$ with
 the property that $f\circ \phi'$ factors through some orbifold chart
 $\phi:W\to V$ of $\mathcal{S}$: $f\circ \phi'=\phi\circ f'$.

 \medskip\noindent {\it Example.} Let $(\ovl{S}, (Z_i,e_i)_{i\in I})$ be
 an orbifold, and let $J\subset I$ be the subset of indices for which
 $e_i>1$. Then the identity of $\ovl{S}$ induces an isomorphism
 $(\ovl{S}, (Z_i,e_i)_{i\in I})\cong (\ovl{S}, (Z_i,e_i)_{i\in J})$. In
 other words, one can `forget' the divisors with index $e_i=1$.

 \noindent Considering $K$-manifolds as orbifolds with an empty
 collection of divisors ($S=\ovl{S}$), one can embed {\it the category of 
 $K$-manifolds as a full subcategory of the category of $K$-orbifolds}.

\noindent It will be convenient to attach to any irreducible divisor
 $D\subset \ovl{S}$ an index $e_D\in \Z_{>0}$ which is $e_i$ if
 $D=Z_i$ and $1$ if $D$ is none of the $Z_i$'s.
\end{para}

\begin{para}\label{para::4.4.3} A morphism $f$ is called an
{\it orbifold-covering} if it is a covering in the sense of
 \ref{def::1.2.3} and 
 for any orbifold chart $\phi': W'\to V'$ of $\mathcal{S}'$, $f\circ
 \phi': W'\to V=f(V')$ is an orbifold chart of $\mathcal{S}\;$ (note that if
 $\mathcal{S}'$ and $\mathcal{S}$ are just $K$-manifolds, an orbifold-covering
 amounts to a temperate \'etale covering, by our definition of orbifold
 charts).
 
 \noindent Orbifold-coverings of $\mathcal{S}$ form a category in an obvious
 way, denoted by $\Cov_\mathcal{S}$.
 
\medskip\noindent The restriction of $f$ above $S$ is then an \'etale
 covering: indeed, this is clear except above a closed subset of
 codimension $\geq 2$ in $S$, thus above the whole of $S$ by purity since
 $S$ is smooth.
 However, we do not know in general whether it
 is temperate (due to the non-local nature of the definition of
 temperate coverings).
 
\medskip\noindent {\it Example.} Assume that $\ch(K) \nmid n$. Then
 $(\A^1, (0, e')) \xrightarrow{z\mapsto z^n} (\A^1,
 (0, e)) $ is a morphism of orbifolds if and only if $e| ne'$, and is
 an orbifold-covering if and only if $e=ne'$.   

 \medskip \noindent {\it Remark.} Let $f:\mathcal{S}'\to \mathcal{S}$ be an
 orbifold-covering. As we just saw, we may assume that all the indices
 $e'_i,e_j$ are $>1$. Under this assumption, $f(Z')\subset Z$ and
$f^{-1}(S)\subset S'$.
\end{para}

\begin{para}\label{para::4.4.4} An orbifold is said to be {\it
 uniformizable}\footnote{we avoid the terminology `good' in 
 order to prevent confusion with good spaces in the sense of Berkovich,
 \cf last section.} if it admits a {\it global} orbifold chart.
 
 However, {\it universal global orbifold charts do not exist}; this is
 no surprise: there is no universal temperate covering of a
 $K$-manifold, except in trivial cases. For the same reason, {\it
 orbifold charts are not locally unique}.
 
 \medskip Let $f:\mathcal{S}'\to \mathcal{S}$ be an orbifold-covering of a
 uniformizable orbifold $\mathcal{S}$. Then the {\it restriction of $f$ to
 $f^{-1}(S)$ is a temperate \'etale covering} of $S$.

 Indeed, let $W\to \ovl{S}$ be a global orbifold chart. Then
 $W_{|S}\times_S f^{-1}(S)\to f'(W'_{|S}\times_S f^{-1}(S))= f^{-1}(S)$
 is an orbifold chart for $\mathcal{S}'$, hence $W_{| S}\times_S
 f^{-1}(S)\to S$ is an orbifold chart for $\mathcal{S}$. Since its image is
 $S$ itself, it is a temperate covering of $S$, of which $f^{-1}(S)\to S$ is 
 a quotient.

\medskip\noindent {\it Example.} Assume that $K$ is algebraically closed
of characteristic zero. Then $(\P^1, (\zeta_i, e_i)_{i=0,\ldots, n})$
is uniformizable if and only if $n\geq 2$ or ($n=1$ and $e_0=e_1$),
just as in the complex case. This actually follows from the complex
 case: one may assume that $K$ has the same cardinality as $\C$, choose
 an abstract isomorphism $K\cong \C$, and use the fact that in the
 complex case, one can find a global algebraic orbifold chart which is a
 finite ramified covering of $\ovl{S}$.
\end{para}

\begin{para}\label{para::4.4.5} Let us again assume that $K$ is of
 characteristic zero. Let $\mathcal{S}=(\ovl{S}, (Z_i,e_i))$ be a
 connected $K$-orbifold, and let $\Gamma$ be a virtually free discrete group acting
 faithfully and properly by automorphisms on $\mathcal{S}$.  The {\it
 orbifold quotient} $\mathcal{S}/\Gamma$ is defined as in the complex
 situation: the underlying analytic space is $\ovl{S}/\Gamma$. The
 ramification divisors are the images of the $Z_i$'s, and the
 irreducible divisors $D$ with non-trivial fixer $\Gamma_D$; the latter
 are endowed with the multiplicity $e_D.| \Gamma_D|$.
 Note that the projection $\mathcal{S}\to \mathcal{S}/\Gamma$ is an
 orbifold-covering.

 Interesting examples with $\ovl{S}$ open in $\P^1$ are discussed in
 \cite{van92:_discr_mumfor}. 
\end{para}

\begin{para}\label{para::4.4.6} One of the applications of $K$-orbifolds would
 be to represent moduli problems for isomorphism classes of objects with
 finite automorphism groups.

 Typical instances occur in the context of $p$-adic uniformization of
 Shimura varieties. Let us briefly mention how.  Recall from
 II.\ref{sub:7.2}, \ref{sub:7.3}
 the following basic commutative diagram

 \begin{equation*}
  \begin{CD}
   \mathcal{M}^0 @>\mathcal{P}>> \mathcal{D}\\
   @V\mathcal{Q}VV @VV\mathcal{Q}V\\
   S @>\mathcal{P}>> \mathcal{D}/\Gamma
  \end{CD}
 \end{equation*}
$S$ is a tube-like domain in the $p$-adic
 analytification of a Shimura variety, the horizontal maps are \'etale,
and $\Gamma$ is an arithmetic subgroup in the semi-simple group
$\ovl{G}^{\mathrm{ad}}$. Under the assumption that $\Gamma$ is small
enough, the maps $\mathcal{Q}$ are topological covering maps.

If one drops that assumption, one has to consider $S$ as an orbifold
quotient (defined by $\mathcal{Q}$ as in \ref{para::4.4.5}). This is
 especially relevant 
in the case of global uniformization, \ie when $S$ is a connected
component of the Shimura variety.

The fake projective planes of Mumford, Ishida, Kato, provide interesting
examples of this situation, \cf II.\ref{sub:7.5} and \cite{opuofpp}.
We shall see several other instances in the sequel.
\end{para}

\subsection{Non-archimedean orbifold fundamental groups.}\label{sub::4.5}

\medskip From now on, we assume that $K$ is {\it algebraically closed of
characteristic zero}. 

\begin{para}\label{para::4.5.1}  Let $\mathcal{S}=(\ovl{S}, (Z_i,e_i))$
 be a uniformizable connected orbifold.

As we have seen in \ref{para::4.4.4}, the restriction above $S$ of an
 orbifold-covering is a temperate \'etale covering. This gives 
rise to a restriction functor 
\begin{equation*}
   \Cov_\mathcal{S} \longrightarrow \Cov_S^{\mathrm{temp}}.
\end{equation*}
\end{para}

\begin{pro}\label{pro::4.5.2} This functor is fully faithful. Its image
 consists of those temperate \'etale coverings $S'\to S$ which extend to
 ramified coverings $\ovl{S}'\to \ovl{S}$, with $\ovl{S}'$ normal, such
 that the ramification index at each irreducible divisor above $Z_i$
 divides $e_i$ (for every $i$).

 \noindent This image is stable under
 taking connected components, fiber products (over $S$), quotients,
 finite disjoint unions, and Galois closure.
\end{pro}

\begin{proof}
 From Kiehl's theorem on the existence of tubular neighborhoods
 and the example at the end of \ref{para::4.4.3}, it follows that
 $\mathcal{S}'$ is 
 determined by $\ovl{S}'$ and by the ramification index $n_{ij}$ at each
 divisor $Z_{ij}$ above $Z_i$ (for every $i$): namely,
 $e'_{ij}=e_i/n_{ij}$. The first assertion then follows immediately from
 theorem \ref{thm::2.1.10}.  The second assertion is a straightforward
 consequence 
 of this characterization of the image.
\end{proof}

 \medskip 
 Let us fix a geometric base-point $\ovl{s}$ of $S
 =\ovl{S}\setminus
 \,(\bigcup Z_i\cup\ovl{S}^{\mathrm{sing}})$ (in dimension one, we could also
 choose a tangential base-point in the sense of 2.2).

\begin{dfn}\label{def::4.5.3} The orbifold fundamental group pointed at
 $\ovl{s}$ 
\begin{equation*}
 \pi_1^{\mathrm{orb}}(\mathcal{S},\ovl{s})
\end{equation*}
 is the topological group attached to $\Cov_\mathcal{S}$ viewed as a
 subcategory of $\Cov_S^{\mathrm{et}}$, \cf 1.4. 
\end{dfn}
This is a prodiscrete group (\ref{cor::1.4.7}). By \ref{thm::1.4.5}, one
has a commutative 
square of functors
\begin{equation*}
 \begin{CD}
  \left\{
  \begin{matrix}
   \text{disjoint unions of}\\
   \text{orbifold-coverings of $S$}
  \end{matrix}
  \right\}
    @>\subset >>
  \left\{
  \begin{matrix}
   \text{disjoint unions of}\\
   \text{temperate coverings of $S$}
  \end{matrix}
  \right\}\\
  @V\wr VV @VV\wr V\\
  \{\pi_1^{\mathrm{orb}}(\mathcal{S},s)\text{-sets}\}
  @>\subset >>
  \{\pi_1^{\mathrm{temp}}(S,s)\text{-sets}\}.
 \end{CD}
\end{equation*}
 Galois orbifold-coverings correspond to open normal subgroups of $
 \pi_1^{\mathrm{orb}}$.

\medskip\noindent{\it Example}. For $r\in | K^\times|$, we set $\matheur{D}_r=
\matheur{D}(0,r^+),\;\matheur{D}_r^\ast=\matheur{D}_r\setminus \{0\},$ and consider the orbifold
$\mathcal{S}_e = (\matheur{D},(0,e))$ (multiplicity $e$ at the origin). Then by
the same calculation as in \ref{pro::2.3.7}, one finds

\begin{equation*}
\pi_1^{\mathrm{orb}}(\mathcal{S}_e)=\pi_1^{\mathrm{alg}}(\matheur{D})\times \Z/e\Z 
\end{equation*} 
(we drop the base point in the notation since the group is abelian, \cf
\ref{pro::1.4.4}).

This can also be expressed in the following way: let $\gamma$ be a
topological generator of the $\widehat{\Z}(1)$-factor in
$\pi_1^{\mathrm{temp}}(\matheur{D}^\ast)$ (the factor corresponding to Kummer
coverings). Then \begin{equation*} \pi_1^{\mathrm{orb}}(\mathcal{S}_e
)=\pi_1^{\mathrm{temp}}(\matheur{D}^\ast,\ovl{s})/\an{\gamma^e}^- , \end{equation*}
where $\an{\gamma^e}^-$ denotes the closure of the (normal) subgroup
generated by $\gamma$.

 
\begin{para}\label{para::4.5.4} By Kiehl's theorem, $\ovl{S}$ is locally
 isomorphic to a product $\matheur{D}^\ast\times Z_i^{\mathrm{smooth}}$ in the
 neighborhood 
 of the non-singular part of $Z_i$; whence a homomorphism
 $\pi_1^{\mathrm{temp}}(\matheur{D}^\ast)\to \pi_1^{\mathrm{temp}}(S,\ovl{s})$,
 well-defined up to conjugation. Let us denote by $\gamma_i$ the image in
 $\pi_1^{\mathrm{temp}}(S,\ovl{s})$ of a topological generator of the
 $\widehat{\Z}(1)$-factor of $\pi_1^{\mathrm{temp}}(\matheur{D}^\ast)$ (``local
 monodromy'' at $Z_i$). The closure $\an{\gamma_i^{e_i}}^-$ of the normal
 subgroup of $\pi_1^{\mathrm{temp}}(S,\ovl{s})$ generated by the
 $\gamma_i^{e_i}$'s is a well-defined closed subgroup, independent of the
 choice of the $\gamma_i$'s.
\end{para}

\begin{thm}\label{thm::4.5.5}
 \begin{enumerate}
  \renewcommand{\theenumi}{\alph{enumi}}
  \item One has $\displaystyle
	\pi_1^{\mathrm{temp}}(S,\ovl{s}) \cong \limproj_{(e_i)}
	\pi_1^{\mathrm{orb}}((\ovl{S},(Z_i,e_i)),\ovl{s}),$ where the
	limit is taken according to the ordering by divisibility.
  \item The homomorphism $\pi_1^{\mathrm{temp}}(S,\ovl{s})\to
	\pi_1^{\mathrm{orb}}(\mathcal{S},\ovl{s})$ induces an injective strict
	homomorphism
	\begin{equation*}
	 \pi_1^{\mathrm{temp}}(S,\ovl{s})/\an{\gamma_i^{e_i}}^-\,\to \,
	  \pi_1^{\mathrm{orb}}(\mathcal{S},\ovl{s})
	\end{equation*}
	with dense image. This is an isomorphism if $S$ is the analytification
	of a smooth algebraic $K$-variety $S^{\mathrm{alg}}$; in particular
	$\pi_1^{\mathrm{temp}}(S,\ovl{s})/\an{\gamma_i }^-\,\cong \,
	\pi_1^{\mathrm{temp}}({\ovl{S}},\ovl{s}) $.
 \end{enumerate}
\end{thm}
 
\begin{proof}
  (a) reflects the fact that $\Cov_S^{\mathrm{temp}}= \liminj
 \Cov_{(\ovl{S},(Z_i,e_i)),\ovl{s})}$ which follows from the previous
 proposition.

  (b) The previous proposition shows more precisely that a
 $\pi_1^{\mathrm{temp}}(S,\ovl{s})$-set comes from a 
 $\pi_1^{\mathrm{orb}}(\mathcal{S},\ovl{s})$-set if and only if the
 $\gamma_i^{e_i}$'s act trivially. 
 
 \noindent We shall use lemma \ref{lem::1.4.9}. By construction
 of the topology, open subgroups of
 $\pi_1^{\mathrm{temp}}(S,\ovl{s})$ and
 $\pi_1^{\mathrm{orb}}(\mathcal{S},\ovl{s})$ respectively form a fundamental
 system of neighborhoods of 
 $1$. In fact, this is even true for normal open subgroups
 (existence of Galois closures). It follows that 
 $\pi_1^{\mathrm{temp}}(S,\ovl{s})/\an{\gamma_i^{e_i}}^-$ has the same
 property. The first part of b) then follows from
 \ref{lem::1.4.9}. (ii). The 
 second part follows from \ref{lem::1.4.9}. (iii) and (iv), since
 we have 
 shown that the quotient of $\pi_1^{\mathrm{temp}}(S,\ovl{s})$ by
 any closed normal subgroup (\eg $\an{\gamma_i^{e_i}}^-$) is
 complete when $S$ is algebraic (\ref{pro::2.1.6}).
\end{proof}

\begin{rems}
 (a) \ref{pro::2.3.9} is a special case of the
 last assertion of \ref{thm::4.5.5}. 

 (b) When $S$ is algebraic, the profinite completion of
 $\pi_1^{\mathrm{orb}}(\mathcal{S},\ovl{s})$ admits an 
 algebraic description. 

 (c) The situation is analogous to the classical one (over $\C$, \cf 4.2),
 with at least two noteworthy differences: 

 \begin{itemize}
 \item since there is no universal orbifold-covering, the geometrically
       meaningful object is not $\pi_1^{\mathrm{orb}}$ itself, but its
       discrete quotients,
 \item the non-archimedean $\pi_1^{\mathrm{orb}}$ is a finer invariant
       which takes into account the relative ``position'' of the $Z_i$. 
       For instance, if $\ovl{S}= \P^1_{\C_p}$, the classical
       $\pi_1^{\mathrm{orb}}$ 
       depends only on the numbers $e_i$, not on the branch 
       points $Z_i$ themselves. In the $p$-adic case, this is not the case. For
       instance, if $Z=\{0,1,\lambda,\infty\}$ and every 
       $e_i=2$, then the Legendre elliptic curve with parameter $\lambda$
       provides a global orbifold chart, 
       and it follows from \ref{para::2.3.2} that the orbifold fundamental
       group is either
       \begin{itemize}
	\item a (compact) semi-product of $\Z/2\Z$ by $\widehat{\Z}\times
	      \widehat{\Z}$ if $| \lambda(1-\lambda)| =1$,
	\item a (non-compact) semi-product of $\Z/2\Z$ by $ \Z\times
	      \widehat{\Z}$ if $| \lambda(1-\lambda)| \neq 1$.
       \end{itemize}
 \end{itemize}
\end{rems}

\begin{para}\label{para::4.5.6} Let $\phi:\,\ovl{S}'\to \ovl{S}$ be a global
 orbifold chart, such that the induced map $S'\to S$ is a 
finite Galois \'etale covering with group $G$. Let $\ovl{s}'$ be a
 geometric point of $S'$ above $\ovl{s}$.
\end{para}

\begin{pro}\label{pro::4.5.7} The universal topological covering
 $\widetilde{\ovl{S}'}$ of ${\ovl{S}'}$ is a Galois orbifold-covering of
 $\mathcal{S}$. Let $\Gamma= \Aut(\widetilde{\ovl{S}'}/\mathcal{S})$
 denote the corresponding (discrete) Galois group. Then one has a
 commutative diagram of strict homomorphisms of 
 topological groups, in which the rows are exact sequences
 \begin{equation*}
  \begin{CD}
   1 @>>> \pi_1^{\mathrm{temp}}(\ovl{S}',\ovl{s}') @>>>
   \pi_1^{\mathrm{orb}}(\mathcal{S},\ovl{s}) @>>> G @>>> 1\\
   @. @VVV @VVV @VVV @.\\
   1 @>>> \pi_1^{\mathrm{top}}(\ovl{S}',\ovl{s}') @>>> \Gamma @>>> G @>>> 1.
  \end{CD}
 \end{equation*}
\end{pro}

\begin{proof}
 Let us denote by $\mathcal{S}'$ the orbifold $\ovl{S}'$ with all 
 $\phi^{-1}(Z_i)$ endowed with multiplicity 
 $e_i=1$. The pull-back by $\phi$ provides the vertical arrows of a
 commutative diagram of functors
 \begin{equation*}
 \begin{CD}
  \Cov_\mathcal{S} @>\subset >> \Cov_S^{\mathrm{temp}}\\
  @VVV @VVV\\
  \Cov_{\mathcal{S}'}=\Cov_{\ovl{S}'}^{\mathrm{temp}}@>\subset>>
  \Cov_{S'}^{\mathrm{temp}}.
 \end{CD}
 \end{equation*}
 One concludes by \ref{cor::1.4.12} that there is an exact sequence

 \begin{equation*}
 1\to \pi_1^{\mathrm{temp}}(\ovl{S}',\ovl{s}')\to
  \pi_1^{\mathrm{orb}}(\mathcal{S},\ovl{s}) \to G \to 1.
 \end{equation*}
 On the other hand, it follows from \ref{pro::1.1.3} that the restriction
 of $\widetilde{\ovl{S}'}$ above $S'$ is nothing 
 but the universal topological covering $\tilde{  S'}$ of ${S'}$. By
 \ref{lem::2.1.2}, $\tilde S'\to S$ is a Galois temperate 
 covering, and it follows from that and \ref{pro::4.5.2} that
 $\widetilde{\ovl{S}'}\to \mathcal{S}$ is a Galois orbifold-covering. It
 is then clear that the corresponding epimorphism
 $\pi_1^{\mathrm{orb}}(\mathcal{S},\ovl{s})\to \Gamma$ gives rise to the above
 commutative diagram.
\end{proof} 

\begin{para}\label{para::4.5.8} Assume that $S$ is the analytification of a
 smooth algebraic $K$-variety $S^{\mathrm{alg}}$. According to
 \ref{thm::3.5.12}, 
 discrete representations of $\pi_1^{\mathrm{temp}}(S,\ovl{s})$ give rise
 to algebraic vector bundles with regular integrable 
 connection $(M,\nabla)$ on $S^{\mathrm{alg}}$. 
\end{para}

\begin{pro}\label{pro::4.5.9} An algebraic vector bundle with regular integrable
connection $(M,\nabla)$ on $S^{\mathrm{alg}}$ comes from a discrete representation $\rho$ of $\pi_1^{\mathrm{orb}}(\mathcal{S},\ovl{s})$ if and
only if there is a finite global orbifold chart $\phi:\,\ovl{S}'\to \ovl{S}$ such that $(\phi^\ast(M))^\nabla$ extends to a
locally constant sheaf on $\ovl{S}'$.   
\end{pro}

This follows from \ref{para::3.5.6} and \ref{thm::4.5.5}. 

\medskip\noindent {\it Remark.} If $S'/S$ is Galois, one can restrict
$\rho$ to the open normal subgroup $ 
\pi_1^{\mathrm{temp}}(\ovl{S}',s')$ of
$\pi_1^{\mathrm{orb}}(\mathcal{S},\ovl{s})\,$ (\ref{pro::3.5.7}), and this restriction
factors through $\pi_1^{\mathrm{top}}(\ovl{S}',s')$.

 
 
\subsection{$p$-adic uniformizing differential equations.}\label{sub::4.6}

\begin{para}\label{para::4.6.1} We now develop the non-archimedean counterpart
 of 4.3. We assume that $\ovl{S}$ is a 
 smooth projective {\it curve} over $K$. Then
 $Z=\{\zeta_0,\ldots ,\zeta_n\}$ is just a finite set of points and
 $S=\ovl{S}\setminus Z$.

 We assume moreover that $\mathcal{S}$ is
uniformizable. To imitate the complex theory,
we shall assume that {\it there exists a finite Galois global orbifold
 chart $\phi: \ovl{S}'\to \ovl{S}$ such that $\ovl{S}'$ 
is a Mumford curve} (of any genus $\geq 0$). 

The universal topological covering $\widetilde{\ovl{S}'}$ of
 ${\ovl{S}'}$ will replace the universal 
 orbifold-covering, which does not exist in the non-archimedean
 situation. By assumption, there is an embedding  
  \begin{equation*}
\widetilde{\ovl{S}'} \subset \P^1,
\end{equation*}
 which we fix, so that the group $\Aut\, \widetilde{\ovl{S}'}\;$
 becomes a Lie subgroup of $PGL_2(K)$.

The orbifold-covering $\widetilde{\ovl{S}'}\to \ovl{S}$ gives rise to a
 projective discrete representation 
\begin{equation*}
\ovl{\rho}:\,\pi_1^{\mathrm{orb}}(\mathcal{S},s) \to PGL_2(K).
\end{equation*}
 Its image $\Gamma$ is contained in the normalizer in $PGL_2(K)$ of the 
Schottky group $\pi_1^{\mathrm{top}}({\ovl{S}'},s')$, hence discrete; it
 acts properly on $\widetilde{\ovl{S}'}$, with quotient space 
$\widetilde{\ovl{S}'}/\Gamma=\ovl{S}$.  
\end{para}

\begin{para}\label{para::4.6.2} We fix a non-constant rational function
 $z\in K(\ovl{S})$, and assume for simplicity that $z$ has no 
 pole on $S$. 
 
 \noindent Having identified $\widetilde{\ovl{S}'}$ with a domain of
 $\P^1$ allows us to consider $z: \tau \mapsto z(\tau)$ as a meromorphic
 function defined on that domain (and analytic on 
 $\widetilde{  S'}\subset \widetilde{\ovl{S}'}$). Notice that the derivation
 $\frac{d}{dz}$ extends canonically to $\mathcal{O}(\widetilde{S'})$.  
\end{para}

\begin{pro}\label{pro::4.6.3} There is a unique differential equation of
 the form
 \begin{equation*}
  y''- qy=0,\;\;\; q\in K(\ovl{S}),
 \end{equation*}
 the {\rm uniformizing differential equation}, such that a local inverse
 $\tau(z)$ of the map $z(\tau)$ is given by the quotient of two solutions in
 $\mathcal{O}(\widetilde{ S'})$ of this differential equation. It is
 fuchsian and has singularities only at $Z=\{\zeta_0,\ldots ,\zeta_n\}$. 

 \hfill\break If $\ovl{S}=\P^1$ and $\zeta_n=\infty$, the exponents are
 rational numbers given 
 by the Riemann scheme:
\begin{equation*}
 \begin{pmatrix}
  \zeta_0 & \ldots & \zeta_{n-1} & \zeta_n=\infty \\
  &&&\\\frac{1}{2}(1+\frac{1}{e_0})&
  \ldots &\frac{1}{2}(1+\frac{1}{e_{n-1}})& \frac{-1}{2}(1+\frac{1}{e_{n}}) \\
  \\
  \frac{1}{2}(1-\frac{1}{e_0})&
  \ldots &\frac{1}{2}(1-\frac{1}{e_{n-1}})& \frac{-1}{2}(1-\frac{1}{e_{n}}) \\
 \end{pmatrix}.
\end{equation*}
\end{pro}

 \noindent The proof is the same as in the complex case: $q$ is nothing
 but the Schwarz derivative $\{\tau, z\}= 
 (\frac{d\tau}{dz})^{1/2}\frac{d^2}{dz^2}[(\frac{d\tau}{dz})^{-1/2}]$,
 which is single-valued on $\ovl{S}$. A basis of solutions of the
 uniformizing differential equation is given by
 $(\frac{d\tau}{dz})^{-1/2}$ and $\tau.(\frac{d\tau}{dz})^{-1/2}$.

\begin{pro}\label{pro::4.6.4} The projective connection attached to the
 uniformizing differential corresponds to the projective discrete
 representation $\ovl{\rho}$ in the sense of \ref{para::3.5.13}: it is
 in the image 
 of the temperate \'etale Riemann-Hilbert functor and its projective
 monodromy group is the discrete group $\Gamma$.
\end{pro}

\begin{proof}
 A basis of solutions of the symmetric square of the uniformizing 
 differential equation is given by $\frac{1}{{d\tau/
 dz}},\,\frac{\tau}{{d\tau/ dz}},\,\frac{\tau^2}{{d\tau/dz}} $. These 
 are elements of $\mathcal{O}(\widetilde{  S'})$ (in fact multivalued 
 meromorphic functions on $\ovl{S}'$). Criterion \ref{pro::4.5.9} then
 shows that 
 this symmetric square comes from a discrete representation of 
 $\pi_1^{\mathrm{orb}}(\mathcal{S},s) $. It is clear that this representation is 
 nothing but the composition of $\ovl{\rho}$ with the 
 standard homomorphism $PGL_2\to GL_3$. 
\end{proof}

\noindent [One could also argue by using an analogue of \ref{pro::4.5.9} for
projective connections.]
 
\medskip {\it Remark.} The connection itself attached to the
uniformizing differential equation corresponds to a discrete
representation of $\pi^{\mathrm{temp}}(S,\ovl{s})$ (\cf
\ref{pro::3.5.14}), which 
does not necessarily factor through $\pi_1^{\mathrm{orb}}(\mathcal{S},s) $ [it
factors through an orbifold fundamental group, up to replacing the
$e_i$'s by appropriate multiples].

 
\subsection{$p$-adic Shimura orbifolds and uniformizing differential
  equations.}\label{sub::4.7}

\begin{para}\label{para::4.7.1} We are now ready for a synthesis, relating the
 theory of $p$-adic period mappings surveyed in II with the 
theory of $p$-adic orbifolds. 

Let us consider a maximal order
$\mathcal{B}$ in an indefinite quaternion algebra
$B$ over
$\Q$, endowed with a positive involution $\ast$. ``Fake elliptic
 curves'' are $\ast$-polarized abelian surfaces with multiplication by
$\mathcal{B}$. They are classified by a smooth projective Shimura curve
$\mathcal{X}^+$ defined over $\Q$, \cf
 II.\ref{para:7.4.2}, which should actually be considered as an orbifold 
 $(\mathcal{X}^+, (\zeta_i,e_i))$, because the action of $\Gamma^+$ on
the Poincar\'e upper half-place $\mathfrak{h}$ has fixed points (the
$\zeta_i$'s are their images in $\mathcal{X}^+$).

 The Gauss-Manin connection $(\mathcal{H}, \nabla_{\mathrm{GM}})$
 attached to the 
 ``universal'' fake elliptic curve (in the sense of 
 orbifolds) is a connection with logarithmic poles at the $\zeta_i$'s. It
 has been analyzed in  II.\ref{sub:7.6}. The choice 
 of a quadratic splitting field
 $\Q(\sqrt{-d})/\Q$ for $B$ decomposes $(\mathcal{H},
 \nabla_{\mathrm{GM}})$ into 
 two pieces of rank two (the ``$+$ and $-$ 
 eigenspaces'' for the action of $\sqrt{-d}$).  It follows from the
 theory of the complex period mapping that a multivalued 
 section $\tau :\,\mathcal{X}^+(\C)\setminus 
 \{\zeta_i\}\to \mathfrak{h}$ of the projection $\mathfrak{h}\to
 \mathcal{X}^+(\C)$ is given by a quotient of two non-proportional 
solutions of $\nabla_{\mathrm{GM},+}$. One derives that {\it the uniformizing
 differential equation for $(\mathcal{X}^+,
(\zeta_i,e_i))$ is projectively equivalent to $\nabla_{\mathrm{GM},+}$}.

On the other hand, there is a canonical $\ovl{\Q}$-structure in
the $\C$-space of solutions at a given point $s$ (the Betti $\mathrm{H}^1( -
 /\ovl{\Q})$ of the fake elliptic curve with parameter $s$), which is
stable under (global) monodromy.
\end{para}
 
\begin{para}\label{para::4.7.2} According to the \v{C}erednik-Drinfeld theory
 surveyed in II.\ref{sub:7.4}, \ref{sub:7.6} (and the description of
 the $p$-adic period mapping given there as a quotient of solutions of
 the Gauss-Manin connection), there is an entirely parallel $p$-adic
 description of this geometric situation for any critical prime $p$, \ie
 any $p$ which divides the discriminant of $B$: $\mathcal{X}^+(\C_p)$ is
 a $p$-adic orbifold of the kind considered in 4.6, and {\it the
 $p$-adic uniformizing differential equation for $(\mathcal{X}^+,
 (\zeta_i,e_i))$ is projectively equivalent to
 $\nabla_{\mathrm{GM},+}$}, viewed 
 as a $p$-adic connection (\cf II.\ref{cor:7.2.6}, II.\ref{para:7.2.7}).

 By the theory of $p$-adic Betti lattices (\cf  II.\ref{para:7.6.2}),
 there is a canonical $\ovl{\Q}$-structure in the $\C_p$-space of
 solutions at a given point $s$, which is stable under (global) monodromy.
\end{para}

\begin{para}\label{para::4.7.3} Following  II.\ref{sub:7.5}, all this generalizes to
 Shimura curves defined by quaternion algebras $B$ over totally real
number fields $E$, and a ``congruence subgroup'' $\Gamma$ (here, we drop
the general assumption of II.\ref{sec:7} that $\Gamma$ is small enough). We
recall that these Shimura curves are not modular if $E\neq \Q$, but are
twisted forms of modular Shimura curves defined by quaternion algebras
$B^\bullet = B\otimes_E K$ over suitable totally imaginary quadratic
extensions $K/E$.  One can thus still attach a Gauss-Manin connection to
such a situation, and this connection is related to the uniformizing
differential equations both in the complex and in the $v$-adic
sense\footnote{if $p$ divides the level of the congruence group, the
orbifold is not quite of the kind considered in 4.6: the top space is
not the set of ordinary points for a Schottky group, but a finite
\'etale covering of such a space, \cf II.\ref{para:7.4.7}; however, the
construction of the uniformizing differential equation can still be done
in the same way, due to the fact that the finite \'etale coverings of
the Drinfeld space which occur are $PGL_2(E_v)$-equivariant.}, for {\it
critical places $v$ of $E$}, \cf II.\ref{sub:7.5}, II.\ref{sub:7.6} (we
fix embeddings $\C \leftarrow \ovl{\Q} \to \C_p$ for every critical
$v\mid p$, compatible with the privileged places $\infty_0$ and $v$ of
$E$, and denote by $\ovl{B}$ the quaternion algebra over $E$ obtained
from $B$ by interchanging the local invariants at $\infty_0$ and
$v$). We summarize these results as follows, using the terminology of
orbifolds:


\end{para}

\begin{thm}\label{thm::4.7.4} The Shimura curve $\mathcal{X}_\Gamma$ is
 defined over $\ovl{\Q}$, and may be considered, in a canonical way,
 both as a complex orbifold and as a $p$-adic orbifold (for any critical
 $v\mid p$).

 \noindent The complex uniformizing differential equation is defined over
 $\ovl{\Q}$ and coincides with the $p$-adic uniformizing differential
 equation. It is projectively equivalent to a 
 direct summand of the Gauss-Manin connection: the period mapping (\resp
 the $p$-adic period mapping) is given by a suitable quotient
 $\tau $ (\resp $\tau_p$) of two non-proportional solutions. The
 projective monodromy group of the uniformizing differential 
 equation over $\C$ (\resp over $\C_p$) is $\Gamma$ (\resp is a discrete
 arithmetic group $\Gamma_p$ in $PGL_2(\C_p)$ 
 attached to the quaternion algebra $\ovl{B}$).

 \noindent Moreover, there is a canonical $\ovl{\Q}$-structure in the
 space of solutions of the complex (\resp $p$-adic) uniformizing
 differential equation at any $\ovl{\Q}$-point, which is stable under
 global monodromy.
\end{thm}
Let us recall that the $p$-adic projective monodromy group, based at a
geometric point $\ovl{s}$, can be understood in a 
sense close to the usual intuitive one, \cf \ref{para::3.5.6}. In loose
words: there is a smooth projective curve $S$ and a finite 
orbifold-covering
$S\to \mathcal{X}_\Gamma$ such that the projective monodromy is given by
the values taken by $\tau_p$ when one passes 
from a point $ s'$ lying above $ s$ to another $ s''$ by all possible
paths on the Berkovich space $S_{\C_p}^{\mathrm{an}}$.

\begin{para}\label{para::4.7.5} Let us illustrate this section with a concrete
 example (\cf \cite{poaptg} for more details). We first 
define a quadrangle
$ABCD$ with circular edges in $\mathfrak{h}$, as follows 
%%% PICTURE 3c
\begin{figure}[h]
 \begin{picture}(250,250)(0,-20)
  \put(0,-20){\includegraphics[scale=0.75,keepaspectratio,clip]{fig3-3c.eps}}
  \put(242,7){\makebox(0,0)[c]{{\tiny $\Sigma'$}}}
  \put(202,32){\makebox(0,0)[c]{{\footnotesize $\Sigma$}}}
  \put(180,80){\makebox(0,0)[c]{{\footnotesize $\Sigma''$}}}
  \put(182,-12){\makebox(0,0)[c]{{\small $-1$}}}
 \end{picture}
 \caption{}
\end{figure}
 
 $L$ is the half line of slope $-1/2$ through the origin,
 
 $\Sigma$ (\resp $\Sigma'$) is the unit half-circle centered at the
 origin (\resp at $2-\sqrt 3$),

 $D$ is the point $L\cap \Sigma'$,

 $\Sigma''$ is the half-circle centered on the horizontal axis,
 orthogonal to $\Sigma$ and passing through $D$, 

 $B$ (\resp $C$) is the intersection point of the vertical axis with
 $\Sigma$ (\resp $\Sigma'$), 

 $A$ is the intersection point of $\Sigma$ with $\Sigma'$ (specifically,
 $A=- \frac{1+2i}{\sqrt 5}$).

\medskip The quadrangle $ABCD$ is a hyperbolic quadrangle with angles
 $(\frac{\pi}{2},\frac{\pi}{2},\frac{\pi}{2},\frac{\pi}{6})$. The
 symmetries around the edges of this polygon generate a group of
 M\"obius transformations which contains a 
 subgroup $\Gamma^\ast$ of index $2$ of conformal transformations. This
 fuchsian group has already attracted the attention 
 of many a mathematician (J.F. Michon, M.F. Vigneras
 \cite[p. 123]{adadq}, F. Beukers, D. Krammer \cite{aeoaafg}, N. Elkies 
 \cite{elkies98:_shimur}, D. Kohel... ). It is $PSL_2(\R)$-conjugated to the 
 discrete group generated by the following unimodular 
matrices (up to sign)
\begin{equation*}
\begin{pmatrix}
 0&1\\-1&0
\end{pmatrix},\,
\begin{pmatrix}
 0&-\frac{\sqrt 5+1}{2}\\
 \frac{\sqrt 5-1}{2}&0
\end{pmatrix},\,
\begin{pmatrix}
 -\sqrt 3& 1+\sqrt 5 \\ 1-\sqrt 5&\sqrt 3
\end{pmatrix},\,
\begin{pmatrix}
 \frac{\sqrt 3+\sqrt {15}}{2}& -2\\2&\frac{\sqrt 3-\sqrt {15}}{2}
\end{pmatrix}
\end{equation*}
of respective orders $2,2,2,6$, whose product in this order is $1$. The
 orbifold quotient of $\mathfrak{h}$ by $\Gamma^\ast$ has 
been determined by Krammer --- and, later but independently, by Elkies: it is 
\begin{equation*}
\bigl(\P^1; (0,e_0=2), (1,e_1=2), (81, e_{81}=2), (\infty , e_\infty =
 6)\bigr),
\end{equation*}
 the points $A,B,C,D$ being mapped to $0,1,81,\infty$ respectively. 

Moreover, these authors have computed the uniformizing
equation\footnote{it was also found, earlier, by D. and G. Chudnovsky 
\cite[p.193]{chudnovsky89:_comput}, who did not give 
the proof.}. It is actually easier to write down the differential 
equation obtained from the uniformizing equation by 
multiplying all solutions by $P^{1/4}$, where $P=z(z-1)(z-81)$.  It is 
\begin{equation*}
 18P. y'' + 9 P' y' + (z-9). y =0.\tag{$\ast$}
\end{equation*}
\end{para}

\begin{para}\label{para::4.7.6} In general, it is quite hard to compute
 uniformizing differential equations (due to the presence of
 accessory parameters, when there are
 more than three singularities). The method of Krammer and Elkies takes
 advantage of the fact that the orbifold
 under consideration is actually a Shimura orbifold (the method, which
 seems to go back to H.P.F. Swinnerton-Dyer, exploits 
 the Hecke correspondences). 

 In fact, $\Gamma^\ast$ coincides with the group denoted by the same
 symbol in  II.\ref{para:7.4.2}, namely $\{g\in B^\times \mid 
 g\mathcal{B}=\mathcal{B}g, \Nr(g)>0\}$, for the quaternion algebra $B$ of
 discriminant $3.5$ over $\Q$  \cite[IV.3.B,C]{adadq}
 (\cf \cite[8]{poaptg} for more details on the coincidence of $\Gamma$
 with the group denoted by $W^+$ in \cite{aeoaafg}, up to 
 conjugation), and is the orbifold fundamental group of the Shimura
 orbifold $\mathcal{X}^\ast$:
 \begin{equation*}
  \mathcal{X}^\ast(\C)= \mathfrak{h}/\Gamma^\ast.
 \end{equation*} 
\end{para}

\begin{para}\label{para::4.7.7} Let us now consider the critical prime
 $5$ (similar phenomena occur for $p=3$). Note that there is a 
``confluence'' between the singularities $1$ and $81$ in characteristic
 $5$. Let $\ovl{B}$ be the (definite) quaternion 
algebra with discriminant $3$, and let $\ovl{\mathcal{B}}$ be a maximal
 order ($\cong \Z\bigl[1,i,\frac{i+j}{2},\frac{1+ij}{2}\bigr]$ with 
$i^2=-1,j^2=-3,ij=-ji\,$). Let $\Gamma_5^\ast$ be the image in
 $PGL_2(\Q_5)$ of the normalizer $\{g\in \ovl{B}^\times \mid 
 g\ovl{\mathcal{B}}[\frac{1}{5}]=\ovl{\mathcal{B}}[\frac{1}{5}]g\}$. 
 This is a discrete subgroup, and one has the \v{C}erednik uniformization 
\begin{equation*}
 \mathcal{X}^\ast(\C_p)= \Omega_{\C_5}/\Gamma_5^\ast
\end{equation*} 
 by the Drinfeld space (II.\ref{para:7.4.3}), which represents
 $\mathcal{X}^\ast$ as a $p$-adic orbifold.

\medskip It follows from what precedes that viewed as a $5$-adic
 differential equation, $(\ast)$ is obtained from the $p$-adic
differential uniformizing equation attached to $\mathcal{X}^\ast$ by
multiplying all solutions by $(z(z-1)(z-81))^{1/4}$. In particular, it
has projective global monodromy group $\Gamma_5^\ast$. This group is
generated by the {\it conjugacy classes} of local monodromies at
$0,1,81$, as follows from \ref{pro::2.3.9}\footnote{however, it is not
clear that 
it is generated by three such elements, as is stated erroneously in
loc.~cit.}.
\end{para}

\newpage
\section{$p$-adic triangle groups.}\label{sec::5}
\begin{abst}
In loose words, $p$-adic triangle groups are projective monodromy groups
of those hypergeometric differential equations which are in the image of
the \'etale Riemann-Hilbert functor. We give a purely geometric
description of them, and then construct all ``arithmetic'' $p$-adic
triangle groups using the \v{C}erednik-Drinfeld uniformization
of Shimura curves.
\end{abst}

\subsection{Review of Schwarz triangle groups.}\label{sub::5.1}

\begin{para}\label{para::5.1.1} A {\it Schwarz orbifold} is a complex orbifold
 of the form 
 \begin{equation*}
  \mathcal{S}(e_0,e_1,e_\infty)=(\P^1; (0,
   e_0), (1,e_1), (\infty, e_{\infty})).
 \end{equation*}
 It is uniformizable. The universal orbifold-covering is $\P^1, \A^1$ or 
 the open unit disk $\matheur{D}$ according to whether $\sum  \frac{1}{e_i}$ is
 $>1,\;=1,\;$ or $<1$ respectively.

 \noindent We fix a base-point $\ovl{s}$ (which can be a tangential
 base-point).
 
\medskip\noindent The {\it Schwarz triangle group\footnote{in conformity
 with the previous chapter, we do 
 not allow the $e_i$'s to be infinite. Thus we restrict our attention to
 the so-called cocompact triangle groups.}} $\Delta(e_0,e_1,e_\infty)$
 may be defined as $\pi_1^{\mathrm{orb}}(\mathcal{S}(e_0,e_1,e_\infty),s)$.

\noindent As an abstract group, it has generators 
$\gamma_0,\gamma_1, \gamma_\infty$ and relations
 $\gamma_0^{e_0}=\gamma_1^{e_1}=\gamma_\infty^{e_\infty}= \Pi \gamma_i = 
1$. A permutation of $e_0,e_1,e_\infty)$ does not change the group, only
 the orbifold by a homography interchanging 
 $0,1$ and $\infty$.  
\end{para}

\begin{para}\label{para::5.1.2} The uniformizing differential equation is
 fuchsian with singularities at $0,1,\infty$, and Riemann scheme 
\begin{equation*}
\begin{pmatrix}
 0 &  1 & \infty \\
 &&&\\
 \frac{1}{2}(1+\frac{1}{e_0})& \frac{1}{2}(1+\frac{1}{e_{1}})&
 \frac{-1}{2}(1+\frac{1}{e_{\infty}}) \\ 
 &&& \\
 \frac{1}{2}(1-\frac{1}{e_0})& \frac{1}{2}(1-\frac{1}{e_{1}})&
 \frac{-1}{2}(1-\frac{1}{e_{\infty}})
\end{pmatrix} .
\end{equation*}

 \medskip Multiplying solutions of this equation by  
\begin{equation*}
 z^{\frac{1}{2}(-1+\frac{1}{e_0})}(z-1)^{\frac{1}{2}(-1+\frac{1}{e_1})}
\end{equation*}
transforms this equation into a (projectively equivalent) differential equation of hypergeometric type
\begin{equation*}
 HGDE(a,b,c)\;\;\;\;\;\;\;\;\;\;\;\;\;\; z(z-1)y''+(c-(a+b+1)z)y'-ab y=0,  
\end{equation*}
with $1-c=\frac{1}{e_0},\;c-a-b=\frac{1}{e_1},\; \pm(a-b)=\frac{1}{e_\infty}$, whose Riemann scheme is    
 \begin{equation*}
  \begin{pmatrix}
   0 &  1 & \infty \\
   &&&\\
   0& 0& \frac{1}{2}(1-\frac{1}{e_0}-\frac{1}{e_1}+\frac{1}{e_{\infty}}) \\
   &&&\\
   \frac{1}{e_0}&\frac{1}{e_{1}}&
   \frac{1}{2}(1-\frac{1}{e_0}-\frac{1}{e_1}-\frac{1}{e_{\infty}})
  \end{pmatrix} .
 \end{equation*}
 The triangle group $\Delta(e_0,e_1,e_\infty)$ may thus be interpreted
 as the {\it projective monodromy group pointed at $\ovl{s}$ of this
 hypergeometric differential equation}.
\end{para}

\begin{para}\label{para::5.1.3} The quotient $\tau$ of a suitable basis of
 solutions of $ HGDE(a,b,c)$ sends the upper half-plane bijectively 
onto a curvilinear triangle $\mathcal{T}$ with vertices
 $\tau(0),\tau(1),\tau(\infty)$ and respective angles 
$\pi/e_0,\pi/e_1,\pi/e_\infty$. By the Schwarz reflection principle, the
 map $\tau$ can be extended to the lower half-plane 
through each of the ``segments'' $]0,1[$, $]1,\infty[$, $]\infty,0[$ by
 conformal reflections $s_\infty,s_0,s_1$ respectively. The products 
$\gamma_0=s_1s_\infty,\,\gamma_1=s_\infty s_0,\,\gamma_\infty=s_0s_1 $
 give the standard generators of $\Delta(e_0,e_1,e_\infty)$. By
 repeating these reflections, the occuring curvilinear triangles do not
 overlap (except on their edges) and provide a tiling of the universal
orbifold-covering ($\P^1, \A^1$ or $\matheur{D}$).
\end{para}

\begin{para}\label{para::5.1.4} A similar situation arises for any
 hypergeometric differential equation with rational parameters
 $(a,b,c)$.  However, if some of the differences of the exponents
$1-c,c-a-b,a-b$ is not the inverse of an integer, the triangles which
occur by iterating Schwarz reflections may overlap. Moreover, if the
local projective monodromies are of respective order $e_0, e_1,
e_\infty$ at $0,1,\infty$, the global projective monodromy group need
not be $\Delta(e_0,e_1,e_\infty)$. This is related to the fact that a
homomorphic image of $\Delta(e_0,e_1,e_\infty)$ which preserves the
order of $\gamma_0,\gamma_1,\gamma_\infty$ need not be isomorphic to
$\Delta(e_0,e_1,e_\infty)$ (\cf \cite{greenberg81:_homom_psl} for a
detailed discussion 
of homomorphic images of triangle groups\footnote{we are indebted to
M. Kapovich for this reference.}).
\end{para}

\begin{para}\label{para::5.1.5} The case $\sum  \frac{1}{e_i}>1$ (universal
 covering $\P^1$) corresponds to the Schwarz list of $HGDE$ with finite
 (projective) monodromy groups: the list of triples
 $(e_0,e_1,e_\infty)$ is, up to permutation:

 \medskip
 $(2,2,n)$ corresponding to a dihedral group,

 $(2,3,3)$ corresponding to the tetrahedral group,

$(2,3,4)$ corresponding to the octahedral group,

 $(2,3,5)$ corresponding to the icosahedral group.
 
 \medskip In the case $\sum  \frac{1}{e_i}=1$ (universal covering $\A^1$), there is a global orbifold chart given by an elliptic curve with complex
 multiplication. There are three possibilities for $(e_0,e_1,e_\infty)$, up to permutation:
 $(2,3,6)$,
 $(2,4,4)$,
 $(3,3,3)$.
 
 \medskip All the other triples belong to the ``hyperbolic case'':
 $\Delta(e_0,e_1,e_\infty)$ is a fuchsian group.
 
 For more detail on these topics, we refer to the  good treatises
 \cite{f2}, \cite{ntatg}, \cite{hfml}.

\end{para}

\begin{para}\label{para::5.1.6} As was remarked by J. Wolfart, Belyi's
 criterion for algebraic curves to be defined over $\ovl{\Q}$ can be
 expressed in terms of Schwarz orbifolds as follows:

 a complex projective curve $\ovl{S}$ is defined over $\ovl{\Q}$ if and
 only if it is the base of an orbifold-covering $(\ovl{S}, (Z_i,e_i))$
 of some Schwarz orbifold.
\end{para}


\subsection{$p$-adic Schwarz orbifolds and triangle groups.}\label{sub::5.2}

\begin{para}\label{para::5.2.1} A {\it $p$-adic Schwarz orbifold} is an
 orbifold over some complete subfield $K$ of $\C_p$  of the form 
\begin{equation*}
\mathcal{S}(e_0,e_1,e_\infty)=(\P^1, (0, e_0), (1,e_1), (\infty, e_{\infty})).
\end{equation*}
\noindent Up to replacing $K$ by a finite extension, it is uniformizable
 (\cf example in \ref{para::4.4.4}). However, there is no universal
 orbifold-covering (\ref{para::4.4.4}).
 
We want our $p$-adic triangle groups to be {\it discrete}
groups of fractional linear $p$-adic transformations. Thus it is not
 possible in general to define them as orbifold 
fundamental groups, but only as appropriate discrete quotients. This
 leads us to imitate the alternative definition \ref{para::5.1.2} 
of triangle groups as projective hypergeometric monodromy groups. 
\end{para}

\begin{para}\label{para::5.2.2} Let $e_0, e_1, e_\infty$ be integers $>1$, and
 define  
\begin{equation*}
 a=
 \frac{1}{2}\Bigl(1-\frac{1}{e_0}-\frac{1}{e_1}+\frac{1}{e_{\infty}}\Bigr),\,
 b=
 \frac{1}{2}\Bigl(1-\frac{1}{e_0}-\frac{1}{e_1}-\frac{1}{e_{\infty}}\Bigr),\, 
 c= 1- \frac{1}{e_0}.
\end{equation*} 
 We now consider the hypergeometric differential equation 
$HGDE(a,b,c)$ as a $p$-adic differential equation.
\end{para}

\begin{dfn}\label{def::5.2.3} Assume that $ HGDE(a,b,c) $ is in the
 image of the \'etale Riemann-Hilbert functor, 
\ie the \'etale sheaf of germs
 of solutions is locally constant on $\P^1_{\C_p}\setminus\{0,1,\infty\}$
 (\cf \ref{thm::3.4.6}).  If the corresponding projective monodromy group
\begin{equation*}
\Im\bigl[\pi_1^{\mathrm{et}}(\P^1\setminus\{0,1,\infty\},\ovl{s})\to
 PGL_2(\C_p)\bigr]
\end{equation*} is discrete and finitely generated, it is then called
a {\rm $p$-adic triangle group}, and denoted by
$\Delta_p(e_0,e_1,e_\infty)$.

\medskip
In the sequel, we say that ``{\it $\Delta_p(e_0,e_1,e_\infty)$ exists}''
to express that the assumptions of \ref{def::5.2.3} are 
fulfilled (this depends on the four numbers $p,e_0,e_1,e_\infty$).

\end{dfn}

\begin{pro}\label{pro::5.2.4}
 \begin{enumerate}
  \renewcommand{\theenumi}{\alph{enumi}}
  \item If $\Delta_p(e_0,e_1,e_\infty)$ exists,
	there exists a finite Galois covering 
	$\phi:\,\ovl{S}' \to \P^1$ ramified above $\{0,1,\infty\}$, such that
	the pull-back of $ HGDE(a,b,c) $ by $\phi$ has a full 
	set of multivalued analytic solutions over
	$S'=\phi^{-1}(\P^1\setminus\{0,1,\infty\})$.
  \item If $\Delta_p(e_0,e_1,e_\infty)$ exists, the homomorphism
	\begin{equation*}
	 \Im[\pi_1^{\mathrm{et}}(\P^1\setminus\{0,1,\infty\},\ovl{s})\to
	 PGL_2(\C_p)]
	\end{equation*}
	factors through a surjective homomorphism 
	\begin{equation*}
	 \pi_1^{\mathrm{orb}}(\mathcal{S}(e_0,e_1,e_\infty)) \to
	  \Delta_p(e_0,e_1,e_\infty). 
	\end{equation*}
  \item Conversely, if the projective connection attached to $HGDE(a,b,c)$
	comes from a projective representation
	\begin{equation*}
	 \pi_1^{\mathrm{orb}}(\mathcal{S}) \to  PGL_2(\C_p)
	\end{equation*}
	with discrete image, this image is a $p$-adic triangle group
	$\Delta_p(e_0,e_1,e_\infty)$.
  \item Every $p$-adic triangle group
	$\Delta_p(e_0,e_1,e_\infty)$ is a discontinuous subgroup of
	$PGL_2(K)$ for some finite extension $K$ of $\Q_p$.
 \end{enumerate}
\end{pro}

\begin{proof}
  (a) Because the Wronskian of $ HGDE(a,b,c) $ is algebraic, the
 determinant of the representation  
 $\rho:\,\pi_1^{\mathrm{et}}(\P^1_{\C_p} \setminus
 \{0,1,\infty\},\ovl{s})\to 
 GL_2(\C_p)$ factors through a finite representation of
 $\pi_1^{\mathrm{temp}}(\P^1_{\C_p}\setminus\{0,1,\infty\},\ovl{s})$.
 Under the assumptions of \ref{def::5.2.3}, $\Im\,\ovl{\rho}$ is
 finitely 
 generated, and it follows that $\Im \,\rho$ is finitely generated. By
 \ref{pro::3.5.10}, one derives that $\rho$ factors through
 $\pi_1^{\mathrm{temp}}(\P^1\setminus\{0,1,\infty\},\ovl{s})$. This
 implies a), \cf \ref{para::3.5.6} (where a concrete description
 of the monodromy is 
 also given).

   (b) We have just seen that $\rho$ factors through
 $\pi_1^{\mathrm{temp}}(\P^1_{\C_p}\setminus\{0,1,\infty\},\ovl{s})$. To
 check that the associated projective representation factors through 
 $\pi_1^{\mathrm{orb}}(\mathcal{S}(e_0,e_1,e_\infty))$ is a matter of
 projective local monodromy --- or more plainly, of difference of
 exponents --- at $0,1,\infty$ (\cf 4.5); this is clear.

  (c) Again because the Wronskian of $HGDE(a,b,c)$ is algebraic, the
 fact that the projective connection attached to $HGDE(a,b,c)$ comes
 from a projective representation of
 $\pi_1^{\mathrm{temp}}(\P^1\setminus\{0,1,\infty\},\ovl{s})$ implies
 that $ HGDE(a,b,c) $ itself comes from a representation of
 $\pi_1^{\mathrm{temp}}(\P^1\setminus\{0,1,\infty\},\ovl{s})$ (\cf the
 discussion in \ref{para::3.5.13}). By \ref{pro::3.5.10}, the
 image is finitely generated, and 
 the assertion follows.

  (d) If $\Delta_p(e_0,e_1,e_\infty)$ exists, there exists a finite
 Galois global orbifold chart $\ovl{S}'\to \P^1$ such that pull-back of
 $ HGDE(a,b,c) $ on the universal topological covering
 $\widetilde{\ovl{S}'}$ becomes trivial (as a connection). Notice that
 $\ovl{S}'$ is a projective curve defined over some number field, and a
 fortiori over some finite extension $K/\Q_p$.  Up to replacing $K$ by a
 finite extension, we may assume that $\widetilde{\ovl{S}'}$ is also
 defined over $K$ (and that the image in $\ovl{S}'$ of some geometric
 point above $\ovl{s}$ is $K$-rational). Thus $\rho$ factors through a
 representation of
 $\pi_1^{\mathrm{temp}}(\P^1_K\setminus\{0,1,\infty\},\ovl{s})\to
 GL_2(K)$, which makes clear that $\Delta_p(e_0,e_1,e_\infty)\subset
 PGL_2(K)$. Since $K$ is locally compact, any discrete subgroup of
 $PGL_2(K)$ is discontinuous \cite[1.6.4]{sgamc}.
\end{proof}

\medskip\noindent{\it Remark.} By (b), if $\Delta_p(e_0,e_1,e_\infty)$
exists, there {\it exists} a finite Galois global 
orbifold chart $\phi:\,\ovl{S}' \to \P^1$ such that the quotient of two
non-proportional solutions of $ HGDE(a,b,c) $ 
defines a multivalued meromorphic function on $\ovl{S}'$. However, one
should not believe that {\it any} finite Galois 
global orbifold chart $\phi:\,\ovl{S}' \to \P^1$ has this property.

\noindent Indeed, we shall see below that the diadic triangle group
$\Delta_2(4,4,4)$ exists. The Fermat curve $F_4$ over 
$\C_2$ provides a finite Galois global orbifold chart, with
$\pi_1^{\mathrm{top}}(F_4)=1$ (because $F_4$ has tree-like reduction). 
But the corresponding projective hypergeometric connection does not
become trivial over $F_4$, since $\Delta(4,4,4)$ is 
infinite.  

\begin{para}\label{para::5.2.5} The case $\sum \frac{1}{e_i}>1$ corresponds
 to the list of $p$-adic $HGDE$ with finite (projective) monodromy
 groups. This is the same as the Schwarz list: $(2,2,n)$, $(2,3,3)$,
$(2,3,4)$, $(2,3,5)$, up to permutation, and yields finite $p$-adic
triangle groups for every $p$ which are isomorphic to the corresponding
finite Schwarz triangle groups.

\medskip\noindent On the other hand, there is {\it no $p$-adic triangle
 group with $\sum  \frac{1}{e_i}=1$}. Indeed, there is a global
 algebraic orbifold chart given by an elliptic curve with complex
 multiplication.  Viewed over $\C_p$, these have good reduction, and it
 follows from \ref{para::4.5.6} and \ref{para::2.3.2} that the orbifold
 fundamental group is 
 compact. Any discrete representation would then be finite, and it would
 follow that the corresponding $HGDE$ has a basis of algebraic
 solutions. This is an algebraic property, which can be translated back
 over $\C$. One finds that such a $HGDE$ would belong to the Schwarz
 list, a contradiction.

\medskip In the sequel, we restrict our attention to the ``hyperbolic
 case'' $\sum  \frac{1}{e_i}<1$. It is not clear a 
 priori that there exists any hyperbolic $p$-adic triangle group.
\end{para}

\subsection{Mumford-Schwarz orbifolds.}\label{sub::5.3}

\begin{para}\label{para::5.3.1} Our aim is to give a more geometric
 description of $p$-adic triangle groups, independent of the theory of
 monodromy of $p$-adic differential equations. This involves the notion
 of Mumford-Schwarz orbifold.
\end{para}

\begin{dfn}\label{def::5.3.2} A Mumford-Schwarz orbifold is a $p$-adic
 Schwarz orbifold such that there exists a finite 
Galois global orbifold chart $\ovl{S}' \to \P^1$, $\ovl{S}'$ being a
 Mumford curve.  
\hfill\break In other
words, it is the orbifold quotient of a Mumford curve $\ovl{S}'$ by a
 finite group of automorphisms $G$ such that $\ovl{S}'/G= \P^1$ and the
Galois covering $\ovl{S}' \to \P^1$ is ramified exactly above $0,
1,\infty$ with ramification index $e_0,e_1,e_\infty$ respectively.
\end{dfn}

\medskip\noindent According to \ref{pro::4.5.7}, one has an exact sequence
\begin{equation*}
1\to \pi_1^{\mathrm{temp}}(\ovl{S}',\ovl{s}')\to
 \pi_1^{\mathrm{orb}}(\mathcal{S}(e_0,e_1,e_\infty),\ovl{s}) \to G \to 1. 
\end{equation*}

\noindent On the other hand, because $\ovl{S}'$ is a Mumford curve,
$\pi_1(\ovl{S}',\ovl{s}')$ is a Schottky group embedded in 
$PGL_2(K) $, and the universal topological covering
$\widetilde{\ovl{S}'}$ can be 
identified with the set $\Omega\subset \P^1_{\C_p}$ of ordinary points
for $\pi_1(\ovl{S}',\ovl{s}')$. The subgroup  
$\Gamma$ of all liftings of the elements of $G \subset \Aut (\ovl{S}')$
in $PGL_2(K)$ is a subgroup of finite index inside 
the normalizer of
$\pi_1(\ovl{S}',\ovl{s}')$ in $PGL_2(K)$; it sits in an extension 

\begin{equation*}
 1\to \pi_1^{\mathrm{top}}(\ovl{S}',\ovl{s}')\to \Gamma \to G \to 1.
\end{equation*}
{\it Remark.} A Mumford-Schwarz orbifold is hyperbolic if and only if some (\resp every) curve $\ovl{S}'$
as above has genus $\geq 2$. 

\begin{pro}\label{pro::5.3.3} Let $ \mathcal{S}=(\P^1; (0,
 e_0), (1,e_1), (\infty, e_{\infty})) $ be a Mumford-Schwarz orbifold. Then the
 $p$-adic triangle group $\Delta_p(e_0,e_1,e_\infty)$ exists, and
 coincides with the discrete group $\Gamma\;$ (in particular, $\Gamma$
 does not depend on the choice of $\ovl{S}'/ \P^1$). Moreover, one has a
 commutative diagram
 
 \begin{equation*}
  \begin{CD}
   1 @>>> \pi_1^{\mathrm{temp}}(\ovl{S}',\ovl{s}') @>>>
  \pi_1^{\mathrm{orb}}(\mathcal{S}(e_0,e_1,e_\infty),\ovl{s}) @>>> G @>>> 1\\
  @. @VVV @VVV @VVV @.\\
  1 @>>>\pi_1^{\mathrm{top}}(\ovl{S}',\ovl{s}')@>>>
  \Delta_p(e_0,e_1,e_\infty)@>>> G @>>> 1
  \end{CD}
 \end{equation*}
 with surjective vertical maps and exact arrows.
\end{pro}

\begin{proof}
  Let us consider the $p$-adic uniformizing differential equation
 attached to the orbifold-covering $\widetilde{\ovl{S}'}\to
 \mathcal{S}(e_0,e_1,e_\infty)$ (3.6). It is fuchsian with singularities
 at $0,1,\infty$, and Riemann scheme
 \begin{equation*}
 \begin{pmatrix}
 0 &  1 & \infty \\
 &&&\\
 \frac{1}{2}(1+\frac{1}{e_0})&  \frac{1}{2}(1+\frac{1}{e_{1}})&
 \frac{-1}{2}(1+\frac{1}{e_{\infty}}) \\
 &&& \\
 \frac{1}{2}(1-\frac{1}{e_0})&  \frac{1}{2}(1-\frac{1}{e_{1}})&
 \frac{-1}{2}(1-\frac{1}{e_{\infty}})     
 \end{pmatrix} .
 \end{equation*}

 \medskip Multiplying solutions of this equation by  
 \begin{equation*}
 z^{\frac{1}{2}(-1+\frac{1}{e_0})}(z-1)^{\frac{1}{2}(-1+\frac{1}{e_1})}
 \end{equation*}
 transforms the equation into a (projectively equivalent) $p$-adic
 differential equation of hypergeometric type $ HGDE(a,b,c) $ with
 $1-c=\frac{1}{e_0},\;c-a-b=\frac{1}{e_1},\;
 \pm(a-b)=\frac{1}{e_\infty}$.  By \ref{pro::4.6.4}, the projective connection
 attached to $ HGDE(a,b,c) $ is the projective connection attached to to
 the projective representation given by the composed homomorphism
 $\pi_1^{\mathrm{orb}}(\mathcal{S}(e_0,e_1,e_\infty),\ovl{s})\to \Gamma \to
 PGL_2(K)$, which has discrete image. By \ref{pro::5.2.4}. (b), we
 conclude that 
 $\Delta_p(e_0,e_1,e_\infty)$ exists and coincides with $\Gamma$. The
 commutative diagram in the statement comes from \ref{pro::4.5.7}.
\end{proof}

\begin{para}\label{para::5.3.4} F. Kato calls the $p$-adic triangle groups
 which arise in this way ``$p$-adic triangle groups of Mumford type''
\cite{kato:_schwar_mumfor}. Inspired by F. Herrlich's study of automorphisms of
Mumford curves (using graphs of groups), he has developed
geometrico-combinatorial techniques to construct Mumford-Schwarz
orbifolds, and even a complete classification scheme, including a
description of $p$-adic triangle groups ``of Mumford type'' as explicit
amalgamated sums. As a by-product, he has proved the following
\end{para}

\begin{thm}[F. Kato]\label{thm::5.3.5} For $p= 2, 3$ and $5$, there are
 infinitely many hyperbolic Mumford-Schwarz orbifolds.  For $p>5$, there
 is none.
\end{thm}

\begin{para}\label{para::5.3.6} The following theorem explains how to obtains
 all $p$-adic triangle groups from $p$-adic triangle groups ``of Mumford
 type''.
\end{para}

\begin{thm}\label{thm::5.3.7} $\Delta_p(e_0,e_1,e_\infty)$ exists if and
 only if $\mathcal{S}(e_0,e_1,e_\infty)$ is a finite orbifold-covering
 of some Mumford-Schwarz orbifold
 $\mathcal{S}(e'_0,e'_1,e'_\infty)$. Moreover, the latter can be chosen
 in such a way that
 $\Delta_p(e_0,e_1,e_\infty)=\Delta_p(e'_0,e'_1,e'_\infty)$ as a subgroup
 of $PGL_2(\C_p)$.
\end{thm}

\begin{proof}
 (a) Assume that we have a finite orbifold-covering $h:
 \mathcal{S}\to \mathcal{S}'$ of a Mumford-Schwarz orbifold by a 
 Schwarz orbifold. The pull-back by $h$ of the hypergeometric projective
 connection $ HGPC(a',b',c') $ related to 
 $e'_0,e'_1,e'_\infty$ as above is the hypergeometric projective
 connection $HGPC(a,b,c) $ related to $e_0,e_1,e_\infty$. 
 Because $\mathcal{S}'$ is a Mumford-Schwarz orbifold, its uniformizing
 differential equation, hence $ HGPC(a',b',c') $, comes 
 from a representation of the temperate fundamental group, and it follows
 that so does $ HGPC(a,b,c) $. Moreover, the monodromy group of
 $HGPC(a,b,c) $ is a subgroup of finite index of the monodromy group of
 $HGPC(a',b',c') $, which is the discrete finitely generated group
 $\Delta_p(e'_0,e'_1,e'_\infty)$. Therefore the monodromy 
 group of $HGPC(a,b,c) $ is itself discrete finitely generated, and is
 $\Delta_p(e_0,e_1,e_\infty)$.  

 (b) Conversely, assume that $\Delta_p(e_0,e_1,e_\infty)$ exists. Let 
 $\Sigma \to \P^1$ be the ramified covering corresponding to the quotient 
 $\Delta_p(e_0,e_1,e_\infty)$ of
 $\pi_1^{\mathrm{orb}}(\mathcal{S},\ovl{s}) $. By 
 definition of $\Delta_p(e_0,e_1,e_\infty)$ as projective monodromy
 group, any quotient of two non-proportional solutions of 
 $HGDE(a,b,c)$ defines an \'etale map
 \begin{equation*}
 \tau: \;\Sigma \to \P^1_{\C_p}
 \end{equation*} which is equivariant under
 $\Delta_p(e_0,e_1,e_\infty)$.

 On the other hand, $\Delta_p(e_0,e_1,e_\infty)$ being
 a finitely generated discontinuous subgroup of $PGL_2(K)$ for some
 finite extension $K/\Q_p$, it admits a normal 
 subgroup of finite index which is a Schottky group $\Gamma$. Let
 $\Omega\subset \P^1_{\C_p}$ denote the set of ordinary 
 points for $\Gamma$. Since $\Delta_p(e_0,e_1,e_\infty)$ normalizes
 $\Gamma$, it preserves $\Omega$ and the quotient 
 $H=\Delta_p(e_0,e_1,e_\infty)/\Gamma$ is a group of automorphisms of the
 Mumford curve $\Omega/\Gamma$. We denote by $T$ 
 the quotient of this Mumford curve by $H$ (a smooth projective curve).

 We claim that
 \begin{equation*}
  \tau(\Sigma) \subset \Omega.
 \end{equation*}
 Indeed, otherwise $\tau$ would have a value $\tau(s)$ which is a limit
 point for $\Gamma$, \ie an accumulation point of an infinite sequence
 of points $\gamma_n(x)$, $x\in \P^1$. Since $\tau$ is \'etale, hence
 open, we may assume that $x\in \Im\,\tau$, say $x=\tau(t)$. By
 $\Gamma$-equivariance of $\tau$, we see that $\tau(s)$ is an
 accumulation point of the infinite sequence $\tau(\gamma_n(t))$. This
 contradicts the fact that the set $\{\gamma_n(t)\}$ is discrete in
 $\Sigma$ and that $\tau$ is \'etale (hence locally invertible around
 $s$).

 Next, being $\Delta_p(e_0,e_1,e_\infty)$-equivariant, $\tau$ induces an
 analytic map
 \begin{equation*}
  \ovl{\tau}: \, \Sigma/\Delta_p(e_0,e_1,e_\infty)=\P^1
   \to \Omega/\Delta_p(e_0,e_1,e_\infty)=T.
 \end{equation*}
 This map is necessarily algebraic and finite, and $T\cong \P^1$ (we
 normalize the latter isomorphism by imposing that
 $\ovl{\tau}(\{0,1,\infty\}\subset \{0,1,\infty\}$). Notice that because
 $\tau$ is $\Delta_p(e_0,e_1,e_\infty)$-equivariant, 
 it follows that
 \begin{equation*}
  \tau(\Sigma)= \Omega 
 \end{equation*} 
 and that $\tau$ is a finite \'etale covering (not a topological
 covering, unless $\Sigma=\Omega$). We have a cartesian 
 square 
 \begin{equation*}
 \begin{CD}
  \Sigma @>>> \Omega \\
  @VVV @VVV\\
  \P^1 @>>> \P^1
 \end{CD}
 \end{equation*}
 in which the vertical maps are Galois
 orbifold-coverings with group $\Delta_p(e_0,e_1,e_\infty)$, and define
 orbifold quotients $\mathcal{S}$ and $\mathcal{S}'$ of $\Sigma$ and
 $\Omega$ respectively. Since $\tau$ is \'etale, we see that
 $\mathcal{S}$ is actually a finite orbifold-covering of
 $\mathcal{S}'$. This forces $\mathcal{S}'$ to be a Schwarz orbifold
 (like $\mathcal{S}$), hence a Mumford-Schwarz orbifold by construction.
\end{proof}

\begin{cor}\label{cor::5.3.7}
 \begin{enumerate}
  \renewcommand{\theenumi}{\alph{enumi}}  
  \item Assume that $\Delta_p(e'_0,e'_1,e'_\infty)$ is a $p$-adic
	triangle group ``of Mumford type''. Assume that
	$\Delta(e_0,e_1,e_\infty)$ is a subgroup of finite index of the
	corresponding ``complex'' triangle group
	$\Delta(e'_0,e'_1,e'_\infty)$. Then $\Delta_p(e_0,e_1,e_\infty)$ 
	exists, and can be realized as a subgroup of finite index of
	$\Delta_p(e'_0,e'_1,e'_\infty)$.
  \item If $\Delta_p(e_0,e_1,e_\infty)$ exists, and if the
	corresponding ``complex'' triangle group $\Delta(e_0,e_1,e_\infty)$ is
	maximal in its commensurability class of triangle groups, then
	$\Delta_p(e_0,e_1,e_\infty)$ is ``of Mumford type''.
 \end{enumerate}
\end{cor}
\begin{proof}
 (a) If $\Delta(e_0,e_1,e_\infty)$ is a subgroup of finite index of
 $\Delta(e'_0,e'_1,e'_\infty)$, there is a finite morphism
 $\P^1\to \P^1$, defined over a number field, which underlies a
 finite orbifold-covering $\mathcal{S}(e_0,e_1,e_\infty)\to
 \mathcal{S}(e'_0,e'_1,e'_\infty)$ in the complex sense, but also
 in the $p$-adic sense (for every embedding of the number field
 in $\C_p$). The assertion comes from the ``if'' part of the
 theorem.

 (b) If $\Delta_p(e_0,e_1,e_\infty)$ exists, then, similarly, by the
 ``only if'' part of \ref{para::5.3.6}, there is a finite
 morphism $\P^1\to 
 \P^1$ defined over a number field, which underlies a finite
 orbifold-covering $\mathcal{S}(e_0,e_1,e_\infty)\to
 \mathcal{S}(e'_0,e'_1,e'_\infty)$ in the $p$-adic sense, and
 also in the complex sense. Hence
 $\Delta(e_0,e_1,e_\infty)\subset \Delta(e'_0,e'_1,e'_\infty)$
 with equality if and only if
 $\mathcal{S}(e_0,e_1,e_\infty)=\mathcal{S}(e'_0,e'_1,e'_\infty)$.
\end{proof}

\begin{cor}\label{cor::5.3.8} There is no infinite $p$-adic triangle
 group for $p>5$.
\end{cor}
\begin{proof}
 This follows from \ref{thm::5.3.5} and \ref{para::5.3.6}.  
\end{proof}


\subsection{Arithmetic triangle groups. Review of Takeuchi's
  list.}\label{sub::5.4}

\begin{para}\label{para::5.4.1} (Cocompact) arithmetic triangle groups
 are hyperbolic triangle groups which are {\it commensurable with the
 unit groups of quaternion algebras $B$ over number fields}.

 In order to make this definition more precise, let us first place
 ourselves in greater generality. Let $E$ be a number field, and let $B$
 be a quaternion algebra over $E$. Let $V$ be a finite set of places of
 $E$ containing the places at infinity and at least one place (possibly
 infinite) where $B$ is unramified. Let us define
 \begin{equation*}
  G_{(V)} = \prod \,B_v^\times/E_v^\times = \prod \, PGL_2(E_v)
 \end{equation*} where the product runs over all places $v\in V$ which
 are unramified for
 $B$. Let $\mathcal{O}_{E}(V)$ be the ring of so-called $V$-integers of
 $E$ (elements which are integers at the places not in $V$),
 let $\mathcal{B}(V)$ be a maximal $\mathcal{O}_{E}(V)$-order in
 $B$, and let $PSL_1(\mathcal{B}(V))$ be the image in $G_{(V)}$ of the
 group of elements of $\mathcal{B}(V)$ of reduced norm $1$. 
 This is a cocompact discrete subgroup of $ G_{(V)}$, but the projection
 of $PSL_1(\mathcal{B}(V))$ in any proper factor of 
 $G_{(V)}$ is dense \cite[IV.1.1]{adadq}. 

 Any subgroup of $G_{(V)}$ which is commensurable to
 $PSL_1(\mathcal{B}(V))$ is called an {\it arithmetic subgroup of $G_{(V)}$}. 

 \noindent In particular, arithmetic subgroups of $PGL_2(\R)$ are
 subgroups which are commensurable with the unit groups
 $PSL_1(\mathcal{B})$ of quaternion algebras over {\it totally real}
 number fields $E$, which are {\it split at exactly one place at
 infinity $\infty_0$} (and one can take $V=\{\infty_0\}$).
\end{para}

\begin{para}\label{para::5.4.2} The complete list of cocompact
 arithmetic triangle groups (there are $76$ such groups, up to
 permutation of $e_0,e_1,e_\infty$), and of the corresponding quaternion
 algebras, has been established by K. Takeuchi; moreover, he has shown
 that two arithmetic triangle groups are commensurable with each other
 if and only if they come from the same quaternion algebra. Here is the
 list \cite[Table 1]{ccoatg}:

\renewcommand{\arraystretch}{1.5}
\begin{longtable}{lcp{0.63\textwidth}}
 $E$ & $\mathrm{disc}(B)$ & $ (e_0,e_1,e_\infty)\text{ up to permutation}$ \\
 \hline
 $\Q$ & $2.3$ & $(2,4,6)\, (2,6,6)\, (3,4,4)\, (3,6,6)$ \\
 $\Q(\sqrt 2)$ & $v_2$ &
 $(2,3,8)\, (2,4,8)\, (2,6,8)\, (2,8,8)\, (3,3,4)\, (3,8,8)\,$ 
 $(4,4,4)\, (4,6,6)\, (4,8,8)$ \\
 $\Q(\sqrt 3)$ & $v_2$ &
 $(2,3,12)\, (2,6,12)\, (3,3,6)\, (3,4,12)\, (3,12,12)\, (6,6,6)\,$ \\
 $\Q(\sqrt 3)$ & $v_3$ & $(2,4,12)\, (2,12,12)\, (4,4,6)\, (6,12,12)\,$ \\
 $\Q(\sqrt 5)$ & $v_2$ & $(2,4,5)\, (2,4,10)\, (2,5,5)\, (2,10,10)\,
 (4,4,5)\, (5,10,10)$ \\
 $\Q(\sqrt 5)$ & $v_3$ & $(2,5,6)\, (3,5,5)$ \\
 $\Q(\sqrt 5)$ & $v_5$ & $(2,3,10)\, (2,5,10)\, (3,3,5)\, (5,5,5)$ \\
 $\Q(\sqrt 6)$ & $v_2$ & $(3,4,6)$ \\
 $\Q(\cos \pi/7)$ & $1$ &
 $(2,3,7)\, (2,3,14)\, (2,4,7)\, (2,7,7)\, (2,7,14)\, (3,3,7)\,$ 
 $(7,7,7)$ \\ 
 $\Q(\cos \pi/9)$ & $1$ & $(2,3,9)\, (2,3,18)\, (2,9,18)\, (3,3,9)\,
 (3,6,18)\, (9,9,9)$ \\
 $\Q(\cos \pi/9)$ & $v_2.v_3$ &
 $(2,4,18)\, (2,18,18)\, (4,4,9)\, (9,18,18)$ \\
 $\Q(\cos \pi/8)$ & $v_2$ &
 $(2,3,16)\, (2,8,16)\, (3,3,8)\, (4,16,16)\, (8,8,8)$ \\
 $\Q(\cos \pi/10)$ & $v_2$ & $(2,5,20)\, (5,5,10)$ \\
 $\Q(\cos \pi/12)$ & $v_2$ & 
 $(2,3,24)\, (2,12,24)\, (3,3,12)\, (3,8,24)\, (6,24,24)\,$ $(12,12,12)$ \\
 $\Q(\cos \pi/15)$ & $v_3$ & $(2,5,30)\, (5,5,15)$ \\
 $\Q(\cos \pi/15)$ & $v_5$ &
 $(2,3,30)\, (2,15,30)\, (3,3,15)\, (3,10,30)\, (15,15,15)$ \\
 $\Q(\sqrt 2, \sqrt 5)$ & $v_2$ & $(2,5,8)\, (4,5,5)$ \\
 $\Q(\cos \pi/11)$ & $1$ & $(2,3,11).$\\
% \caption{}
\end{longtable}

\medskip\noindent ($v_p$ denotes the unique place of $E$ above
 $p$. Notice that all fields 
 $E$ appearing in this table are Galois --- even abelian --- over $\Q$).
\end{para}

\begin{para}\label{para::5.4.3} We do not know whether all arithmetic
 triangle groups are 
 congruent subgroups (actually, unit groups of quaternion algebras seem
 to be unsolved hard cases of the congruence subgroup problem, \cf 
 \cite{rapinchuk92:_congr}). Nevertheless, in each commensurability
 class of arithmetic 
 triangle groups, there is a (unique) triangle group of the form
 \begin{equation*} \Gamma^\ast= \{g\in B^\times\mid 
 g\mathcal{B}=\mathcal{B}g, \Nr(g)>0\}: \end{equation*} this is the one
 which appears first in the each of the above rows, \cf
 \cite[Table 3]{ccoatg}. The corresponding Schwarz orbifolds are Shimura
 orbifolds $\mathcal{X}^\ast$. More generally, any Schwarz orbifold whose
 orbifold fundamental group is a congruence group, is a Shimura orbifold.
\end{para}


\subsection{Arithmetic $p$-adic triangle groups. A $p$-adic analogue of
Takeuchi's list.}\label{sub::5.5} 

\begin{para}\label{para::5.5.1} Arithmetic $p$-adic triangle groups are
 hyperbolic $p$-adic triangle groups which are {\it commensurable with
 the $p$-unit groups of quaternion algebras $\ovl{B}$ over number
 fields}.\footnote{In fact, these were the first $p$-adic triangle groups
 to be constructed \cite{poaptg}, and the existence of non-arithmetic
 hyperbolic $p$-adic triangle groups remained unclear for a while, before
 Kato's construction.}
 
 \medskip Let now $V$ be the set of $p$-adic and infinite places of
 $E$. It follows from the discussion of \ref{para::5.4.1} that
 arithmetic subgroups of 
 $G_{(V)}= PGL_2(E_v)$ are subgroups which are commensurable with the
 $p$-unit groups $PSL_1(\ovl{\mathcal{B}}[\frac{1}{p}])$ of quaternion
 algebras over {\it totally real} number fields $E$, which are {\it
 totally definite at infinity and split at exactly one place $v$ above
 $p$ }.
\end{para}

\begin{para}\label{para::5.5.2} Let $\Delta_p(e_0,e_1,e_\infty)\subset
 PGL_2(\C_p)$ be an arithmetic $p$-adic triangle group. Thus, there is a
 totally real number field $E \subset \C_p$, and a totally definite
 quaternion algebra $\ovl{B}$ over $E$, split at one place $v$ above $p$,
 and a maximal order $\ovl{\mathcal{B}}$, such that
 $\Delta_p(e_0,e_1,e_\infty)\subset PGL_2(E_v)$ is commensurable, up to
 conjugation, to
 $PSL_1\bigl(\ovl{\mathcal{B}}[\frac{1}{p}]\bigr)$. Instead of the 
 latter, it is equivalent and more convenient to work with the image
 $\Gamma_p^\ast$ in $PGL_2(E_v)$ of the normalizer $\bigl\{g\in
 \ovl{B}^\times\mid 
 \;g\ovl{\mathcal{B}}[\frac{1}{p}]=\ovl{\mathcal{B}}[\frac{1}{p}]g\bigr\}$.

 Let us show that the corresponding complex triangle group
 $\Delta_p(e_0,e_1,e_\infty)$ $\subset$ $PGL_2(\C)$ is arithmetic, attached
 to the same totally real number field $E$ (embedded into $\C$ via some
 real place $\infty_0$), and to the quaternion algebra $B$ over $E$
 obtained from $\ovl{B}$ by interchanging the local invariants at
 $\infty_0$ and $v$ (in particular $B$ is definite at every place at
 infinity except $\infty_0$, and {\it ramified at $v$} --- which turns out
 to be the unique place $v_p$ of $E$ above $p$ according to Takeuchi's
 classification).
 
 In the case $\Delta_p(e_0,e_1,e_\infty)=\Gamma_p^\ast$, this follows
 from \v{C}erednik's uniformization explained in  II.\ref{sub:7.4},
 II.\ref{sub:7.5}: in this 
 case, the Schwarz orbifold is a Shimura orbifold $\mathcal{X}^\ast$ and
 $\Delta(e_0,e_1,e_\infty)=\Gamma^\ast$.

 In general, the Schwarz
 orbifold $\mathcal{S}(e_0,e_1,e_\infty)$ is close to be a Shimura
 orbifold.  If fact, according to \ref{thm::5.3.7}, the $p$-adic Schwarz
 orbifold 
 $\mathcal{S}(e_0,e_1,e_\infty)$ is a finite orbifold-covering of a
 Mumford-Schwarz orbifold $\mathcal{S}(e'_0,e'_1,e'_\infty)$ with
 $\Delta_p(e_0,e_1,e_\infty)=\Delta_p(e'_0,e'_1,e'_\infty)$. Let $\Gamma$
 be a torsion-free subgroup of finite index of
 $\Delta_p(e'_0,e'_1,e'_\infty)\cap \Gamma_p^\ast$. Let $\Omega \subset
 \P^1_{\C_p}$ be the set of ordinary points for $\Gamma$. Then the
 Mumford curve $S=\Omega/\Gamma$ is simultaneously a finite orbifold
 covering of $\mathcal{S}(e'_0,e'_1,e'_\infty)$ and of $\mathcal{X}^\ast$
 (via the embeddings $\Gamma \subset \Delta_p(e'_0,e'_1,e'_\infty)$ and
 $\Gamma \subset \Gamma_p^\ast$ respectively).

 All these finite orbifold-coverings are defined over number
 fields $\subset \C$. Looking at the corresponding complex orbifolds, it
 is now clear that $\pi_1^{\mathrm{orb}}(\mathcal{S}(e_0,e_1,e_\infty))$
 $=\Delta(e_0,e_1,e_\infty)$ is commensurable with
 $\pi_1^{\mathrm{orb}}(\mathcal{X}^\ast)= \Gamma ^\ast$, whence the claim.
\end{para}

\begin{para}\label{para::5.5.3} Let us show, conversely, that if
 $\Delta(e_0,e_1,e_\infty)$ appears in Takeuchi's list, and if $B/E$ is
ramified at the place $v\mid p$, then $\Delta_p(e_0,e_1,e_\infty)$
exists and is an arithmetic $p$-adic triangle group; more precisely, it
is commensurable with the $p$-unit group of the quaternion algebras
$\ovl{B}/E$ obtained from $B$ by interchanging the local invariants at
$\infty_0$ (the indefinite place at infinity) and $v_p$ (in particular
$\ovl{B}$ is totally definite, and split at $v_p$).
 
 In the case $\Delta(e_0,e_1,e_\infty)=\Gamma^\ast$, \ie if
 $(e_0,e_1,e_\infty)$ is the first triple occurring on a row of 
 Takeuchi's table, this follows from \ref{thm::4.7.4}.     

 In general, $\Delta(e_0,e_1,e_\infty)$ is commensurable to an arithmetic
 triangle group $\Delta(e'_0,e'_1,e'_\infty)$ of the form $\Gamma^\ast$,
 and it is not difficult to deduce that $\Delta_p(e_0,e_1,e_\infty)$
 exists and is commensurable to
 $\Delta_p(e'_0,e'_1,e'_\infty)=\Gamma_p^\ast$. The point is that the
 corresponding respective projective connections $HGPC(a,b,c)$ and
 $HGPC(a',b',c')$ become isomorphic after suitable finite pull-backs.
\end{para}

\begin{para}\label{para::5.5.4} Putting \ref{para::5.5.2} and
 \ref{para::5.5.3} together, we conclude 
 that the arithmetic $p$-adic triangle groups
 $\Delta_p(e_0,e_1,e_\infty)$ are exactly those corresponding to the
 triples $(e_0,e_1,e_\infty)$ in Takeuchi's list when $v_p$ divides
 $\mathrm{disc}(B)$.  From that list, one gets:
\end{para}

\begin{thm}\label{thm::5.5.5} \cite{poaptg}. Up to permutation of the
 indices, there are $45$ arithmetic $2$-adic triangle groups, 
 $16$ arithmetic
 $3$-adic triangle groups, $9$ arithmetic $5$-adic triangle groups (and
 no arithmetic $p$-adic triangle group for $p>5$).

 More precisely, they are given by the following tables:

\begin{center}
 \renewcommand{\arraystretch}{1.5}
 \begin{longtable}{lcp{0.66\textwidth}}
 \multicolumn{3}{l}{\it Arithmetic $2$-adic triangle groups:}\\ \hline
 $E$ &$\mathrm{disc}(\ovl{B})$& $(e_0,e_1,e_\infty)
 \text{ \rm up to permutation}$ \\ \hline
 $\Q$ & $3$ & $(2,4,6)\,(2,6,6)\,(3,4,4)\,(3,6,6)$ \\
 $\Q(\sqrt 2)$ & $1$ &
 $(2,3,8)\, (2,4,8)\, (2,6,8)\, (2,8,8)\, (3,3,4)$\, $(3,8,8)\, 
 (4,4,4)\, (4,6,6)\, (4,8,8)$ \\
 $\Q(\sqrt 3)$ & $1$ &
 $(2,3,12)\, (2,6,12)\, (3,3,6)\, (3,4,12)\, (3,12,12)$\, $(6,6,6)$ \\
 $\Q(\sqrt 5)$ & $1$ &
 $(2,4,5)\, (2,4,10)\, (2,5,5)\, (2,10,10)\, (4,4,5)$\, $(5,10,10)$ \\
 $\Q(\sqrt 6)$ & $1$ & $(3,4,6)$ \\
 $\Q(\cos \pi/9)$ & $v_3$ & $(2,4,18)\, (2,18,18)\, (4,4,9)\, (9,18,18)$ \\
 $\Q(\cos \pi/8)$ & $1$ & $(2,3,16)\, (2,8,16)\, (3,3,8)\, (4,16,16)\,
 (8,8,8)$ \\ 
 $\Q(\cos \pi/10)$ & $1$ & $(2,5,20)\, (5,5,10)$  \\
 $\Q(\cos \pi/12)$ & $1$ &
 $(2,3,24)\, (2,12,24)\, (3,3,12)\, (3,8,24)\, (6,24,24)$\, $(12,12,12)$ \\
 $\Q(\sqrt 2, \sqrt 5)$ & $1$ & $(2,5,8)\, (4,5,5)$  \\ &&\\
 \multicolumn{3}{l}{\it Arithmetic $3$-adic triangle groups:} \\ \hline
 $\Q$ & $2$ & $(2,4,6)\, (2,6,6)\, (3,4,4)\, (3,6,6)$  \\
 $\Q(\sqrt 3)$ & $1$ & $(2,4,12)\, (2,12,12)\, (4,4,6)\, (6,12,12)$ \\
 $\Q(\sqrt 5)$ & $1$ & $(2,5,6)\, (3,5,5)$\\
 $\Q(\cos \pi/9)$ & $v_2$ & $(2,4,18)\, (2,18,18)\, (4,4,9)\, (9,18,18)$ \\
 $\Q(\cos \pi/15)$ & $1$ & $(2,5,30)\, (5,5,15)$ \\ && \\
 \multicolumn{3}{l}{\it Arithmetic $5$-adic triangle groups:}\\ \hline
 $\Q(\sqrt 5)$ & $1$ & $(2,3,10)\, (2,5,10)\, (3,3,5)\, (5,5,5)$ \\
 $\Q(\cos \pi/15)$ & $1$ &
 $(2,3,30)\, (2,15,30)\, (3,3,15)\, (3,10,30)\, (15,15,15)$ \\
 \end{longtable} 
\end{center}
\end{thm}

\begin{rems}
 (a) In these lists, one notice that $p$ divides
 at least one of the numbers $e_0,e_1,e_\infty$. This means that the
 parameters $a,b,c$ of the corresponding hypergeometric differential
 equation are not $p$-adic integers. In the literature on $p$-adic
 hypergeometric functions, on the contrary, $a,b,c$ are always assumed to
 be $p$-adic integers, for obvious convergence reasons.

 It is well-known that a fuchsian $p$-adic differential equations with
 Frobenius structure has exponents in $\Q\cap \Z_p$. Since this is not
 the case here, there is no Frobenius structure in the standard
 sense. However, there is some hidden Frobenius structure due to the
 ``geometric origin'' of the differential equation. We leave to the
 reader the intriguing question of finding out the right formulation of
 this hidden Frobenius structure (hint: the Frobenius structure is
 explicit at the level of the period mapping, \cf \cite{tpuosc}).

 \medskip
 (b) Takeuchi has also given a table with the various
 inclusions of triangle groups in each commensurability class. When
 $\Delta(e_0,e_1,e_\infty)$ is maximal (and $p$ ramified in the
 corresponding quaternion algebra), the corresponding $p$-adic triangle
 group $\Delta_p(e_0,e_1,e_\infty)$ is ``of Mumford type'', according to
 \ref{cor::5.3.7}.~(b).

 Arithmetic $p$-adic triangle groups not ``of Mumford type'' do exist,
 and are related to Shimura curves with {\it level divisible by $p$}.
 One can construct such $p$-adic triangle groups using Drinfeld's finite
 \'etale coverings $\Upsilon^n$ of the Drinfeld ``half space'' $\Omega$,
 which uniformize $p$-adically Shimura curves with level $p^n.M\,$ ($p$
 prime to $M$), \cf II.\ref{para:7.4.7}.  One has
 $\Gamma_p^+(p^nM)=\Gamma_p^+( M)$, 
 and the orbifold-quotient of a suitable intermediate covering $\Sigma$
 between $\Omega$ and $\Upsilon^n$ by $\Gamma_p^+( M)$ may be a Schwarz
 orbifold but not a Mumford-Schwarz orbifold. It seems that
 $\Delta_2(4,4,4)$ is of this type, for $p=2, \,n=M=1$ (we have not
 checked the details); the corresponding $\Gamma_2^+(1)$ is
 $\Delta_2(3,3,4)$.

 \medskip
 (c) Let $\ovl{S}$ be a smooth projective curve defined
 over $\ovl{\Q}$, and $\zeta_i$ be $\ovl{\Q}$-points on $\ovl{S}$
 endowed with indices $e_i$. We then say that the orbifold
 ${\mathcal{S}}=(\ovl{S}, (\zeta_i,e_i))$ is defined over $\ovl{\Q}$. In
 general, there is no reason for the uniformizing differential equation
 to be defined over $\ovl{\Q}$.

 However, this occurs if either ${\mathcal{S}}$ is a Schwarz orbifold
 (the uniformizing differential equation being essentially of
 hypergeometric type with rational parameters), or a Shimura orbifold
 (the uniformizing differential equation being essentially of Gauss-Manin
 type, \cf \ref{thm::4.7.4}). On the other hand, it is not difficult to
 see that the 
 property for the uniformizing differential equation to be defined over
 $\ovl{\Q}$ depends only on the commensurability class of the orbifold
 fundamental group $\pi_1^{\mathrm{orb}}(\mathcal{S})$.

 A transcendence conjecture of Chudnovsky, modified by Krammer 
 \cite[12]{aeoaafg}, predicts that the Schwarz and the Shimura 
 cases are the only exceptions, up to commensurability. Namely: 
 
 \begin{coj}\label{coj::5.5.6} (Chudnovsky-Krammer). Let ${\mathcal{S} }$
  be a one-dimensional hyperbolic orbifold defined over
  $\ovl{\Q}$. Assume that the uniformizing differential equation is
  defined over $\ovl{\Q}$. Then $\pi_1^{\mathrm{orb}}(\mathcal{S})$ is either
  commensurable to a triangle group, or is arithmetic.
 \end{coj}

 We {\it conjecture} the (straightforward) $p$-adic analogue of
 \ref{coj::5.5.6}.
 
 \medskip
 (d) By the theory of $p$-adic Betti lattices, there is
 a canonical $\ovl{\Q}$-structure in the $\C_p$-space of solutions at
 any point of $\P^1(\ovl{\Q})$ of the hypergeometric differential
 equation $HDGE(a,b,c)$ attached to an arithmetic $p$-adic triangle group
 (\cf \ref{thm::4.7.4}. If the base point is chosen to be one of the
 singularities $0,1,\infty$, one applies \ref{thm::4.7.4}
 to a finite orbifold-covering of the Schwarz orbifold which is a
 Shimura curve attached to a torsion-free congruence subgroup). 
 
 This raises the following (open) questions:

 {\it Is there a $\ovl{\Q}$-structure in the $\C_p$-space
 of solutions of the hypergeometric differential equation $HDGE(a,b,c)$
 attached to a non-arithmetic $p$-adic triangle group? Are
 non-arithmetic $p$-adic triangle groups related to some $p$-adic period
 mappings and suitable higher-dimensional Shimura orbifolds?}
\end{rems} 


\subsection{On special values of hypergeometric functions attached to 
arithmetic $p$-adic triangle groups.}\label{sub::5.6}  

\begin{para}\label{para::5.6.1} Let $a,b,c$ be rational numbers such that
 $c>a+b$. Then the hypergeometric series ${}_2F_1(a,b;c;z)$ 
 converges at $1$ and takes the value (Gauss' formula) 
 \begin{equation*}
  {}_2F_1(a,b;c;1)= \frac{\Gamma(c)\Gamma(c-a-b)}{\Gamma(c-a)\Gamma(c-b)}\;.
 \end{equation*}   

\noindent If one considers ${}_2F_1(a,b;c;z)$ as a multivalued analytic
function on $\P^1\setminus\{0,1,\infty\}$, the 
values at $1$ form a finitely generated subgroup of $\C$; they generate
a $\ovl{\Q}$-subspace of $\C$ of dimension $\leq 2$,
namely:  
\begin{equation*}
 \frac{\Gamma(c)\Gamma(c-a-b)}{\Gamma(c-a)\Gamma(c-b)}\ovl{\Q} \, +\,
  \frac{\Gamma(c)\Gamma(a+b-c)}{\Gamma(a)\Gamma(b)}\ovl{\Q}. 
\end{equation*}
Similarly, its
values at $\infty$ generate the $\ovl{\Q}$-subspace of $\C$  
\begin{equation*}
 \frac{\Gamma(c)\Gamma(b-a)}{\Gamma(c-a)\Gamma(b)}\ovl{\Q} \, +\,
  \frac{\Gamma(c)\Gamma(a-b)}{\Gamma(a)\Gamma(c-b)}\ovl{\Q}. 
\end{equation*}   
\end{para}

\begin{para}\label{para::5.6.2} If $HGDE(a,b,c)$ is the differential equation
 attached as above to an {\it arithmetic} triangle group
$\Delta(e_0,e_1,e_\infty)$, these formulas are special cases of similar
formulas expressing the special values of ${}_2F_1(a,b;c;z)$ at $CM $
points, in terms of Gamma values, up to factors in $\ovl{\Q}$.

The point is that there is a $\ovl{\Q}$-structure in
the complex space of solutions which is stable under monodromy: the
 $\mathrm{H}_{1,\mathrm{B}}({\;\,},\ovl{\Q})$ of the fake elliptic curves $A_z$
 parametrized by $z\in\P^1$ ($\P^1$ being viewed as a Shimura curve) if
the quaternion algebra attached to $\Delta(e_0,e_1,e_\infty)$ is over
$\Q$, \cf II.\ref{para:7.6.3} for the general case. For any point $z\in
\P^1(\ovl{\Q})$, let $V_z$ be the $\ovl{\Q}$-structure in the complex
space of solutions at $z$ consisting of solutions whose Taylor
coefficients are in $\ovl{\Q}$. Then $V_z$ is related to
$\mathrm{H}_{1,\mathrm{B}}(A_{z,\C},\ovl{\Q})$ by the period matrix of
 $A_z$. At any CM 
point $z$ (according to an old observation of Shimura, the branch points
of a Shimura orbifold --- here, $0,1,\infty$ --- are actually CM points),
the Chowla-Selberg formula (I.\ref{thm-ogus-formula}) thus relates the two
$\ovl{\Q}$-structures in terms of $\Gamma$-values. In particular, $V_0$
is generated by ${}_2F_1(a,b,c;z)$ and
$z^{1-c}{}_2F_1(a+1-c,b+1-c,2-c;z)$, and the branch point $0$
corresponds to a CM point (for elliptic curves rather than fake elliptic
curves, this was worked out in I.\ref{coj-no-complex-multi}).
\end{para}

\begin{para}\label{para::5.6.3} There is an entirely analogous phenomenon for
 {\it arithmetic $p$-adic} triangle 
group $\Delta_p(e_0,e_1,e_\infty)$: thanks to II.\ref{thm:7.6.4}, there is a
 $\ovl{\Q}$-structure in 
the $\C_p$-space of solutions which is stable under $p$-adic monodromy:
 the ``$p$-adic Betti lattice''
 $\mathrm{H}_{1,\mathrm{B}}({\;\,},\ovl{\Q})$ of the fake 
 elliptic 
 curves $A_z$ if the quaternion algebra attached to
 $\Delta_p(e_0,e_1,e_\infty)$ is over $\Q$, \cf II.\ref{para:7.6.3} for the
 general case. ``Periods'' should be understood here in the sense of
 I.\ref{sub-p-adic-betti-lattice-supersing}; $p$-adic CM periods
 were computed in terms of $\Gamma_p$-values 
 in I.\ref{thm-ogus-formula}, I.\ref{para-ramif} (for $p\neq 2$), 
 I.\ref{coj-complex-mult}.

 Therefore, if $HGDE(a,b,c)$ is the differential equation attached as
 above to an {\it arithmetic $p$-adic} triangle group
 $\Delta_p(e_0,e_1,e_\infty)$, there is, at least for odd $p$, a
 expression in terms of $\Gamma_p$-values of the $\ovl{\Q}$-space of
 dimension $\leq 2$ generated by the values taken at any given CM point, \eg
 the point $1$, by ${}_2F_1(a,b,c;z)$ viewed as a multivalued $p$-adic
 analytic function on a suitable finite \'etale covering of
 $\P^1\setminus\{0,1,\infty\}$.

 We shall see examples in the next section.
\end{para}

\newpage

\section{The tale of $F(\frac{1}{24},\frac{7}{24},\frac{5}{6};z)$
(Escher's triangle group and its diadic and triadic twins).}\label{sec::6}
 
\markboth{\thechapter. $p$-ADIC ORBIFOLDS AND MONODROMY.}%
{\thesection. THE TALE OF $F(\frac{1}{24},\frac{7}{24},\frac{5}{6};z)$.}

\begin{abst}
 We study the rich geometry related to the simplest cocompact arithmetic
 triangle group: associated Shimura curves, period mapping and monodromy
 --- in the archimedean, diadic and triadic cases. At other primes, the
 story would be similar to that of $F(\frac{1}{2},\frac{1}{2},{1};z)$;
 we do not repeat it.
\end{abst}

\subsection{Escher's triangle group $\Delta(2,4,6)$.}\label{sub::6.1}

\begin{para}\label{para::6.1.1} This is the symmetry group of the
 black-and-white tesselation of the unit disk by hyperbolic 
 triangles with angles $(\frac{\pi}{2}, \frac{\pi}{4}, \frac{\pi}{6})$

%%% PICTURE 6a 
\begin{figure}[h]
 \begin{picture}(300,170)(0,0)
  \put(0,0){\includegraphics[scale=0.43,clip]{fig3-6a1.eps}}
  \put(170,0){\includegraphics[scale=0.6,clip]{fig3-6a2.eps}}
 \end{picture}
 \caption{}
\end{figure}

 M.C. Escher has engraved several celebrated variants of this Schwarz
 tesselation (``cirkel-limiet'', angels and 
 demons), {\it cf}  \cite{escher71:_m} \footnote{he has also used another
 Schwartz tesselation in a picture --- ``cirkel-limiet III'' ---
 brilliantly analyzed by H.S.M. Coxeter \cite{coxeter79:_escher_iii}}.

\medskip By gluing together two basic triangles, we get hyperbolic
 triangles with angles $(\frac{\pi}{3}, \frac{\pi}{4}, 
 \frac{\pi}{4})$ or $(\frac{\pi}{2}, \frac{\pi}{6}, \frac{\pi}{6})$

%%% PICTURE 6b 
\begin{figure}[h]
 \begin{picture}(300,100)(0,-50)
  \put(0,-32){\includegraphics[scale=0.8,clip]{fig3-6b.eps}}
  {\small
  \put(0,-10){\makebox(0,0)[c]{$O$}}
  \put(47,-10){\makebox(0,0)[c]{$Q$}}
  \put(65,34){\makebox(0,0)[c]{$P$}}
  \put(-6,0){\makebox(0,0)[b]{$6$}}
  \put(55,0){\makebox(0,0)[b]{$2$}}
  \put(54,37){\makebox(0,0)[c]{$4$}}

  \put(85,0){\makebox(0,0)[b]{$3$}}
  \put(145,37){\makebox(0,0)[c]{$4$}}
  \put(152,-30){\makebox(0,0)[c]{$4$}}

  \put(237,37){\makebox(0,0)[c]{$2$}}
  \put(174,0){\makebox(0,0)[b]{$6$}}
  \put(269,0){\makebox(0,0)[b]{$6$}}
  }
 \end{picture}
 \caption{}
\end{figure}

\noindent This exhibits $\Delta(3,4,4)$ and $\Delta(2,6,6)$ as subgroups
 of index $2$ of $\Delta(2,4,6)$. 

\noindent Furthermore, by gluing two triangles with angles
 $(\frac{\pi}{2}, \frac{\pi}{6}, \frac{\pi}{6})$, one can get a 
triangle with angles $(\frac{\pi}{3}, \frac{\pi}{6}, \frac{\pi}{6})$,
 which exhibits $\Delta(3,6,6)$ as a subgroup 
 of index $2$ of $\Delta(2,6,6)$. 

\noindent Similarly, the intersection of $\Delta(2,6,6)$ and
 $\Delta(3,4,4)$ in $\Delta(2,4,6)$ is a quadrangle 
 group $\diamondsuit (2,2,3,3)$.
\end{para}

\begin{para}\label{para::6.1.2} The vertices of one of the basic triangles are

\begin{equation*}
 O=(0,0),\;Q=(\sqrt 2 -1, 0),\;P=\Bigl(\sqrt 3 \cos
 \frac{5\pi}{12}, \cos \frac{5\pi}{12}\Bigr).
\end{equation*}

\medskip\noindent
The matrix of the composition of the reflection along side $QP$, \resp
$PO$, followed by the reflection 
along side $OQ$ is $A=\sqrt {-1}\begin{pmatrix}\sqrt 2 & -1 \\1 &-\sqrt 2
\end{pmatrix}$, \resp  $B=\dfrac{1}{2}\begin{pmatrix}\sqrt 3 +\sqrt
{-1} & 0 \\0 & \sqrt 3 -\sqrt {-1}
\end{pmatrix}$. 

\noindent These matrices satisfy $A^2=B^6=(AB)^4=-1$ and, taken up to
 sign, generate the triangle group $\Delta(2,4,6)$, 
 (\cf \cite[II5]{ntatg}.)
\end{para}


\begin{para}\label{para::6.1.3} The group $\Delta(2,4,6)$ enjoys the following
 ``universal property'', discovered by Mennike in 1968: 

\begin{quote}
 {\it for all but a finite number of indefinite ternary quadratic forms
 $q$ with integral coefficients, the 
 (integral) associated unitary group of $q$ is contained in $\Delta(2,4,6)$.}
\end{quote}

 \noindent It coincides with $\Delta(2,4,6)$ in the case of
 $q=3x^2-y^2-z^2$, (\cf \cite[III3]{ntatg}).  These unitary groups
 have been thoroughly investigated by Fricke and Klein.
\end{para}

\subsection{Some Shimura curves attached to $D_{2.3}$.}\label{sub::6.2}

\begin{para}\label{para::6.2.1} The $\Q$-subalgebra of $M_2(\Q (\sqrt
 {-1}, \sqrt 2, \sqrt 3))$ generated by the matrices $A$ and $B$ 
 is the quaternion algebra $B=B_{2.3}$ over $\Q$ with discriminant
 $6$. This quaternion algebra is unique up to 
 isomorphism and admits standard generators $i,j$ and relations 
 $i^2=-1,\;j^2=3,\;ij=-ji$.  

 \noindent It follows that $\Delta(2,4,6)$ is a subgroup of
 $B_{2.3}^\times/\Q^\times$. According to \cite{ccoatg}, 
 $\Delta(2,4,6)$, $\Delta(3,4,4)$, $\Delta(2,6,6)$ and $\Delta(3,6,6)$
 are the only cocompact triangle 
 groups which give rise to a quaternion algebra over $\Q$. 

 \noindent As any indefinite quaternion algebra over $\Q$, $B_{2.3}$ has
 a unique maximal order $\mathcal{B}=\mathcal{B}_{2.3}$ 
 modulo conjugation, \cf \cite[p.99]{adadq}. It can be described
 as $\Z  + \Z i + \Z j + \Z \rho$, where 
 $\rho =\frac{1+i+j+ij}{2}$. According to \cite{ccoatg}, it can also be
 described as {\it the subring of 
 $B_{2.3}$ generated by the elements of $\diamondsuit (2,2,3,3)$} (the
 subgroup of $\Delta(2,4,6)$ of index $4$ 
 considered above). 

 Following the notation of  II.\ref{sub:7.4}, we introduce the following
 arithmetic subgroups of $B^\times/\Q^\times$:

\begin{itemize}
 \item the image $\Gamma^+(N)$ of the group $\{g\in
       (1+N\mathcal{B})^\times \mid \Nr(g)=1\}$ in 
       $PSU(1,1)(\R)$ $=\Aut(\matheur{D}(0,1^-))$, where $\Nr$ stands for the
       reduced norm ($\Gamma^+(1)=PSL_1(\mathcal{B})$ is abbreviated
       into $\Gamma^+$),
 \item the image $\Gamma^\ast$ in $PSU(1,1)(\R)$ of the group 
       $N(\mathcal{B})^+=\{g\in B^\times\mid
       g\mathcal{B}=\mathcal{B}g \mid \Nr(g)>0\}$, \ie 
       $N(\mathcal{B})^+/\Q^\times$. 
\end{itemize}
       
\medskip According to \cite{ccoatg}, one has: 
 
\begin{itemize}
 \item $\Delta(2,4,6)=\Gamma^\ast$, 
 \item $\diamondsuit (2,2,3,3)=\Gamma^+$.
\end{itemize}

These groups also appear in \cite[p.122]{adadq}, together
with the orbifolds

\noindent ${\mathcal{X}^+}_\C^{\mathrm{an}}= \matheur{D}(0,1^-)/\diamondsuit 
 (2,2,3,3)$ and ${\mathcal{X}^\ast}_\C^{\mathrm{an}}=\matheur{D}(0,1^-)/\Delta(2,4,6)$.
 The orbifold Euler characteristics are found to be $-\frac{1}{3}$ and 
 $-\frac{1}{12}$ respectively.  
\end{para}

\begin{para}\label{para::6.2.2} Any element of $N(\mathcal{B})^+$ can be
 written uniquely as a product 
\begin{equation*}
 q.g.(1+i)^{\epsilon_2}(3i+j+ij)^{\epsilon_3},
\end{equation*}
where $q\in \Q, \; g\in SL_1(\mathcal{B}),\;\epsilon_m=0 $
or $1$, \cf \cite{michon84:_courb_shimur}. The quotient group 

 \begin{align*}
  &\Delta(2,4,6)/\diamondsuit(2,2,3,3)\\
  = &\{1,w_2=[1+i],\;w_3=[3i+j+ij],\; w_6=w_2w_3=[-3+3i+2ij]\}\\
  \cong &(\Z/2\Z)^2
 \end{align*}
 is called the Atkin-Lehner group. The fixpoints of the involutions $w_m$
 have been the object of thorough 
 investigations, \cf {\it loc.~cit.} and \cite{shimura67:_const} (\cf also
 II.\ref{para:7.4.4}). In our case, the results are:

 $\bullet$ the two fixpoints of $w_2$ on $\matheur{D}(0,1^-)/\diamondsuit
 (2,2,3,3)$ have the same image in $\matheur{D}(0,1^-)/\Delta(2,4,6)$, which is the
 image of the fixpoint in $\matheur{D}(0,1^-)$ of any maximal 
 embedding of $\Z[\sqrt {-1}]$ in $\mathcal{B}$, 

$\bullet$ the two fixpoints of $w_3$ on $\matheur{D}(0,1^-)/\diamondsuit
(2,2,3,3)$ have the same image in $\matheur{D}(0,1^-)/\Delta(2,4,6)$, which is the
image of the fixpoint in $\matheur{D}(0,1^-)$ of any maximal 
embedding of $\Z[^3\sqrt{1}]$ in $\mathcal{B}$, 

$\bullet$ the two fixpoints of $w_2$ on $\matheur{D}(0,1^-)/\diamondsuit
(2,2,3,3)$ have the same image in $\matheur{D}(0,1^-)/\Delta(2,4,6)$, which is the
image of the fixpoint in $\matheur{D}(0,1^-)$ of any maximal 
embedding of $\Z[\sqrt {-6}]$ in $\mathcal{B}$, 

 From easy considerations of ramification, one draws that 

\begin{itemize}
 \item $\an{\Gamma^+,(1+i)}=\Delta (3,4,4)$
 \item $\an{\Gamma^+,(3i+j+ij)}=\Delta (2,6,6)$
 \item $\an{\Gamma^+,(1+i)}$ is the third intermediate group, which is
       a quadrangle group $\diamondsuit 
       (2,2,2,3)$ (not to be confused with $\diamondsuit (2,2,3,3)$).
\end{itemize}
\end{para}

\begin{lem}\label{lem::6.2.3} The congruence group $\Gamma^+(2)$ is
 torsion-free. The quotient $\Gamma^+/\Gamma^+(2)$ is the 
 tetrahedral group $\Delta(2,3,3)\cong \mathfrak{A}_4$ (of order $12$).
\end{lem}

\begin{proof}
 Recall that $\mathcal{B}=\Z + \Z i + \Z j + \Z \rho$. Any element
 $a+bi+cj+dij$ of finite order in the congruence subgroup
 $SL_1(\mathcal{B})(2)$ of level two must have reduced trace $2a=\pm 2$
 and reduced norm $a^2+b^2-3c^2-3d^2=1$; it is easy to see that the only
 such elements are $\pm 1$, which implies the first assertion.

 The quotient group $\Gamma^+/\Gamma^+(2)$ is represented by $\{1,i,j,ij,
 \frac{1\pm i\pm j \pm ij }{2}\}$ modulo $2$. It admits $\{1,\ovl{i},
 \ovl{j},\ovl{i}\ovl{j}\}$ as a normal subgroup, and
 $\{1,\ovl{\rho},\ovl{\rho}^2\}$ as a non-normal subgroup. its center is
 trivial.  It is thus isomorphic to the tetrahedral group, the group of
 oriented symmetries of the tetrahedron.
%%% PICTURE 6c
\begin{figure}[h]
 \begin{picture}(100,100)(0,0)
  \put(0,0){\includegraphics[scale=0.3,clip]{fig3-6c.eps}}
 \end{picture}
 \caption{}
\end{figure}
\end{proof}

\begin{cor}\label{cor::6.2.4}
 The associated Shimura curve $\mathcal{X}^+(2)$ has genus $g=3$. 
\end{cor}
\begin{proof}
 By definition,
 $\mathcal{X}^+(2)_\C^{\mathrm{an}}=\matheur{D}(0,1^-)/\Gamma^+(2)$ (here we
 prefer to identify the universal 
 covering with the unit disk rather than the Poincar\'e upper half plane
 as in  II.\ref{sub:7.4}). The assertion follows from 
 the computation
 $2-2g=\chi(\matheur{D}(0,1^-)/\Gamma^+(2))
 =12.\chi_{orb}(\matheur{D}(0,1^-)/\Gamma^+)=-4$.
\end{proof}

\begin{lem}\label{lem::6.2.5} The congruence group $\Gamma^+(2)$ is a
normal subgroup of $\Delta(2,4,6)$ $=$ $\Gamma^\ast$.  The quotient 
$\Gamma^\ast/\Gamma^+(2)$ is the extended
octahedral group $\Delta^\ast(2,3,4)$ $\cong$ $\mathfrak{S}_4\times \Z/2\Z$ 
(of order $48$). The generator of the factor $\Z/2\Z$ is 
given by the image $\tilde w_3$ of $(3i+j+ij)$.
\end{lem}

\par We refer to \cite[N. 58]{ford29:_autom} and
 \cite[p.27]{coxeter74:_regul} for a discussion 
 of $\Delta^\ast(2,3,4)$, the group of all symmetries of the octahedron.

%%% PICTURE 6d
\begin{figure}[h]
 \begin{picture}(100,100)(0,0)
  \put(0,0){\includegraphics[scale=0.3,clip]{fig3-6d.eps}}
 \end{picture}
 \caption{}
\end{figure}

\begin{proof}
 Let us first show that $\Gamma^+(2)$ is a normal
 subgroup of $\Delta(2,6,6)=\an{\Gamma^+,(3i+j+ij)}$. It makes 
 sense to reduce the latter modulo $2$, $(3i+j+ij)$ being sent to $1 \in
 \Delta(2,3,3)$; on the other hand, $\Nr$ induces a 
 surjective homomorphism $\Delta(2,6,6)\to 3^\Z/9^\Z\cong \Z/2\Z$, and
 $(3i+j+ij)$ is sent to a generator of $\Z/2\Z$. It follows that
 $\Gamma^+(2)$ is a normal subgroup of $\Delta(2,6,6)$, with quotient
 group $\mathfrak{A}_4\times \Z/2\Z$.

 \noindent We next remark that every element of $\Gamma^+(2)$ can be written 

\begin{align*}
 &a+bi+cj+dij,\text{ with }a,b,c,d \in \Z,\;b \equiv c\equiv d
 \not\equiv a \;(2),\\
 &a^2+b^2-3c^2-3d^2=1.
\end{align*} 
 Since $(1+i)(a+bi+cj+dij)(1+i)^{-1}=a+bi-dj+cij,$ we see that $(1+i)$
 normalizes  $\Gamma^+(2)$, and conclude that $\Gamma^+(2)$ is normal in
 $\Gamma^\ast$. The quotient group sits in an extension 
 \begin{equation*}
  1\to \Delta(2,3,3)\times \Z/2\Z\to G \to \Z/2\Z\to 1,
 \end{equation*}
 the quotient $\Z/2\Z$ being generated by $(1+i)$. Since the image of
 $(1+i)$ in $G$ is not central, we conclude that $G\cong
 \mathfrak{S}_4\times \Z/2\Z$.
\end{proof} 

\begin{cor}\label{cor::6.2.6} The automorphism group of the curve
 $\mathcal{X}^+(2)_\C$ is $\Delta^\ast(2,3,4)$.
\end{cor}

\begin{proof}
 It follows from \ref{cor::6.2.4} and \ref{lem::6.2.5} that $\Aut
 \mathcal{X}^+(2)_\C$ contains $\Delta^\ast(2,3,4)$. On the other hand,
 any automorphism of $\mathcal{X}^+(2)_\C$ lifts to an element of
 $SU(1,1)(\R)$ which normalizes $\mathcal{B}$. The image in
 $PSU(1,1)(\R)$ of such an element lies in $\Gamma^\ast$, whence
 $\Aut\;\mathcal{X}^+(2)_\C\subset
 \Gamma^\ast/\Gamma^+(2)=\Delta^\ast(2,3,4)$.
\end{proof}

\medskip By similar arguments, one shows that the groups
$\Delta(3,4,4)/\Gamma^+(2)$ and
\noindent $\diamondsuit(2,2,2,3)/\Gamma^+(2)$ are the octahedral group
$\Delta(2,3,4)\cong \mathfrak{S}_4$.

\newpage
One reads from the following table of ramification that $\tilde w_3$ is
a hyperelliptic involution of $\mathcal{X}^+(2)_\C$.

%%% PICTURE 6e 
\begin{figure}[h]
 \begin{picture}(300,480)(-10,-20)
  \put(5,-14){\includegraphics[scale=0.8,clip]{fig3-6e.eps}}
  \put(0,-20){\makebox(0,0)[c]{$\chi_{\mathrm{orb}}$}}
  \put(0,0){\makebox(0,0)[r]{$-\frac{1}{12}$}}
  \put(0,67){\makebox(0,0)[r]{$-\frac{1}{6}$}}
  \put(0,158){\makebox(0,0)[r]{$-\frac{1}{3}$}}
  \put(0,203){\makebox(0,0)[r]{$-\frac{1}{2}$}}
  \put(0,271){\makebox(0,0)[r]{{\footnotesize $-1$}}}
  \put(0,338){\makebox(0,0)[r]{{\footnotesize $-2$}}}
  \put(0,407){\makebox(0,0)[r]{{\footnotesize $-4$}}}
 \end{picture}
 \caption{}
\end{figure}
\vfill

\newpage
\begin{para}\label{para::6.2.7} Similar considerations apply to the congruence
 subgroup of level three $\Gamma^+(3)$. The Shimura curve 
$\mathcal{X}^+(3)_\C$ has genus $4$ and its automorphism group is of
 order $72$. We omit the details. 
\end{para}

\begin{para}\label{para::6.2.8} By the general theory of Shimura curves,
 $\mathcal{X}^+(2)$ is a moduli space for polarized abelian surfaces
 with quaternionic multiplication by $\mathcal{B}_{2.3}$ and level-two
 structure, \cf  II.\ref{sub:7.4}. It is defined over $\Q$, as well as
 $\mathcal{X}^\ast$ and the intermediate Shimura curves attached to
 $\Delta(3,4,4),\;\Delta(2,6,6),\;\diamondsuit(2,2,2,3)$ and
 $\diamondsuit(2,2,3,3)$.
 According to Y. Ihara, A. Kurihara \cite{oseoedscatmu}, A. Ogg and
 J.F. Michon, in each of the cases 
 $\Delta(2,4,6),\;\Delta(3,4,4),\;\Delta(2,6,6),\;\diamondsuit(2,2,2,3)$,
 the $\Q$-curve itself is $\P^1_\Q$, whereas for
 $\diamondsuit(2,2,3,3)$, it is the conic $x^2+y^2+3z^2=0$.\footnote{on
 the other hand, we do not know any hyperelliptic equation for
 $\mathcal{X}^+(2)$, even over $\C$.} The Atkin-Lehner involutions are
 defined over $\Q$.
\end{para}

\begin{para}\label{para::6.2.9} Any CM point on $\mathcal{X}^+(2)$
 parametrizes an abelian surface isogenous to the square of an elliptic 
 curve with complex multiplication. It follows from \ref{para::6.2.2}
 that the points 
 of ramification of the ``projection'' $\mathcal{X}^+(2)\to
 \mathcal{X}^\ast$ are CM points. More precisely, the corresponding
 imaginary quadratic fields are  

\begin{itemize}
 \item $\Q[\sqrt {-1}]$ for the points of ramification index $4$ (\eg 
       the image of the point $P\in \matheur{D}(0,1^-)$, \cf \ref{para::6.1.2})
 \item $\Q[ \sqrt{-3}]$ for the points of ramification index $6$ (\eg 
       the image of the point $O\in \matheur{D}(0,1^-)$) 
 \item $\Q[\sqrt {-6}]$ for the points of ramification index $2$ (\eg 
       the image of the point $Q\in \matheur{D}(0,1^-)$). 
\end{itemize}
\end{para}


\subsection{Triadic uniformization.}\label{sub::6.3}

\begin{para}\label{para::6.3.1} The Shimura curve $\mathcal{X}^+(2)$ has good
 reduction at any prime $p\neq 2,3$. As for the critical primes $2$ and
 $3$, the situation for $p=3$ is easier to describe than for $p=2$,
 since $3$ does not divide the level (II.\ref{sub:7.4}). 

 For the triadic \v{C}erednik-Drinfeld uniformization, the relevant
 quaternion algebra is the Hamilton definite quaternion algebra
 $\ovl{B}=B_{2.\infty}$ over $\Q$ with basis $(1,i,j,ij)$ and relations
 $i^2=j^2=-1, ij=-ji$. A maximal order (unique up to conjugation, {\it
 cf.} \cite[5.11]{adadq}) consists of the Hurwitz quaternions 

\begin{equation*}
\ovl{\mathcal{B}}=\ovl{\mathcal{B}}_{2.\infty} = \Z  + \Z i + \Z j + \Z
 \rho,\;\rho =\frac{1+i+j+ij}{2}. 
\end{equation*}
  
\noindent
The group $PGL_1(\ovl{\mathcal{B}})=\ovl{\mathcal{B}}^\times/\pm 1$ is
 finite, isomorphic to the tetrahedral group $\Delta(2,3,3)\cong
 \mathfrak{A}_4$.

\medskip 
Following the notation of  II.\ref{para:7.4.3}, we introduce the following
 $\{3,\infty\}$-arithmetic subgroups of $\ovl{B}^\times/\Q^\times$: 

\medskip $\bullet$ the image $\Gamma_3^+(N)$ of the group $\{g \in
 (1+N\ovl{\mathcal{B}}[\frac{1}{3}])^\times\mid \ord_3 \Nr(g) \;{\rm
 {even}}\}$ in $PGL_2(\Q_3)$ $=\Aut(\Omega_{\C_3})$  ($\Gamma_3^+(1)$
 is abbreviated into
 $\Gamma_3^+$),

 \medskip $\bullet$ the image $\Gamma_3^\ast$ in $PGL_2(\Q_3)$ of the group 
 
 \noindent $N(\mathcal{B} [\frac{1}{3}])=\{g\in \ovl{B}^\times \mid 
 g\ovl{\mathcal{B}}[\frac{1}{3}]=\ovl{\mathcal{B}}[\frac{1}{3}]g\}$,
 \ie $N(\mathcal{B}[\frac{1}{3}])/\Q^\times$.
 
\medskip According to \v{C}erednik-Drinfeld, one has:

\medskip
 $\bullet$
 ${\mathcal{X}^+(2)}_{\C_3}^{\mathrm{an}}=\Omega_{\C_3}/\Gamma_3^+(2)$ 
\medskip

$\bullet$ ${\mathcal{X}^+}_{\C_3}^{\mathrm{an}}=\Omega_{\C_3}/\Gamma_3^+$. 

\medskip

$\bullet$
 ${\mathcal{X}^\ast}_{\C_3}^{\mathrm{an}}=\Omega_{\C_3}/\Gamma_3^\ast$.

 In terms of triadic triangle (\resp quadrangle) groups, one can
 thus write 
 
\medskip
$\bullet$ $\Delta_3(2,4,6)=\Gamma_3^\ast$, 

$\bullet$ $\diamondsuit_3 (2,2,3,3)=\Gamma_3^+$.
\end{para}

\begin{para}\label{para::6.3.2} These uniformizations actually hold over
 $\Q_9$, the quadratic unramified extension of $\Q_3$. The corresponding
 Mumford curves over $\Q_9$ appear in \cite[IX]{sgamc}, without 
 referring to the \v{C}erednik theorem. The results of {\it loc.~cit.}
 allow to give a concrete description of the 
 arithmetic groups in terms of the generators $1,i,j,ij$ of $\ovl{B}$. 

 Let us fix a square root of $-1$ in $\Q_9$. A convenient matrix
 representation %\footnote{this is the same convention as in 
% \cite[6]{bertolini98:_heegn_l_cered}; Gerritzen and Van der Put prefer}
 for the elements of $\ovl{B}\otimes\Q_9$ is 

\begin{equation*}
a+bi+cj+dij \mapsto
 \begin{pmatrix}
  a+b\sqrt{-1} & c+d\sqrt{-1}
  \\-c+d\sqrt{-1} & a-b\sqrt{-1}  
 \end{pmatrix}
\end{equation*}  
 \noindent Of course, $\ovl{D}\otimes\Q_3\cong M_2(\Q_3)$, and any such
 matrix representation is conjugate to the previous after
 $\otimes\Q_9$ one by some $u\in PGL_2(\Q_9)$; $u$ is a kind of triadic
 Cayley transform which sends the Drinfeld space 
 $\Omega$ to the complement $u\Omega$ of the ``circle''
 $\{x+y\sqrt{-1}\mid x,y\in \Q_3, \;x^2+y^2=-1\}$ in the projective line  
 (here, we follow the usual convention that $PGL_2$ acts on the left on
 $\P^1$).
 
 After {\it loc.~cit.}, the congruence level subgroup of level two
 $PGL_1(\ovl{\mathcal{B}}[\frac{1}{3}])(2)$ is generated by 

 \begin{equation*}
  z_{\epsilon\epsilon'}=i+\epsilon j +\epsilon' ij,\;\;\epsilon=\pm 1,
   \epsilon' =\pm 1
 \end{equation*}

\noindent with the only relations $z_{\epsilon\epsilon'}^2=1$. The
 subgroup of index two
$\Gamma_3^+(2)$ consisting of words of even length (which is also the
 kernel of the homomorphism $PGL_1(\ovl{\mathcal{B}}[\frac{1}{3}])(2)\to
 3^\Z/9^\Z\cong \Z/2\Z$ induced by $\Nr$) is a {\it Schottky group of
 rank three} in $PGL_2(\Q_9)$, freely generated by
\begin{equation*}
\gamma_1=z_{++}z_{+-}=
 \begin{pmatrix}
  -1-2\sqrt{-1} & 2 \\-2 & -1+2\sqrt{-1} 
 \end{pmatrix},
\end{equation*}
\begin{equation*}
\gamma_2=z_{++}z_{-+}=
 \begin{pmatrix}
  -1-2\sqrt{-1} & -2\sqrt{-1}
  \\-2\sqrt{-1} & -1+2\sqrt{-1}  
 \end{pmatrix},
\end{equation*}
 \begin{equation*}
  \gamma_3=z_{++}z_{+-}=
   \begin{pmatrix}
    1 & 2-2\sqrt{-1} \\-2-2\sqrt{-1} & 1 
   \end{pmatrix}.
 \end{equation*}
Moreover $PGL_1(\ovl{\mathcal{B}}[\frac{1}{3}])$ is
 the semi-direct product of $PGL_1(\ovl{\mathcal{B}}[\frac{1}{3}])(2)$ 
 and the finite group $PGL_1(\ovl{\mathcal{B}})$, and $\Gamma_3^\ast$ is
 generated by $PGL_1(\ovl{\mathcal{B}}[\frac{1}{3}])$ and the 
 image of $(1+i)$.

At last, the dual graph of the stable reduction modulo $3$ of
 ${\mathcal{X}^+(2)}_{\Q_9}^{\mathrm{an}}$ is 

%%% PICTURE 6f
 \begin{figure}[h]
  \begin{picture}(100,40)(-50,-20)
   \put(-30,0){\makebox(0,0)[c]{$\bullet$}}
   \put(30,0){\makebox(0,0)[c]{$\bullet$}}
   \qbezier(-30,0)(-20,12)(0,12)
   \qbezier(30,0)(20,12)(0,12)
   \qbezier(-30,0)(0,10)(30,0)
   \qbezier(-30,0)(0,-10)(30,0)
   \qbezier(-30,0)(-20,-12)(0,-12)
   \qbezier(30,0)(20,-12)(0,-12)
  \end{picture}
  \caption{}
 \end{figure}
\end{para}

\begin{para}\label{para::6.3.3} The case of the prime $2$ is similar if one
 replaces $\mathcal{X}^+(2)$ by $\mathcal{X}^+(3)$\footnote{to deal with
 the diadic uniformization of $\mathcal{X}^+(2)$ itself would involve a
 certain finite \'etale covering of degree $12$ of $\Omega$.}. This is a
 Mumford curve over $\Q_4$, the unramified quadratic extension of
 $\Q_2$. The relevant quaternion algebra here is $B_{3.\infty}$. A
 maximal order (again unique up to conjugation) is given by $\Z+\Z i+\Z
 \frac{i+j}{2}+\Z \frac{1+ij}{2},\;\;i^2=-1,\;j^2=-3,\;ij=-ji$.
 
\begin{itemize}
 \item $\Delta_2(2,4,6)=\Gamma_2^\ast$, 
 \item $\diamondsuit_2 (2,2,3,3)=\Gamma_2^+$.
\end{itemize}

\end{para}

\subsection{The non-archimedean triangle groups $\Delta_3(2,4,6),
\Delta_2(2,4,6)$ and Herrlich's tree.}\label{sub::6.4}

\begin{para}\label{para::6.4.1} Our aim here is to explore the combinatorics
 of the finitely  
 generated groups $\Delta_3(2,4,6)$ and $ \Delta_2(2,4,6)$. 

\noindent It follows from the above considerations that
 $\Delta_3(2,4,6)$ sits in an exact sequence 
\begin{equation*}
1\to \an{\gamma_1,\gamma_2,\gamma_3}=\Gamma_3^+(2)\to \Delta_3(2,4,6)\to
 \mathfrak{S}_4\times \Z/2\Z =\Aut(\mathcal{X}^+(2))\to 1.
\end{equation*} 
 \noindent This extension is described by the outer action of
 $\mathfrak{S}_4\times \Z/2\Z$ on  $\an{\gamma_1,\gamma_2,\gamma_3}$. 
 Similarly, $\Delta_2(2,4,6)$ is an extension of a group of order $72$,
 $\Aut(\mathcal{X}^+(3))$, by a free group of rank four. In order to
 describe these extensions, we shall use F. Herrlich's work on 
 automorphisms of Mumford curves. 
\end{para}

\begin{para}\label{para::6.4.2} Let $S$ be a Mumford curve over some finite
 extension $K$ of $\Q_p$, and let $\Gamma 
 \subset PGL_2(K)$ be the corresponding Schottky group. It is known that
 any automorphism of $S$ lifts to $PGL_2(K)$: in fact, 
 the normalizer $N(\Gamma)$ of $\Gamma$ in sits $PGL_2(K)$ in an exact
 sequence \cite[IV]{sgamc}: 
\begin{equation*}
 1\to \Gamma\to N(\Gamma)\to  \Aut S\to 1.
\end{equation*} 
 \noindent Herrlich constructs a tree $\mathcal{T}$ endowed with an
 action of $N(\Gamma)$ as follows. One introduces a partition of the set
 of affinoid disks in $\P^1(K)$ into classes of the following form: a
 class contains a unique disk in $\A^1(K)$, of radius denoted by $r$ and
 all affinoid disks of the form $| z-a|\geq r$. The vertices of
 $\mathcal{T}$ are the classes of those disks which contain two fixpoints
 of some non-trivial element of $N(\Gamma)$, whose distance is the radius
 of the disk.  Two vertices are related by an edge if they can be
 represented by disks which are either both maximal, or nested in such a
 way that one is a maximal proper subdisk of the other. Herrlich shows
 that the natural action of $N(\Gamma)$ on affinoid disks induce an
 action without inversion on $\mathcal{T}$. This allows to apply the
 Bass-Serre theory of graphs of groups \cite{aas}, exploiting the fact
 that the stabilizer of vertices or edges belong to the short list of
 finite subgroups of $PGL_2(K)$
 \cite{herrlich80:_ordnun_autom_schot} \cite{herrlich80:_endlic_grupp}.
 When $S/\Aut(S)$ has genus $0$, the quotient $\mathcal{T}/N(\Gamma)$ is
 a finite tree of groups. The group $N(\Gamma)$ can then 
 be identified with its fundamental group, which is an amalgamated sum of
 the stabilizers of the vertices along the stabilizers of the edges.
\end{para}

\begin{para}\label{para::6.4.3} Let us apply this to
 $S=\mathcal{X}^+(2)_{\Q_9}$ with $N(\Gamma)=\Delta_3(2,4,6)$.  One has
 $(g-1)/\sharp(\Aut(S))=\frac{1}{24}$, and an inspection of the list of
 \cite{herrlich80:_ordnun_autom_schot} leaves only one possibility for
 the tree of groups $\mathcal{T}_3/\Delta_3(2,4,6)$, namely:

 \begin{center}
  \begin{picture}(100,35)(-50,-15)
   \put(-30,0){\makebox(0,0)[c]{$\mathfrak{D}_6$}}
   \put(-14.5,0){\makebox(0,0)[r]{$\bullet$}}
   \put(-15,0){\line(1,0){30}}
   \put(14.5,0){\makebox(0,0)[l]{$\circ$}}
   \put(30,0){\makebox(0,0)[c]{$\mathfrak{S}_4$}}
   \put(0,-10){\makebox(0,0)[c]{$\mathfrak{D}_3$}}
  \end{picture}
 \end{center}

 \medskip\noindent where $\mathfrak{D}_n$ are the dihedral groups.  
 In this case, the Herrlich tree $\mathcal{T}_3$ may be considered as a
 triadic analogue of the Schwarz-Escher tesselation:
 
%%% PICTURE 6g
\begin{figure}[h]
 \begin{picture}(180,180)(0,0)
  \put(0,0){\includegraphics[scale=0.5,clip]{fig3-6g.eps}}
 \end{picture}
 \caption{}
\end{figure}


\end{para}
 
\begin{para}\label{para::6.4.4} In the case of $S=\mathcal{X}^+(3)_{\Q_4}$
 with $N(\Gamma)=\Delta_2(2,4,6)$, one still has
 $(g-1)/\sharp(\Aut(S))=\frac{1}{24}$, and an inspection of the list of
 \cite{herrlich80:_ordnun_autom_schot} leaves only two possibilities for
 the tree of groups $\mathcal{T}_2/\Delta_2(2,4,6)$, namely: 

 \begin{center}
  \begin{picture}(100,35)(-50,-15)
   \put(-30,0){\makebox(0,0)[c]{$\mathfrak{D}_6$}}
   \put(-14.5,0){\makebox(0,0)[r]{$\bullet$}}
   \put(-15,0){\line(1,0){30}}
   \put(14.5,0){\makebox(0,0)[l]{$\circ$}}
   \put(30,0){\makebox(0,0)[c]{$\mathfrak{D}_4$}}
   \put(0,-10){\makebox(0,0)[c]{$\mathfrak{D}_2$}}
  \end{picture}
 \end{center}

\noindent and 

 \begin{center}
  \begin{picture}(100,35)(-50,-15)
   \put(-30,0){\makebox(0,0)[c]{$\mathfrak{S}_4$}}
   \put(-14.5,0){\makebox(0,0)[r]{$\bullet$}}
   \put(-15,0){\line(1,0){30}}
   \put(14.5,0){\makebox(0,0)[l]{$\circ$}}
   \put(30,0){\makebox(0,0)[c]{$\mathfrak{S}_4$}}
   \put(0,-10){\makebox(0,0)[c]{$\mathfrak{D}_4$}}
  \end{picture}
 \end{center}

\medskip\noindent In order to select the right one, we remark that
 $\Delta_2(2,4,6)$ contains the finite group generated by $ij$ and the
 quotient of $(\mathcal{B}_{3.\infty})^\times=
 \bigl\{\pm1,\;\pm i,\; \frac{\pm i\pm j}{2},\; \frac{\pm 1\pm ij}{2}\bigr\}$
 by $\pm 1$. This group is
 isomorphic to $\mathfrak{D}_6\cong \Z/2\Z\;\times \mathfrak{D}_3$,
 which is not contained in $\mathfrak{S}_4$.  Since any finite subgroup
 of the amalgamated sum of two groups is conjugated to a subgroup of one
 of these groups \cite[4.3]{aas}, we conclude that
 $\Delta_2(2,4,6)\cong \mathfrak{D}_6
 \ast_{\mathfrak{D}_2}\mathfrak{D}_4$,
 not $\mathfrak{S}_4 \ast_{\mathfrak{D}_4}\mathfrak{S}_4$

\medskip\noindent
The Herrlich tree $\mathcal{T}_2$ may be considered as a diadic analogue of
the Schwarz-Escher tesselation:

%%% PICTURE 6h    
\begin{figure}[h]
 \begin{picture}(180,180)(0,0)
  \put(0,0){\includegraphics[scale=0.5,clip]{fig3-6h.eps}}
 \end{picture}
 \caption{}
\end{figure}

\bigskip\noindent
{\small (this is the set of points $u$ such that 
$j(\frac{e^{\pi i/3}-u}{1-e^{\pi i/3}u})\in [0,1728]$, 
where $j$ denotes the modular invariant).}

\medskip
It is important to pay attention to the fact that the symbol 
$\mathfrak{D}_6 \ast_{\mathfrak{D}_2}\mathfrak{D}_4$ is not free from
ambiguity: namely, there are essentially two non-equivalent ways of
embedding $\mathfrak{D}_2$ into $\mathfrak{D}_6\times \mathfrak{D}_4$:

{\it $1$st way$)$} the centers of $\mathfrak{D}_6$ and $\mathfrak{D}_4$ 
restrict to the same subgroup $\Z/2\Z$
of $\mathfrak{D}_2$. On identifying $\mathfrak{D}_6$ and $\mathfrak{D}_4$ with 
the matrix groups
\begin{equation*}
 \Bigl\langle\begin{pmatrix} 0& 1 \\ 1& 0
\end{pmatrix}, \begin{pmatrix} 0& -1 \\ 1& 1
\end{pmatrix}\Bigr\rangle\text{ and }
 \Bigl\langle\begin{pmatrix} 0& 1 \\ \pm1& 0
\end{pmatrix}\Bigr\rangle
\end{equation*}
respectively, the corresponding amalgam 
$\mathfrak{D}_6 \ast_{\mathfrak{D}_2}\mathfrak{D}_4$ is isomorphic 
to $GL_2(\Z)$;

{\it $2$nd way $)$} the trace of the center of $\mathfrak{D}_6\times 
\mathfrak{D}_4$ on $\mathfrak{D}_2$  is $\mathfrak{D}_2$ itself. The
corresponding amalgam $\mathfrak{D}_6
\ast_{\mathfrak{D}_2}\mathfrak{D}_4$ maps onto  
a semidirect product of $(\Z/3\Z)^2$ and ${\mathfrak{D}_4}$ (the 
latter acting on $(\Z/3\Z)^2$ as a matrix group
$\Bigl\langle
\begin{pmatrix}
 0& 1 \\ \pm1& 0
\end{pmatrix}
\Bigr\rangle$).
\end{para}

The following theorem summaries our results and settles the remaining 
ambiguity:

\begin{thm}\label{thm::6.4.5}
 \begin{enumerate}
  \item $\Delta_3(2,4,6)\cong \mathfrak{D}_6 
        \ast_{\mathfrak{D}_3} \mathfrak{S}_4$.
  \item $\Delta_2(2,4,6)\cong\mathfrak{D}_6\ast_{\mathfrak{D}_2}
	\mathfrak{D}_4$, where the embedding of $\mathfrak{D}_2$ into
	$\mathfrak{D}_6\times \mathfrak{D}_4$ is such that the trace of
	the center of $\mathfrak{D}_6\times \mathfrak{D}_4$ on
	$\mathfrak{D}_2$ is $\mathfrak{D}_2$ itself.
 \end{enumerate} 
\end{thm}

\begin{proof}
 Only the last assertion remains to be established. Let $\mathcal{S}$ be 
 the diadic Schwarz orbifold
 $(\P^1;(0,2),(1,4),(\infty,6))$.  There are surjective homomorphisms 
 (3.5.5, \ref{pro::4.3.3})
 \begin{equation*}
  \pi_1^{\mathrm{temp}}(\P^1\setminus
   \{0,1,\infty\})\to
   \pi_1^{\mathrm{orb}}(\mathcal{S})\to \Delta_2(2,4,6).
 \end{equation*} 
 The images 
 $\bar\gamma_0,\bar\gamma_1,\bar\gamma_\infty\in  \Delta_2(2,4,6)$
 of local monodromy elements $\gamma_0,\gamma_1,\gamma_\infty \in 
 \pi_1^{\mathrm{temp}}(\P^1\setminus
 \{0,1,\infty\})$ are elements of order $2,4,6$ respectively. By 
 \ref{pro::2.3.9}, the smallest closed normal
 subgroup of $\pi_1^{\mathrm{temp}}(\P^1\setminus
 \{0,1,\infty\})$ containing $\gamma_1,\gamma_\infty$ is 
 $\pi_1^{\mathrm{temp}}(\P^1\setminus
 \{0,1,\infty\})$ itself. It follows that the smallest normal
 subgroup of the discrete group $\Delta_2(2,4,6)$ containing 
 $\bar\gamma_1,\bar\gamma_\infty$ is $\Delta_2(2,4,6)$ itself.
 It thus suffices to show that if we where in the case $\mathfrak{D}_6 
 \ast_{\mathfrak{D}_2}\mathfrak{D}_4\cong
 GL_2(\Z)$, then $\bar\gamma_1,\bar\gamma_\infty\in SL_2(\Z)$. This is 
 clear, since any element of order $6$ (\resp $4$) in $ GL_2(\Z)$ is
 conjugate to $\pm
 \begin{pmatrix}
  0 & -1 \\ 1& 1
 \end{pmatrix}$
 (\resp to
 $\pm\begin{pmatrix}
      0& -1 \\ 1& 0
     \end{pmatrix}$).
\end{proof}

\medskip Part (ii) of this theorem allows us to answer at last 
question \ref{que::2.3.11} concerning
local monodromies $\gamma_0,\gamma_1,\gamma_\infty$ at $0,1,\infty$ 
in the temperate fundamental group
$\pi_1^{\mathrm{temp}}(\P^1_{\C_p}\setminus \{0,1,\infty\})$:

\begin{cor}\label{cor::6.4.6} Question \ref{que::2.3.11} has a negative
 answer for $p=2$: one cannot find local monodromies
 $\gamma_0,\gamma_1,\gamma_\infty$ such that 
 $\gamma_0.\gamma_1.\gamma_\infty = 1$.
\end{cor}

\begin{proof}
 Assume that such local monodromies exist.  By means of their images
 $\bar\gamma_0,\bar\gamma_1,\bar\gamma_\infty \in \Delta_2(2,4,6)\cong 
 \mathfrak{D}_6 \ast_{\mathfrak{D}_2}\mathfrak{D}_4 $ (notation of 
 the previous proof), it would be possible to construct
 a homomorphism
 \begin{equation*}
  \phi: \Delta (2,4,6)\to\Delta_2(2,4,6),
 \end{equation*}
 where $\Delta (2,4,6)$ 
 denotes the standard (archimedean) triangle group.
 Notice that the image of $\phi$ cannot be contained in any conjugate 
 of $\mathfrak{D}_6$ or $\mathfrak{D}_4$, since none of
 these finite groups contains elements of order $6$ and of order $4$ 
 simultaneously.

  According to \cite[6.3.5]{aas}, any (cocompact) triangle group 
 $\Delta$ has the
 fixpoint property $(FA)$: any action (without inversion) of $\Delta$ 
 on any tree has a fixpoint. It follows that the
 quotient
 $\Im\phi$ of $\Delta (2,4,6)$ would have property
 $(FA)$ as well. Since $Im\,\phi$ is a subgroup of the amalgam
 $\mathfrak{D}_6 \ast_{\mathfrak{D}_2}\mathfrak{D}_4$, this would 
 imply, according to \cite[6.2, prop. 21]{aas}, that
 $Im\,\phi$ is contained in a conjugate of $\mathfrak{D}_6$ or 
 $\mathfrak{D}_4$, a contradiction. 
\end{proof}

\begin{rem}
 The arithmetic $p$-adic triangle groups 
 $\Delta (e_0,e_1,e_\infty)$ do {\it not} have property
 $(FA)$. Indeed, their description as $p$-unit groups in quaternion 
 algebras shows that they admit a two-dimensional
 representation for which the eigenvalues are non-integral in general. 
 This property contradicts $(FA)$, \cf \cite[6.2, prop. 22]{aas}.
   Actually, Kato's description of $p$-adic triangle groups of Mumford 
 type as amalgams, together with \ref{thm::5.3.7}, implies that
 {\it no hyperbolic $p$-adic triangle group has property $(FA)$} (in 
 contrast to archimedean triangle groups).  Hence no hyperbolic $p$-adic
 triangle group can be generated by three elements $g_0,g_1,g_\infty$ of
 finite order satisfying $g_0g_1g_\infty=1$.
\end{rem}

\subsection{$F(\frac{1}{24}, \frac{7}{24},\frac{5}{6}; z)$ and the period 
mapping. Global analytic triadic continuation.}\label{sub::6.5}

\begin{para}\label{para::6.5.1} Rather than $\Delta(2,4,8)$, we now
 prefer to work with $\Delta(6,2,4)$ (this amounts to permute
 $0,1,\infty$ by $t\mapsto z=\frac{1}{1-t}$, and is more faithful to the
 first picture of this section which places the vertex 
 $O$ at the center of the disk). 

 The triangle group
$\Delta(6,2,4)$ is the projective monodromy group of the hypergeometric
 differential operator 
$L_{\frac{1}{24}, \frac{7}{24},\frac{5}{6}}$. More precisely,
 $L_{\frac{1}{24}, \frac{7}{24},\frac{5}{6}}(y)=0$ 
 is a {\it uniformizing differential equation} for the orbifold
 $\mathcal{X}^\ast$ viewed as the Schwarz orbifold 
\begin{equation*}
\left(\P^1,\;
 \begin{matrix}
  0& 1& \infty \\e_0=6& e_1=2& e_\infty = 4
 \end{matrix}
\right)
\end{equation*}
More explicitly, $F(\frac{1}{24}, \frac{7}{24},\frac{5}{6}; z)$ and
 $z^{1/6}F(\frac{5}{24}, \frac{11}{24},\frac{7}{6}; z)$ are 
independent solutions of $L_{\frac{1}{24}, \frac{7}{24},\frac{5}{6}}$
 and for a suitable constant $\kappa$, the mapping 

\begin{equation*}
\tau(z)= \kappa^{-1}. \frac{z^{1/6}F(\frac{5}{24},
 \frac{11}{24},\frac{7}{6}; z) }{F(\frac{1}{24},
 \frac{7}{24},\frac{5}{6}; z)}
\end{equation*} 

\noindent sends the upper half plane bijectively to the basic triangle
 $OQP$ of the Schwarz tesselation, and 
 $\tau(0)=O,\;\tau(1)=Q,\;\tau(\infty)=P$. 

\begin{figure}[h]
 \begin{picture}(320,100)(0,-15)
  \put(0,0){\includegraphics[scale=0.8,clip]{fig3-6i-1.eps}}
  \put(57,-1){\makebox(0,0)[b]{\small $\bullet$}}
  \put(57,-8){\makebox(0,0)[c]{\small $0$}}
  \put(93,-1){\makebox(0,0)[b]{\small $\bullet$}}
  \put(93,-8){\makebox(0,0)[c]{\small $1$}}
  \put(180,25){\vector(1,0){50}}
  \put(205,35){\makebox(0,0)[c]{$\tau$}}
  \put(250,0){\includegraphics[scale=0.8,clip]{fig3-6i-2.eps}}
  \put(250,-8){\makebox(0,0)[c]{\small $O$}}
  \put(302,-8){\makebox(0,0)[c]{\small $Q$}}
  \put(312,38){\makebox(0,0)[c]{\small $P$}}
 \end{picture}
 \caption{}
\end{figure}

\noindent It extends to a multivalued analytic function on $\P^1
 \setminus \{0,1,\infty\}$ with value in $\matheur{D}(0,1^-)$. The 
inverse mapping
$z(\tau): \matheur{D}(0,1^-)\to \P^1$ is the (single-valued)
$\Delta(2,4,6)$-automorphic uniformizing mapping of the Shimura orbifold
 $\mathcal{X}^\ast $.
\end{para}

\begin{para}\label{para::6.5.2}  Actually, $\tau$ expresses the period mapping
 $\mathcal{P}$ for ``fake elliptic curves'' with multiplication by 
 $\mathcal{B}_{2.3}$, \cf  II.\ref{sub:1.2} and
 II.\ref{sub:7.4}: more precisely, 
 the multivalued function $\tau(z)$ is related to 
 the commutative diagram  
 \begin{equation*}
  \begin{CD}
   \widetilde{{\mathcal{X}^\ast_\C}^{\mathrm{an}}}
    @>{}^{\begin{subarray}\mathcal{P}\\\sim\end{subarray}}>>
   \mathcal{D}=\matheur{D}(0,1^-) @. \quad \subset \;\mathcal{D}^\vee=\P^1 \\
   @V\mathcal{Q}VV @VVV\\
   \left(\P^1,\;
   \begin{matrix}
    0& 1& \infty \\
    6& 2& 4
   \end{matrix}
   \right)
   ={\mathcal{X}^\ast_\C}^{\mathrm{an}}
   @>>> \mathcal{D}/\Delta(2,4,6)
  \end{CD}
 \end{equation*}
by
 \begin{equation*}
  \tau(z)= \mathcal{P}\circ {\mathcal{Q}}^{-1},\; z(\tau )= {\mathcal{Q}}\circ
   \mathcal{P}^{-1}.
 \end{equation*} 
\end{para}

\begin{para}\label{para::6.5.3} Using Gauss' formula for the value of
 hypergeometric functions at $1$, one can compute the constant $\kappa$, 
 \cf \cite[392]{f2}:
 \begin{equation*}
 \kappa=\frac{\Gamma (\frac{7}{6})}{\Gamma
  (\frac{5}{6})}\sqrt\frac{\Gamma (\frac{1}{24})\Gamma
  (\frac{7}{24})\Gamma (\frac{13}{24})\Gamma (\frac{19}{24}) 
 }{\Gamma (\frac{5}{24}) \Gamma (\frac{11}{24})\Gamma
 (\frac{17}{24})\Gamma (\frac{23}{24})}. 
 \end{equation*}

 Using the functional equations of $\Gamma$, one checks that $\kappa
 \sim \Bigl(\frac{\Gamma(\frac{1}{3})}{\Gamma(\frac{2}{3})}\Bigr)^3$
 $\pmod{\ovl{\Q}^\times}$, in conformity with II.\ref{thm:7.6.4} (taking into
 account the fact that $z=0$ corresponds to a fake elliptic curve with
 complex multiplication by an order in $M_2(\Q(\sqrt{-3}))$). 
\end{para}

\begin{para}\label{para::6.5.4} For $p\neq 2,3$, the $p$-adic radius of
 convergence of $F(\frac{1}{24}, \frac{7}{24},\frac{5}{6}; z)$ and 
 $F(\frac{5}{24}, \frac{11}{24},\frac{7}{6}; z)$ is $1$. The unit disk
 $|z|_p<1$ is ordinary or supersingular according 
 to whether $p\equiv 1$ $\pmod{3}$ or not. The ratio 
 $\frac{z^{1/6}F(\frac{5}{24}, \frac{11}{24},\frac{7}{6}; z)
 }{F(\frac{1}{24}, \frac{7}{24},\frac{5}{6}; z)}$ expresses a local 
 $p$-adic period mapping. 

 The $p$-divisible group of the fake elliptic curve parametrized by $z$
 is isogenous to the square of the $p$-divisible group of an elliptic
 curve; the period 
 mapping is thus completely analogous to one of those studied 
 in II.\ref{para:6.3.1}, \ref{para:6.3.3} (Dwork, 
 Gross-Hopkins). From the viewpoint of analytic continuation, this is
 essentially the situation studied in  I.\ref{sec-tale} 
 (for $F(\frac{1}{2},\frac{1}{2},{1}; z)$).
\end{para}

\begin{para}\label{para::6.5.5} For $p=3$ (\resp $p=2$), the
 situation is radically different; the $p$-adic radius of convergence of 
 $F(\frac{1}{24}, \frac{7}{24},\frac{5}{6}; z)$ and
 $F(\frac{5}{24}, \frac{11}{24},\frac{7}{6}; z)$ is $3^{-3/2}$
 (\resp $2^{-6}$).

 The ratio
 $ \frac{z^{1/6}F(\frac{5}{24},\frac{11}{24},\frac{7}{6}; z)}
 {F({\frac{1}{24},{\frac{7}{24}},\frac{5}{6}}; z)}$
 expresses the Drinfeld $3$-adic period mapping: 

 \begin{center}
  \begin{picture}(120,125)(0,-10)
   \put(45,105){\makebox(0,0)[c]{${}^\mathcal{P}$}}
   \put(45,95){\makebox(0,0)[c]{$\sim$}}
   \put(0,90){\makebox(0,0)[c]{$\til{\mathcal{X}^+(2)_{\C_3}^{\mathrm{an}}}$}}
   \put(0,75){\vector(0,-1){15}}
   \put(0,45){\makebox(0,0)[c]{${\mathcal{X}^+(2)}_{\C_3}^{\mathrm{an}}$}}
   \put(0,30){\vector(0,-1){15}}
   \put(0,0){\makebox(0,0)[c]{${\mathcal{X}^{\ast\,\mathrm{an}}_{\C_3}}2$}}
   \put(30,90){\vector(1,0){30}}
   \put(80,90){\makebox(0,0)[dc]{$\mathcal{D}$}}
   \put(97,90){\makebox(0,0)[c]{$\subset$}}
   \put(110,90){\makebox(0,0)[dl]{$\mathcal{D}^{\vee}=\P^1$}}
   \put(80,75){\vector(0,-1){60}}
   \put(80,0){\makebox(0,0)[c]{$\mathcal{D}/\Delta_3(2,4,6)$}}
   \put(25,0){\vector(1,0){15}}
  \end{picture}
 \end{center}

 \noindent 
 Here $\mathcal{D}$ is the image of the Drinfeld space $\Omega_{\C_3}$
 by a suitable homography (it would be interesting to 
 determine this homography explicitly). That this homography is actually
 defined over $\ovl{\Q}$ 
 comes from II.\ref{thm:7.6.4} (ramified variant), I.\ref{para-ramif},
 \ref{thm-cm-period} (the point is
 that the triadic periods, in the sense of
 I.\ref{sub-p-adic-betti-lattice-supersing}, of the abelian surface 
 with complex multiplication by $M_2(\Q[\sqrt{-3}])$ parametrized by
 $z=0$, belong to $\ovl{\Q}$).

 \bigskip Although $F(\frac{5}{24}, \frac{11}{24},\frac{7}{6};
 z)/F(\frac{1}{24}, \frac{7}{24},\frac{5}{6}; z)$ has very small triadic
 radius of convergence, its pull-back on the Shimura curve
 $\mathcal{X}^+(2)$ extend to a global multivalued meromorphic
 function, with poles only above $0,1,\infty$ (in other words, its
 pull-back on $\widetilde{\mathcal{X}^+(2)}_{\C_3} \cong \Omega_{\C_3}$
 is single-valued).

 Note however that the same is not true for $F(\frac{1}{24},
 \frac{7}{24},\frac{5}{6}; z)$ and $F(\frac{5}{24},
 \frac{11}{24},\frac{7}{6}; z)$ individually, for reasons of exponents
 at $\infty$.  Nevertheless, they enjoy a similar property, on replacing
 $\mathcal{X}^+(2)$ by a suitable ramified finite covering.

 \medskip Similarly, in the diadic case: the pull-back of
 $F(\frac{5}{24}, \frac{11}{24},\frac{7}{6}; z)/F(\frac{1}{24},
 \frac{7}{24},\frac{5}{6}; z)$ on the Shimura curve $\mathcal{X}^+(3)$
 extends to a global multivalued meromorphic functions (with poles only
 above $0,1,\infty$).
\end{para}
 

\subsection{A triadic analogue of Gauss' formula, and special values of
$F(\frac{1}{24}, \frac{7}{24},\frac{5}{6}; z)$ at CM points.}\label{sub::6.6}


\begin{para}\label{para::6.6.1} For our hypergeometric function, Gauss'
 formula reads
 \begin{equation*}
  F\Bigl(\frac{1}{24}, \frac{7}{24},\frac{5}{6};1\Bigr)=
 \frac{\Gamma(\frac{5}{6})\Gamma(\frac{1}{2})}
 {\Gamma(\frac{19}{24})\Gamma(\frac{13}{24})}.
 \end{equation*}

\noindent If one considers $F(\frac{1}{24}, \frac{7}{24},\frac{5}{6};
 z)$ as a multivalued analytic function on
 $\P^1\setminus\{0,1,\infty\}$, the values at $1$ generate a
 $\ovl{\Q}$-subspace of $\C$ of dimension $\leq 2$, namely:  
\begin{equation*}
 \frac{\Gamma(\frac{5}{6})\Gamma(\frac{1}{2})}
  {\Gamma(\frac{19}{24})\Gamma(\frac{13}{24})}\ovl{\Q} \,+ \,
  \frac{\Gamma(\frac{5}{6})\Gamma(-\frac{1}{2})}
  {\Gamma(\frac{1}{24})\Gamma(\frac{7}{24})}\ovl{\Q}.
\end{equation*}
Similarly, its values at $\infty$ generate the $\ovl{\Q}$-subspace 

\begin{equation*}
 \frac{\Gamma(\frac{5}{6})\Gamma(\frac{1}{4})}
  {\Gamma(\frac{19}{24})\Gamma(\frac{7}{24})}\ovl{\Q} \, +\,
  \frac{\Gamma(\frac{5}{6})\Gamma(-\frac{1}{4})}
  {\Gamma(\frac{1}{24})\Gamma(\frac{13}{24})}\ovl{\Q}.
\end{equation*}   
\end{para}

\begin{para}\label{para::6.6.2} Let us consider our hypergeometric
 function as a triadic function $F(\frac{1}{24},
 \frac{7}{24},\frac{5}{6}; 1)_3$, and more precisely as a multivalued
 $3$-adic analytic function on a suitable finite \'etale covering of
 $\P^1\setminus\{0,1,\infty\}$. 
 
 The values taken by $F(\frac{1}{24}, \frac{7}{24},\frac{5}{6}; z)_3$ at
 $1$ belong to $\ovl{\Q}$: this follows from \ref{para::5.6.3}, together
 with the fact that the ``periods'' (in the sense of
 I.\ref{sub-p-adic-betti-lattice-supersing}) attached to the 
 abelian surface with complex multiplication by $M_2(\Q[\sqrt{-24}])$ 
 parametrized by $z=1$, belong to $\ovl{\Q}$, by I.\ref{para-ramif}, 
 I.\ref{thm-cm-period}. 

 \medskip In the same way, one shows that the values taken by
 $F(\frac{1}{24}, \frac{7}{24},\frac{5}{6}; z)_3$ at $\infty$ belong to 
 $(\Gamma_3(\frac{3}{4})/\Gamma_3(\frac{1}{4}))^{-1/2}
 \ovl{\Q}=\Gamma_3(\frac{1}{4})\ovl{\Q}$.
 Conjecture I.\ref{coj-no-complex-multi} predicts that
 $\Gamma_3(\frac{1}{4})$ is a transcendental number.
\end{para}

\begin{para}\label{para::6.6.3} In \cite[Table 2]{elkies98:_shimur}, one finds a
 table of rational CM points $z$ (rational points which correspond to
 abelian surfaces with complex multiplication by an order in
 $M_2(\Q(\sqrt{-d}))$ for some fundamental discriminant $-d$); with the
 notation of {\it loc.~cit.}, they are given by $z=B/(B-A)= \pm B /|
 C|$. Some of them are: 

\vspace{5mm}
\begin{center}
 \renewcommand{\arraystretch}{1.2}
 \begin{tabular}{c @{\quad\qquad}l}
 $ z            $ & $ d=-\mathrm{disc}   $ \\ \hline 
 $ -3^3/2^2 7^2 $ & $ d=2^2.3.7$ \\
 $ -5^3/3^7     $ & $ d=2^3.5  $ \\
 $ 2^{10}/7^4   $ & $ d=3.17   $ \\
 $ -2^{10}/3^7  $ & $ d=19.    $ 
 \end{tabular} 
\end{center}
\vspace{3mm}

 Notice that for the first value of $z$, $F(\frac{1}{24},
 \frac{7}{24},\frac{5}{6}; \frac{-3^3}{2^2 7^2})_3$ converges. Moreover,
 since $(\frac{2^2.7}{3})=+1$, it follows from I.\ref{para-ramif},
 I.\ref{thm-cm-period}, that 
 the triadic ``periods'' attached to the abelian surface $A_{-3^3/ 2^2 7^2}$ 
 with complex multiplication by $M_2(\Q[\sqrt{-d}])$
 parametrized by $z=1$, is algebraic (like the periods of $A_0$). By 
 \ref{para::5.6.3}, it follows that the triadic number
 $F(\frac{1}{24}, \frac{7}{24},\frac{5}{6}; \frac{-3^3}{2^2 7^2})_3 $ is
 algebraic (can one compute this algebraic number?). 
 
 \medskip 
 On the other hand since
 $\Q(\sqrt{-d})\neq \Q(\sqrt{-3})$ in that case ($d=2^2.3.7 $), the real
 number $F(\frac{1}{24}, \frac{7}{24},\frac{5}{6}; \frac{-3^3}{2^2 7^2})$
 is transcendental; this follows from Wolfart's theorem 
 \cite{wolfart88:_werteh_funkt}
 (whose proof has been recently completed by work of Cohen-W\"ustholz
 \cite{cohen01:_number_theor} and Edixhoven-Yafaev).

 \medskip It can be proved, a contrario, that if the complex evaluation
 $F(\frac{1}{24}, \frac{7}{24},\frac{5}{6}; \xi)$ is algebraic for some
 algebraic value $\xi$ in the disc of convergence, then any $p$-adic
 evaluation $F(\frac{1}{24}, \frac{7}{24},\frac{5}{6}; \xi)_p$ which makes
 sense is also algebraic: indeed, by 
 \cite{wolfart88:_werteh_funkt}, \cite{cohen01:_number_theor}, such a
 $\xi$ parametrizes a CM abelian surface by some order in
 $M_2(\Q[\sqrt{-3}])$, and one concludes again using $p$-adic Betti
 lattices as in I.\ref{sub-p-adic-betti-lattice-supersing}.
 
 This is actually a very special case of a general, 
 but conjectural, ``principle of global relations'' for special values of
 solutions of Gauss-Manin connections, \cf \cite{tdmeigdvpdg}.
\end{para}
